\documentclass[11pt]{article}

\usepackage[margin=1in]{geometry}
\usepackage{amsmath,amsfonts,amssymb}
\usepackage{graphicx}
\usepackage{booktabs}
\usepackage{multirow}
\usepackage{array}
\usepackage{url}
\usepackage{hyperref}
\usepackage{xcolor}
\usepackage{caption}
\usepackage{float}

\newcommand{\todo}[1]{\textcolor{red}{\textbf{[TODO: #1]}}}
\newcommand{\pending}[1]{\textcolor{blue}{#1}}
\DeclareMathOperator*{\argmin}{arg\,min}
\DeclareMathOperator*{\argmax}{arg\,max}

\title{Measuring the Imperceptibility and Verification-Leakage of\\
Pivotal Objective Perturbation (POP) Audio Cloaking\\
Across a Demographically Balanced, Multi-Corpus Speaker Set}

\author{
  \todo{Author names}\\
  \todo{Affiliation}
}

\date{\today}

\begin{document}
\maketitle

\begin{abstract}
Voice cloning systems can now reproduce a target speaker's identity
from only a few seconds of reference audio, motivating a line of
\emph{proactive} defenses that perturb a speaker's recordings before
they are ever published, so that any model trained or conditioned on
the released audio fails to reproduce the speaker's true voice. We
present a large-scale empirical evaluation of one such defense,
Pivotal Objective Perturbation (POP)~\cite{zhang2024mitigating}, on a
demographically balanced, multi-corpus dataset of 4,867 real speech
recordings drawn from LibriSpeech~\cite{panayotov2015librispeech},
ASVspoof~2021~\cite{yamagishi2021asvspoof}, and
FakeAVCeleb~\cite{khalid2021fakeavceleb}, the latter balanced across
five descent groups and two genders. We reconstruct POP's official
attack pipeline end to end -- a pretrained VITS~\cite{kim2021vits}
surrogate, grapheme-to-phoneme text conditioning, and a 200-step
projected-gradient attack under an $\ell_\infty$ budget of
$\varepsilon=8/255$ -- and apply it to every recording in the corpus.
We then measure two properties that are conflated in much of the
cloaking literature but are conceptually distinct: (i) whether the
perturbation is \emph{imperceptible}, i.e. whether the cloaked
recording still verifies as the same speaker under three independent
speaker-verification back-ends (SpeechBrain
ECAPA-TDNN~\cite{desplanques2020ecapa}, Resemblyzer, and
WavLM-based x-vectors~\cite{chen2022wavlm}), reported via Equal Error
Rate (EER), Receiver Operating Characteristic Area Under the Curve
(ROC-AUC), the minimum normalized Detection Cost Function (minDCF),
True-Accept-Rate at fixed False-Accept-Rate operating points
(TAR@FAR), and Cohen's $d'$; and (ii) whether the perturbation
degrades the underlying signal quality, reported via the Short-Time
Objective Intelligibility measure (STOI), Perceptual Evaluation of
Speech Quality (PESQ), and Scale-Invariant Signal-to-Distortion Ratio
(SI-SDR). We describe in full detail the trial-construction protocol
(genuine vs.\ impostor trials), the three verification back-ends, the
two zero-shot voice-cloning systems and one voice-conversion system
used elsewhere in our broader evaluation
(SV2TTS~\cite{jia2018transfer}, F5-TTS~\cite{chen2024f5tts}, and
Seed-VC), and the four cloaking techniques we implemented end to end
on identical hardware (POP, ProtectYourAudio~\cite{liu2023protecting},
attack-vc~\cite{huang2021defending}, and
AntiFake~\cite{yu2023antifake}), including the practical engineering
obstacles each one presented. Across all three verifiers, POP-cloaked
audio is verified as the genuine speaker with EER~$\leq$~0.16\% and
TAR@1\%FAR~=~100.00\%, uniformly across ten FakeAVCeleb
descent$\times$gender groups (Section~\ref{sec:results}), while STOI
(0.971), PESQ (2.92), and SI-SDR (16.3~dB) confirm the perturbation
is acoustically near-transparent -- exactly the imperceptibility
POP is designed to achieve, though this measures imperceptibility, not
protective efficacy against downstream re-synthesis
(\S\ref{sec:limitations}). We additionally report the resynthesis-based
protective-efficacy evaluation this distinction calls for
(Section~\ref{sec:protective-efficacy}): cloning from POP-cloaked
source audio meaningfully degrades a zero-shot TTS attacker's output
(SV2TTS EER up 15--86\% relative over the clean-audio ceiling) but has
only a marginal effect against a voice-conversion attacker (Seed-VC
EER up 5--15\% relative), indicating POP's protective effect is
system-dependent rather than uniform. We report the identical
evaluation for attack-vc against all three downstream systems
(Section~\ref{sec:protective-efficacy-attackvc}), finding an almost
identical pattern -- strong against SV2TTS (EER up 15--95\% relative),
negligible against Seed-VC (EER up 2--6\% relative), and a nuanced
in-between result for F5-TTS whose EER rises moderately (10--28\%
relative) while its more security-relevant TAR@1\%FAR barely moves --
despite attack-vc and POP being mechanistically unrelated attacks. Both
techniques' protective efficacy is also demographically uneven in a
way that predates cloaking entirely: women and, system-dependently,
certain Asian descent groups are consistently harder for these
downstream systems to reproduce even from clean, uncloaked audio, and
cloaking widens rather than introduces this pre-existing gap. AntiFake
breaks this pattern in two ways (Sections~\ref{sec:results-antifake},
\ref{sec:protective-efficacy-antifake}): its imperceptibility fails
outright (EER 22--66\%, roughly two orders of magnitude worse than POP
or attack-vc), and its protective effect is strong against \emph{all
three} downstream systems rather than only SV2TTS -- SV2TTS (EER up
179--814\% relative), Seed-VC (EER up 116--1{,}485\% relative), and
F5-TTS (EER up 226--2{,}199\% relative, the largest effect of any
(system, verifier) cell in this project) -- a tradeoff consistent with
AntiFake's decoy-target attack optimizing embedding displacement
directly rather than staying within a fixed imperceptibility budget the
way POP and attack-vc do. AntiFake's demographic pattern is also
inverted relative to the other two techniques for SV2TTS and Seed-VC:
the same Caucasian descent groups that show the \emph{worst}
imperceptibility also show the \emph{strongest} downstream protection
under those two systems, indicating a single displacement mechanism
drives both outcomes together, unlike POP/attack-vc where the
demographic pattern in protective efficacy is inherited from the
downstream pipeline rather than the cloak itself; F5-TTS breaks even
this pattern; its gender gap nearly vanishes and its descent ranking
partially inverts, which we report as an open question rather than
force into the same account. We complete the four-technique comparison
with ProtectYourAudio (Sections~\ref{sec:results-protectyouraudio},
\ref{sec:protective-efficacy-protectyouraudio}), whose discrete,
budget-constrained frequency-band-masking search achieves
imperceptibility indistinguishable from attack-vc's verifier by
verifier (EER below 2.1\%, TAR@1\%FAR above 96.2\%) while protecting
measurably more broadly across downstream systems than either POP or
attack-vc -- Seed-VC EER up 15--100\% relative and F5-TTS EER up
62--151\% relative, both far above POP's and attack-vc's
near-negligible Seed-VC/F5-TTS effects, though still well short of
AntiFake's uniform, imperceptibility-sacrificing protection.
ProtectYourAudio's demographic pattern follows POP's and attack-vc's
pipeline-inherited direction, not AntiFake's inverted one,
demonstrating that near-perfect imperceptibility and broader
cross-system protection are not as sharply opposed as POP's and
attack-vc's results alone would suggest.
\end{abstract}

\section{Introduction}
\label{sec:intro}

Modern zero-shot text-to-speech (TTS) and voice-conversion (VC)
systems can clone a speaker's identity from only a handful of seconds
of reference audio, without any speaker-specific fine-tuning. This
capability, while enabling legitimate applications such as
accessibility tools, dubbing, and personalized assistants, also lowers
the barrier for impersonation attacks: fraud via cloned voices,
non-consensual synthetic media, and circumvention of voice-based
authentication. Two broad families of countermeasure have emerged.
\emph{Reactive} defenses attempt to detect already-synthesized
deepfake audio after the fact~\cite{asvspoof2021doi}.
\emph{Proactive} (or \emph{cloaking}) defenses instead act
\emph{before} an attacker ever obtains clean reference audio: the
speaker (or a platform on their behalf) applies an imperceptible
perturbation to a recording prior to publication, such that any
downstream cloning attempt -- inference-time or training-time --
degrades.

This paper is part of a broader project reproducing and stress-testing
four representative cloaking defenses end to end, from scratch, on a
single shared hardware and evaluation stack, so that differences in
reported effectiveness across the original papers cannot be attributed
to differences in dataset, verifier, or evaluation protocol. In this
paper we report the first complete leg of that evaluation: Pivotal
Objective Perturbation (POP)~\cite{zhang2024mitigating}, a defense
against \emph{training-time} voice cloning that perturbs each
recording so that a text-to-speech model trained or fine-tuned on the
cloaked audio produces low-quality synthetic speech at inference time.
Unlike inference-time cloaking (e.g. AntiFake, ProtectYourAudio,
attack-vc, discussed in Section~\ref{sec:cloaking}), POP does not aim
to make the \emph{cloaked recording itself} fail speaker verification
-- indeed, POP's own design goal is for the cloaked recording to
remain perceptually and, we show, verification-indistinguishable from
the original. What POP aims to break is a downstream model trained on
that recording. This distinction matters for how the evaluation in
this paper should be read: the metrics we report here characterize
POP's \emph{imperceptibility} (does the cloak avoid being flagged as
a different speaker on its own?) and its \emph{signal-level cost}
(how much does the perturbation degrade objective audio quality?), not
yet its \emph{protective efficacy} against a downstream cloning
attacker, which requires an additional synthesis step we describe as
future work in Section~\ref{sec:limitations}.

Our contributions are:
\begin{itemize}
  \item A from-scratch, faithful reimplementation of POP's official
    attack pipeline (surrogate VITS model, phonemization, PGD attack)
    applied to a corpus more than 3$\times$ larger and demographically
    far more diverse than POP's original evaluation.
  \item A demographically balanced, tri-corpus evaluation dataset
    (\S\ref{sec:dataset}) spanning read speech (LibriSpeech), spoofed
    logical-access telephony speech (ASVspoof 2021), and
    celebrity/YouTube speech balanced across five descent groups and
    two genders (FakeAVCeleb), with per-group breakdowns throughout.
  \item A full accounting of every evaluation metric used in this
    project (\S\ref{sec:metrics}), including exact formulas,
    hyperparameters, and the rationale for including impostor trials
    at all -- a point that is frequently under-explained in the
    cloaking literature.
  \item A description of every model choice in the broader project
    (three speaker verifiers, three voice-cloning/conversion systems,
    four cloaking techniques) with justification for each
    (\S\ref{sec:verifiers}--\S\ref{sec:cloaking}), including practical
    engineering obstacles (GPU memory limits, missing system
    dependencies, silent data-loader bugs) encountered while building
    all four cloaking pipelines on identical, resource-constrained
    hardware.
\end{itemize}

\section{Related Work}
\label{sec:related}

\subsection{Voice cloning}
Zero-shot TTS systems such as SV2TTS~\cite{jia2018transfer} decouple a
speaker-encoder network (trained for speaker verification) from a
sequence-to-sequence synthesizer, allowing the synthesizer to be
conditioned on a speaker embedding computed from a few seconds of
unseen reference audio. Later systems improved naturalness and speaker
similarity through flow-based and diffusion-based
architectures~\cite{kim2021vits, chen2024f5tts}, and voice-conversion
systems such as Seed-VC decouple linguistic content from timbre so
that a source utterance can be re-rendered in a target speaker's voice
without any text input at all.

\subsection{Inference-time cloaking}
AntiFake~\cite{yu2023antifake} and attack-vc~\cite{huang2021defending}
both perturb a victim recording so that a speaker encoder's embedding
of the perturbed audio is pushed toward a \emph{decoy} speaker and away
from the victim, using an ensemble of surrogate encoders (AntiFake) or
a single voice-conversion model's own encoder/decoder (attack-vc) as
the optimization target. ProtectYourAudio~\cite{liu2023protecting}
instead searches, greedily, over a discrete set of frequency-band
masking operations to find the combination that most degrades a
voice-conversion model's reconstruction while remaining within a
perceptual-similarity budget. All three are inference-time defenses:
their threat model is an attacker who feeds the (cloaked) recording
directly into a cloning system \emph{now}, and the defense's goal is
for that immediate act of cloning to fail.

\subsection{Training-time cloaking}
POP~\cite{zhang2024mitigating} instead targets an attacker who
\emph{fine-tunes} or trains a TTS model on a corpus that includes the
cloaked recording. Building on the "unlearnable examples" line of
work in vision (error-minimizing perturbations that cause a model to
memorize a spurious shortcut rather than learn generalizable
features)~\cite{huang2021unlearnable}, POP crafts a small,
fixed-position perturbation that minimizes a surrogate TTS model's own
reconstruction loss on that exact recording, so that a model
fine-tuned on it overfits to the perturbation rather than learning the
speaker's true acoustic characteristics, degrading the quality of any
\emph{newly synthesized} utterance in that voice.

\section{Threat Model and Trial-Construction Protocol}
\label{sec:threat}

We adopt the standard speaker-verification evaluation protocol used
throughout the cloaking literature and throughout this project: every
trial is a pair of recordings and a binary label indicating whether
they were produced by the same speaker.

\paragraph{Why impostor trials are necessary.}
A verification score in isolation is uninterpretable. Whether a
similarity of, say, 0.65 constitutes ``the same speaker'' or ``a
different speaker'' depends entirely on how that score compares to the
score distribution for pairs that are \emph{known} to be different
speakers. Concretely: a fingerprint lock is not validated by testing it
once with the owner's own finger; it is validated by testing it with
the owner's finger \emph{and} with many other people's fingers, so
that one can measure both the correct-accept rate and the
false-accept rate. Every metric we report in \S\ref{sec:metrics} --
EER, ROC-AUC, minDCF, TAR@FAR, $d'$ -- is a function of \emph{both}
the genuine-trial score distribution and the impostor-trial score
distribution; none of them is definable from genuine trials alone.

\paragraph{Genuine trials.} For every recording $x_i$ in the evaluation
corpus (\S\ref{sec:dataset}) belonging to speaker $s_i$, and its
cloaked counterpart $\hat{x}_i$ produced by a given cloaking technique,
we construct one genuine trial $(\hat{x}_i, x_i, \text{label}=1)$: does
the cloaked recording still verify against the speaker's own real,
uncloaked recording?

\paragraph{Impostor trials.} For the same $\hat{x}_i$, we draw
$k=10$ \emph{other} speakers $s_j \neq s_i$ uniformly at random (with a
fixed seed for reproducibility) from the \emph{entire} 4,867-speaker
pool described in \S\ref{sec:dataset}, and for each drawn speaker
select one of their real, uncloaked recordings $x_j$ uniformly at
random, giving ten impostor trials $(\hat{x}_i, x_j, \text{label}=0)$
per cloaked recording. Every impostor comparison is against
\emph{genuine, unmodified} speech -- the impostor pool is never itself
cloaked -- so that the resulting error rates characterize how the
cloak affects recognizability against the population of real voices an
attacker or verification system would actually encounter.

At the corpus scale used in this paper (4,867 cloaked recordings),
this yields $4{,}867$ genuine trials and $48{,}670$ impostor trials per
verification back-end per cloaking technique -- large enough to
estimate low-false-accept-rate operating points (down to 0.01\% FAR)
with non-degenerate trial counts, which is the reason for scoring each
cloaked clip against ten impostors rather than one.

\section{Dataset}
\label{sec:dataset}

\subsection{Corpus construction and justification}
We assembled a single, unified pool of genuine (``clean\_bonafide'')
speech recordings from three independent public corpora, chosen
specifically so that findings are not an artifact of any one
recording condition, language style, or speaker population:

\begin{itemize}
  \item \textbf{LibriSpeech}~\cite{panayotov2015librispeech}
    (\texttt{train-clean-100}): clean, read audiobook speech,
    included as a near-ideal-recording-condition baseline and because
    it is the most widely used corpus in the TTS/VC literature, making
    our results directly comparable to prior work.
  \item \textbf{ASVspoof 2021}~\cite{yamagishi2021asvspoof} (Logical
    Access, bonafide-labeled trials only, per the official
    \texttt{trial\_metadata.txt} protocol): telephony-channel speech
    that has passed through lossy codec and transmission simulation,
    included to test whether cloaking effectiveness and verification
    leakage are robust to channel degradation the defender does not
    control -- a scenario neither AntiFake, ProtectYourAudio, nor POP's
    own papers evaluate.
  \item \textbf{FakeAVCeleb}~\cite{khalid2021fakeavceleb}: real
    (genuine-audio) recordings from the \texttt{RealVideo-RealAudio}
    and \texttt{FakeVideo-RealAudio} subsets (the latter's audio track
    is untouched even though its video was deepfaked, per the
    dataset's own construction), deduplicated by content hash and
    included specifically for its demographic metadata, absent from
    both other corpora, so that we can measure whether cloaking
    effectiveness and identity leakage are uniform across descent and
    gender groups or systematically weaker for some -- a fairness
    question none of the four original cloaking papers examine.
\end{itemize}

Table~\ref{tab:dataset} reports the exact composition of the resulting
4,867-recording pool, combining the overall per-source counts with the
descent$\times$gender breakdown of FakeAVCeleb -- the only one of the
three source corpora carrying that demographic metadata. The 500
unique FakeAVCeleb speakers are close to evenly split across five
descent groups (African, East Asian, South Asian, Caucasian/American,
Caucasian/European) and two gender labels (men, women); the small
residual imbalance (e.g.\ 404 African-men vs.\ 353 East-Asian-men
recordings) reflects the underlying availability of genuine
(non-deepfaked) audio in the source dataset after deduplication, not
an intentional sampling choice.

\begin{table}[H]
\centering
\caption{Composition of the clean\_bonafide evaluation corpus, with the
FakeAVCeleb subset further broken out by descent and gender.}
\label{tab:dataset}
\begin{tabular}{lrrr}
\toprule
Source / FakeAVCeleb group & Men & Women & Total (\# speakers) \\
\midrule
LibriSpeech (train-clean-100)         & \multicolumn{2}{c}{--} & 500 (207) \\
ASVspoof 2021 (LA, bonafide)          & \multicolumn{2}{c}{--} & 500 (62) \\
\addlinespace
\multicolumn{4}{l}{\textit{FakeAVCeleb, by descent:}} \\
\quad African                    & 404 & 394 & 798 \\
\quad Asian (East)               & 353 & 399 & 752 \\
\quad Asian (South)              & 371 & 395 & 766 \\
\quad Caucasian (American)       & 392 & 391 & 783 \\
\quad Caucasian (European)       & 377 & 391 & 768 \\
\quad \textit{FakeAVCeleb total} & \textit{1{,}897} & \textit{1{,}970} & 3{,}867 (500) \\
\midrule
\textbf{Grand total}             & & & \textbf{4{,}867 (769)} \\
\bottomrule
\end{tabular}
\end{table}

\section{Reference Baselines: Verifier Calibration and the Synthesis Ceiling}
\label{sec:baselines}

Before reporting POP's effect (\S\ref{sec:results}), we first establish
two reference points against which any cloaking result should be read:
(i) how well each of the three verifiers separates genuine from
impostor pairs on completely unmodified, real speech, with no cloaking
and no synthesis involved at all -- the best-case ceiling any
protection scheme is implicitly measured against; and (ii) how much
identity a modern voice-cloning system extracts from \emph{clean,
uncloaked} reference audio -- the synthesis-side ceiling that a
protective cloak (evaluated via the resynthesis protocol of
\S\ref{sec:limitations}) would need to push down.

\subsection{Verifier calibration: clean\_bonafide against itself}
\label{sec:calibration}
Table~\ref{tab:calibration} reports each verifier's EER and TAR@1\%FAR
when genuine and impostor trials are constructed entirely from
unmodified clean\_bonafide recordings (no cloaking, no synthesis),
per source corpus. This calibration pass used an earlier,
independently-sampled trial protocol (\texttt{verify/clean\_bonafide\_metrics.py})
rather than the harmonized one-genuine/ten-impostor-per-clip protocol
of \S\ref{sec:threat} used everywhere else in this paper, so absolute
trial counts (\#T/\#I) differ from Table~\ref{tab:results}; each row's
own EER/TAR/minDCF is nonetheless internally valid as a calibration
reference.

\begin{table}[H]
\centering
\caption{Baseline verifier calibration: clean, unmodified recordings against themselves.}
\label{tab:calibration}
\begin{tabular}{llccc}
\toprule
Corpus & SV & EER & TAR@1\%FAR & \#T/\#I \\
\midrule
LibriSpeech    & SB-ECAPA    & 0.0032 & 0.9977 & 434/1500 \\
LibriSpeech    & Resemblyzer & 0.0250 & 0.9654 & 434/1500 \\
LibriSpeech    & WavLM       & 0.0364 & 0.8802 & 434/1500 \\
\addlinespace
ASVspoof 2021  & SB-ECAPA    & 0.0691 & 0.7429 & 319/1500 \\
ASVspoof 2021  & Resemblyzer & 0.1506 & 0.4890 & 319/1500 \\
ASVspoof 2021  & WavLM       & 0.1289 & 0.4169 & 319/1500 \\
\addlinespace
FakeAVCeleb (overall) & SB-ECAPA    & 0.4036 & 0.1131 & 2864/3000 \\
FakeAVCeleb (overall) & Resemblyzer & 0.3196 & 0.1110 & 2864/3000 \\
FakeAVCeleb (overall) & WavLM       & 0.2732 & 0.1480 & 2864/3000 \\
\bottomrule
\end{tabular}
\end{table}

Two patterns are worth flagging before they resurface in
Section~\ref{sec:results}. First, LibriSpeech -- clean, single-channel,
studio-adjacent read speech -- is trivially separable by all three
verifiers (EER~$\leq$~3.6\%, and SB-ECAPA in particular reaches
EER~=~0.32\% with TAR@1\%FAR~=~99.77\%, the strongest calibration
result of any verifier/corpus pair in this table), while ASVspoof~2021's
telephony-channel degradation and FakeAVCeleb's in-the-wild YouTube
recording conditions substantially compress each verifier's ceiling
(FakeAVCeleb EER 27--40\%): a cloaking technique's absolute EER on
FakeAVCeleb cannot be directly compared to its EER on LibriSpeech,
since the two corpora do not share the same achievable ceiling to
begin with. Second, WavLM and SB-ECAPA are noticeably \emph{weaker}
than Resemblyzer at separating genuine from impostor FakeAVCeleb pairs
in this uncloaked baseline (FakeAVCeleb TAR@1\%FAR: 11--15\% for all
three, versus Resemblyzer's own 48.9\% on the easier ASVspoof~2021
corpus) -- underscoring that FakeAVCeleb is already a hard
verification corpus even before any cloaking is applied.

\subsection{Synthesis ceiling: clean\_bonafide vs.\ voice-cloned speech}
\label{sec:synth-ceiling}
Table~\ref{tab:synth-ceiling} reports the same EER/TAR@1\%FAR/minDCF
metrics when the ``degraded'' side of each trial is instead a
\emph{synthesized clone} of the clean\_bonafide recording (no
cloaking applied to the source), produced by the three voice
cloning/conversion systems of Section~\ref{sec:synth}, following the
identical harmonized trial-construction protocol of
\S\ref{sec:threat} (one genuine + ten impostor trials per synthesized
clip). This is the ceiling a resynthesis-based protective-efficacy
evaluation (\S\ref{sec:limitations}) would need POP, once combined
with a downstream synthesis step, to push \emph{down} from.

\begin{table}[H]
\centering
\caption{clean\_bonafide vs.\ synthesized/converted speech, no cloaking applied (reference ceiling).}
\label{tab:synth-ceiling}
\begin{tabular}{llcccc}
\toprule
System & SV & EER & TAR@1\%FAR & minDCF & \#T/\#I \\
\midrule
SV2TTS (zero-shot TTS)   & SB-ECAPA    & 0.1415 & 0.6007 & 0.7480 & 1{,}500/15{,}000 \\
SV2TTS (zero-shot TTS)   & Resemblyzer & 0.0547 & 0.8453 & 0.4440 & 1{,}500/15{,}000 \\
SV2TTS (zero-shot TTS)   & WavLM       & 0.1667 & 0.3469 & 0.9207 & 1{,}500/15{,}000 \\
\addlinespace
F5-TTS (zero-shot TTS)   & SB-ECAPA    & 0.0166 & 0.9768 & 0.3339 & 4{,}866/48{,}660 \\
F5-TTS (zero-shot TTS)   & Resemblyzer & 0.0223 & 0.9632 & 0.3938 & 4{,}866/48{,}660 \\
F5-TTS (zero-shot TTS)   & WavLM       & 0.0924 & 0.6658 & 0.7668 & 4{,}862/48{,}620 \\
\addlinespace
Seed-VC (voice conversion) & SB-ECAPA    & 0.0311 & 0.9316 & 0.5147 & 4{,}867/48{,}670 \\
Seed-VC (voice conversion) & Resemblyzer & 0.0499 & 0.8436 & 0.6614 & 4{,}867/48{,}670 \\
Seed-VC (voice conversion) & WavLM       & 0.1601 & 0.2509 & 0.9975 & 4{,}867/48{,}670 \\
\bottomrule
\end{tabular}
\end{table}

Two observations frame how POP's numbers in Section~\ref{sec:results}
should be interpreted. First, all three cloning systems leak
substantial identity from \emph{unprotected} audio: F5-TTS in
particular reproduces a target's voice well enough that SB-ECAPA and
Resemblyzer both exceed 96\% TAR@1\%FAR against clean, uncloaked
reference recordings -- this is the real-world threat any of the four
cloaking techniques in this paper is meant to disrupt. Second, the
same cross-verifier ordering seen throughout this project reappears
here: WavLM is consistently the hardest of the three to fool by a
synthesized clone (lowest TAR@1\%FAR for all three systems), and
Resemblyzer or SB-ECAPA the easiest, depending on the system -- a
pattern any future resynthesis-based POP evaluation should expect to
inherit, independent of whether the source audio was cloaked.

\section{Speaker Verification Models}
\label{sec:verifiers}

Every genuine/impostor trial (\S\ref{sec:threat}) is scored
independently by three speaker-verification back-ends, never combined
or ensembled into a single score: each back-end produces its own
embedding space, its own similarity score, and consequently its own
independently computed EER/ROC-AUC/minDCF/TAR@FAR/$d'$. Reporting three
back-ends side by side, rather than one, is deliberate: a cloaking
technique that fools one embedding space but not another would be
substantially over- or under-stated by any single-model evaluation, a
methodological gap in several of the papers we build on (each of which
evaluates against only one or two verifiers of their own choosing).

\begin{table}[H]
\centering
\caption{Speaker verification back-ends used for genuine/impostor scoring.}
\label{tab:verifiers}
\begin{tabular}{p{2.4cm}p{2.6cm}p{7.2cm}}
\toprule
Model & Embedding & Rationale \\
\midrule
SpeechBrain ECAPA-TDNN (SB-ECAPA)~\cite{desplanques2020ecapa}
  & 192-d, channel-attention TDNN, trained on VoxCeleb
  & The current de facto standard open-source speaker-verification
    architecture; used as the (white-box or surrogate) target encoder
    by both AntiFake and attack-vc, making it directly comparable to
    the original papers' own evaluations. \\
Resemblyzer (GE2E d-vector)
  & 256-d LSTM d-vector, generalized end-to-end loss
  & Lightweight, widely used as the surrogate encoder inside
    ProtectYourAudio's own vendored pipeline and in AUTOVC-style
    voice-conversion systems; included so our evaluation covers the
    exact embedding family several of the cloaking methods were
    originally tuned against. \\
WavLM-based x-vector~\cite{chen2022wavlm}
  & 512-d x-vector head over WavLM-base-plus self-supervised features
  & A self-supervised-pretraining-based verifier, architecturally
    unrelated to the CNN/TDNN and LSTM encoders above; included to
    test whether cloaking effectiveness transfers to a verifier family
    none of the four cloaking methods were designed or tuned against
    (a black-box generalization test in the spirit of AntiFake's own
    ensemble-transferability argument, but applied adversarially in
    the other direction, from the defender's side). \\
\bottomrule
\end{tabular}
\end{table}

Across every system-vs-clean\_bonafide comparison run in this project
to date (Section~\ref{sec:results} and our companion TTS/VC
evaluations), WavLM has consistently been the hardest of the three to
fool (largest residual EER, lowest TAR@1\%FAR degradation) and
Resemblyzer the easiest, with SB-ECAPA falling in between -- a pattern
we revisit when interpreting the POP results in
Section~\ref{sec:results}.

\section{Voice Cloning and Conversion Systems}
\label{sec:synth}

Independently of the cloaking evaluation in this paper, our broader
project also measures whether \emph{uncloaked} speech survives
resynthesis by contemporary voice-cloning systems, to calibrate how
much identity a modern cloning system can extract from clean reference
audio in the first place -- the ceiling against which any cloaking
defense's protective effect must be read. Table~\ref{tab:synth} lists
the three systems used for that purpose and the rationale for each.

\begin{table}[H]
\centering
\caption{Voice cloning / conversion systems used for the (uncloaked) clean\_bonafide-vs-synthesis calibration.}
\label{tab:synth}
\begin{tabular}{p{2.2cm}p{1.6cm}p{8.4cm}}
\toprule
System & Type & Rationale \\
\midrule
SV2TTS~\cite{jia2018transfer}
  & Zero-shot TTS
  & The original, foundational ``speaker encoder + Tacotron-style
    synthesizer'' zero-shot cloning architecture that both AntiFake and
    ProtectYourAudio explicitly target in their own evaluations;
    included as the historical/baseline cloning threat. \\
F5-TTS~\cite{chen2024f5tts}
  & Zero-shot TTS (flow-matching)
  & A current state-of-the-art diffusion/flow-matching zero-shot TTS
    system, substantially stronger than SV2TTS at speaker-similarity
    and naturalness; included so cloaking effectiveness is measured
    against a modern, not just historical, cloning threat. \\
Seed-VC
  & Zero-shot voice conversion
  & Unlike the two TTS systems above, a voice-\emph{conversion} system
    (content and target-timbre audio in, converted audio out, no text
    required); included so our evaluation is not confined to the
    text-to-speech threat model, since several cloaking defenses
    (attack-vc, ProtectYourAudio's AutoVC/Chou baselines) are evaluated
    against voice conversion specifically, not TTS. \\
\bottomrule
\end{tabular}
\end{table}

All three systems are also used for the resynthesis-based
protective-efficacy evaluation of Section~\ref{sec:protective-efficacy}
-- cloning \emph{from POP-cloaked source audio} with the identical
systems and protocol used here to establish the clean-audio ceiling,
so the two results are directly comparable; the same three systems are
used identically for the corresponding evaluations of the other three
cloaking techniques in this broader project
(Section~\ref{sec:protective-efficacy-attackvc} for attack-vc).

\paragraph{Fixed synthesis prompt.} For both TTS systems (SV2TTS,
F5-TTS), every single cloning call in this project -- for every
recording, cloaked or uncloaked, across all four cloaking techniques
-- uses the exact same fixed target text, so that the only thing that
varies between conditions is the cloned \emph{identity}, never the
linguistic content being spoken:

\begin{quote}
\itshape
``This is a test of the voice cloning system used to evaluate audio
protection methods. Please listen carefully and notice how closely
the synthesized voice matches the original speaker across different
words, tones, and sentence rhythms. Every recording, whether real or
generated, carries subtle traces of the speaker's identity that these
systems are designed to detect.''
\end{quote}

Seed-VC, being a voice-\emph{conversion} rather than text-to-speech
system, takes no text prompt at all; it instead uses a single fixed
source recording (held constant across every clean\_bonafide/cloaked
target, for the same content-held-constant reason as the text prompt
above) to supply the linguistic content, with only the target-timbre
input varying per recording.

\section{Cloaking Techniques}
\label{sec:cloaking}

We implemented all four cloaking techniques end to end, from each
technique's official public source code, on identical hardware (a
single 4~GB NVIDIA vGPU partition), rather than re-using each paper's
own reported numbers. This section documents the mechanism and
rationale for including each technique, and the practical obstacles
encountered building each pipeline, since these materially affect what
can and cannot be evaluated on a given hardware budget.

\begin{table}[H]
\centering
\caption{The four cloaking techniques evaluated in this project.}
\label{tab:cloaking}
\begin{tabular}{p{1.8cm}p{2.3cm}p{7.9cm}}
\toprule
Technique & Threat model & Mechanism and rationale for inclusion \\
\midrule
POP~\cite{zhang2024mitigating}
  & Training-time (unlearnable audio)
  & PGD perturbation minimizing a pretrained VITS surrogate's own
    mel-reconstruction loss on the exact clip, at a fixed audio patch
    position; included as the only training-time (rather than
    inference-time) defense in our study, and the subject of this
    paper. \\
ProtectYourAudio~\cite{liu2023protecting}
  & Inference-time (voice-conversion spoofing)
  & Greedy discrete search over frequency-band masking operations
    (zeroing, additive noise, Gaussian blur) against a Chou-style
    AdaIN-VC surrogate; included as a non-gradient-based, black-box-
    search alternative to the gradient-based methods below, and
    because it explicitly optimizes for the perturbed audio to
    \emph{remain} usable/intelligible, unlike a pure adversarial
    attack. \\
attack-vc~\cite{huang2021defending}
  & Inference-time (voice-conversion spoofing)
  & Decoy-target PGD attack (the first published adversarial attack on
    voice conversion) with three variants (end-to-end, embedding,
    feedback); we use the paper's own best-performing and most
    efficient variant, the embedding attack, which perturbs the input
    directly against the target model's speaker encoder without
    requiring a full encoder-decoder forward pass. \\
AntiFake~\cite{yu2023antifake}
  & Inference-time (TTS/VC spoofing, ensemble)
  & Decoy-target PGD attack optimized jointly against three surrogate
    encoders (RTVC, AutoVC, Coqui YourTTS) simultaneously, for
    black-box transferability; included as the most heavyweight
    (three-encoder ensemble) of the four techniques. On the original
    4~GB vGPU this project started on, loading all three encoders
    simultaneously consumed $\approx$2.9--3.0~GB before a single
    optimization step ran, reliably exhausting available memory; we
    verified this was a hardware-capacity limitation, not a software
    defect, by reproducing the identical failure on an idle GPU. Cloak
    generation became feasible once the underlying hardware was
    upgraded to a 12~GB vGPU partition; cloaking is now complete for
    4,844 of the 4,867-recording corpus ($\approx$8 minutes/clip, all
    1,000 attack iterations against all three encoders; 21 permanent
    per-clip failures, 2 recordings never attempted), reported in
    Section~\ref{sec:results-antifake}. \\
\bottomrule
\end{tabular}
\end{table}

\subsection{Engineering obstacles common to all four pipelines}
Beyond the AntiFake memory ceiling above, reproducing each technique
end to end (rather than relying on the original papers' released
model checkpoints and reported numbers) surfaced several
implementation-level issues worth recording for reproducibility:
(i) POP's own reference implementation silently \emph{drops} any
training example whose phonemized transcript exceeds a fixed length
threshold from its data loader, with no warning and no error --
undetected, this stalls an unattended batch job indefinitely, since
the affected examples are never marked complete; we patched this by
raising the threshold to cover legitimate long utterances and
explicitly excluding the small number of genuine transcription
failures (Whisper hallucinations on short clips) that remained.
(ii) None of the four techniques' reference implementations are
robust to a shared, memory-constrained GPU: cuDNN/cuFFT internal
errors and CUDA out-of-memory failures occurred transiently under
concurrent load from multiple cloaking jobs and had to be
disambiguated from genuine per-file failures before being treated as
permanent. (iii) Several dependencies (an espeak-ng phonemizer
backend, in particular) required sudo-free workarounds on
infrastructure without administrator access.

\section{Evaluation Metrics}
\label{sec:metrics}

We report two families of metric, corresponding to the two
conceptually distinct properties introduced in
Section~\ref{sec:intro}: verification-leakage metrics (computed from
the genuine/impostor trial scores of \S\ref{sec:threat}, once per
verifier in Table~\ref{tab:verifiers}) and signal-quality metrics
(computed directly from the reference and cloaked waveforms, verifier-
and trial-independent). Table~\ref{tab:metrics-overview} summarizes
both families before we define each metric precisely below.

\begin{table}[H]
\centering
\caption{Overview of the two evaluation-metric families.}
\label{tab:metrics-overview}
\begin{tabular}{p{2.6cm}p{4.6cm}p{5.8cm}}
\toprule
Family & Metrics & What it measures \\
\midrule
Verification leakage & EER, ROC-AUC, pAUC, minDCF, TAR@FAR, $d'$, overlap coefficient
  & Whether the cloaked recording is still recognized as the same speaker,
    relative to real impostors, by each of the three verifiers. \\
Signal quality & STOI, PESQ, SI-SDR
  & How much the cloaking perturbation changed the waveform itself,
    independent of any speaker-identity model. \\
\bottomrule
\end{tabular}
\end{table}

\subsection{Verification-leakage metrics}

Let genuine trial scores $\{t_k\}$ and impostor trial scores $\{i_k\}$
(\S\ref{sec:threat}) be the cosine similarities produced by a given
verifier for all trials involving a given cloaking technique. Let
$\mathrm{FAR}(\tau)$ and $\mathrm{TAR}(\tau) = 1-\mathrm{FRR}(\tau)$ be
the false-accept and true-accept rates at decision threshold $\tau$,
obtained by sweeping $\tau$ over the full receiver operating
characteristic (ROC) curve.

\paragraph{Equal Error Rate (EER).} The threshold $\tau^\ast$ at which
$\mathrm{FAR}(\tau^\ast) = \mathrm{FRR}(\tau^\ast)$; we report the
common value at that point. A lower EER means the two score
distributions are better separated (genuine and impostor trials are
easier to tell apart at some threshold); in the context of a cloak
that is \emph{supposed} to remain imperceptible, a low EER between
cloaked and genuine recordings indicates the cloak is behaving as
designed (still recognizable), whereas we would expect a
\emph{resynthesis}-based evaluation of a successfully protective
cloak to show the opposite: a high EER (or low TAR@FAR) between the
resulting synthetic clone and the genuine speaker.

\paragraph{TAR@FAR.} The true-accept rate at a fixed, operationally
meaningful false-accept rate (we report 1\%, 0.1\%, and 0.01\% FAR,
obtained by interpolating the ROC curve), the standard way to
summarize verification performance at security-relevant, low-FAR
operating points rather than at a single, arbitrarily chosen
threshold.

\paragraph{Minimum normalized Detection Cost Function (minDCF).}
\begin{equation}
C_{\mathrm{det}}(\tau) = C_{\mathrm{miss}} \cdot P_{\mathrm{target}} \cdot P_{\mathrm{miss}}(\tau)
  + C_{\mathrm{fa}} \cdot (1 - P_{\mathrm{target}}) \cdot \mathrm{FAR}(\tau)
\end{equation}
with $P_{\mathrm{miss}}(\tau) = 1 - \mathrm{TAR}(\tau)$, minimized over
$\tau$ and normalized by the cost of the better trivial
(always-accept / always-reject) baseline. We use
$P_{\mathrm{target}}=0.05$, $C_{\mathrm{miss}}=1.0$,
$C_{\mathrm{fa}}=10.0$, reflecting an identity-security deployment
where a false accept (an impostor, or an attacker's clone, passing as
the real speaker) is an order of magnitude costlier than a false
reject.

\paragraph{Cohen's $d'$ (sensitivity index).}
\begin{equation}
d' = \frac{\mu_{\mathrm{genuine}} - \mu_{\mathrm{impostor}}}
          {\sqrt{\tfrac{1}{2}\left(\sigma^2_{\mathrm{genuine}} + \sigma^2_{\mathrm{impostor}}\right)}}
\end{equation}
the standardized distance between the genuine- and impostor-score
distributions' means; larger $|d'|$ indicates better separation
regardless of any particular decision threshold.

\paragraph{Overlap coefficient.} The shared area between the genuine-
and impostor-score histograms (shared bin edges over the pooled score
range), ranging from 0 (perfectly separated) to 1 (identical
distributions); an intuitive, threshold-free complement to $d'$.

\paragraph{Full and partial ROC-AUC.} The full ROC-AUC is
$P(\text{genuine score} > \text{impostor score})$ for a randomly drawn
pair; we report it only as a supplementary figure since it is
dominated by the high-FAR region of the curve and can look
misleadingly strong even when the security-relevant low-FAR region is
weak. We additionally report the partial AUC restricted to
FAR~$\le$~1\% and FAR~$\le$~0.1\%, normalized to a $[0,1]$ scale where
0.5 is chance-level separation in that region, as the more decision-
relevant summary.

\paragraph{Bootstrap confidence intervals.} For EER and TAR@1\%FAR we
report 95\% confidence intervals from 1{,}000 bootstrap resamples,
resampling genuine and impostor trials separately (preserving each
class's sample size) with replacement, the standard approach for
verification-trial bootstrapping.

\subsection{Signal-quality metrics}
\label{sec:quality-metrics}

Unlike the leakage metrics above, the following three metrics compare
the reference (genuine) and cloaked waveforms of the \emph{same}
recording directly, sample-aligned, and do not involve any of the
three verifiers in Table~\ref{tab:verifiers} at all -- a common
misconception we clarify explicitly given it arose during internal
review of this evaluation. They are consequently identical across all
three verifier rows of our results table (Table~\ref{tab:results}):
they were never computed per-verifier to begin with.

\paragraph{Short-Time Objective Intelligibility (STOI).} Compares the
time-frequency envelope correlation of the reference and degraded
signal in short overlapping frames, producing a score in
$[0,1]$ (higher is more intelligible); computed with the
\texttt{pystoi} reference implementation.

\paragraph{Perceptual Evaluation of Speech Quality (PESQ).} An
ITU-T~P.862 standard perceptual model, originally developed for
telephony codec quality assessment, producing a score roughly in
$[-0.5, 4.5]$ (higher is better perceived quality); we use wideband
mode at 16~kHz, computed with the \texttt{pesq} reference
implementation.

\paragraph{Scale-Invariant Signal-to-Distortion Ratio (SI-SDR).}
\begin{equation}
\alpha = \frac{\langle \hat{x}, x \rangle}{\lVert x \rVert^2}, \qquad
\mathrm{SI\text{-}SDR} = 10 \log_{10}
  \frac{\lVert \alpha x \rVert^2}{\lVert \hat{x} - \alpha x \rVert^2}
\end{equation}
where $x$ is the mean-removed reference signal and $\hat{x}$ the
mean-removed cloaked signal; $\alpha x$ is the projection of $\hat{x}$
onto $x$, so the ratio compares the energy of the (optimally scaled)
reference component against the residual distortion, in dB (higher is
less distortion). Self-implemented (no external dependency required),
following the standard formulation used throughout the source
separation and speech enhancement literature.

Because clean\_bonafide is natively sampled at 16~kHz and POP's VITS
surrogate operates at 24~kHz, we resample the cloaked output back to
16~kHz before computing any of the three metrics above (PESQ, in
particular, only supports exactly 8~kHz or 16~kHz, making this
resampling step not optional).

\section{Experimental Setup for the POP Evaluation}
\label{sec:setup}

For every one of the 4,867 recordings in Table~\ref{tab:dataset}, we
applied POP's official attack (surrogate: the paper's own released
LJSpeech-pretrained VITS checkpoint; $\varepsilon=8/255$; 200 PGD
iterations; fixed-position patch mode, matching the paper's own
default configuration exactly) to produce a cloaked counterpart. Where
no official transcript existed for a recording (ASVspoof~2021 and
FakeAVCeleb, neither of which ships transcripts; only LibriSpeech
does), we generated one with Whisper (\texttt{openai/whisper-base})
ASR, since POP's surrogate is a text-conditioned model and requires a
transcript to compute its text/audio alignment. We then scored every
cloaked recording against its own genuine original (1 trial) and ten
randomly drawn impostors from elsewhere in the 4,867-recording pool
(10 trials) with each of the three verifiers in
Table~\ref{tab:verifiers}, and additionally computed STOI, PESQ, and
SI-SDR directly between each cloaked recording and its genuine
original.

Of the 4,867 recordings, 3 were excluded from the attack corpus after
manual inspection revealed their Whisper-generated transcripts were
hallucinations (250--515 phonemized characters per second of audio,
against a corpus median of 17 characters per second -- physically
impossible speaking rates indicating ASR failure rather than a
legitimately long utterance); the remaining 4,864 were cloaked and
evaluated successfully.

\section{Results}
\label{sec:results}

Table~\ref{tab:results} reports the full verification-leakage and
signal-quality metrics for clean\_bonafide vs.\ POP-cloaked, pooled
across all three source corpora (4,864 cloaked recordings, 53,503
SB-ECAPA trials, 53,504 Resemblyzer trials, and 44,154 WavLM trials --
WavLM's trial count is lower because a subset of clips required a
rescoring pass after transient CUDA out-of-memory failures during the
original run, described in \S\ref{sec:limitations}). Every verifier
reports an EER below 0.2\% and a TAR@1\%FAR of 100.00\%: POP-cloaked
audio is, by every verifier used in this paper, indistinguishable from
its genuine source recording. STOI (0.971), PESQ (2.92), and SI-SDR
(16.3~dB) confirm this acoustically -- the perturbation is a
high-fidelity, low-audibility degradation, not a gross corruption of
the signal.

\begin{table}[H]
\centering
\caption{clean\_bonafide vs.\ POP-cloaked: verification leakage and
signal-quality metrics, pooled across all three source corpora.}
\label{tab:results}
\begin{tabular}{lccccccc}
\toprule
SV & EER & TAR@1\% & minDCF & $d'$ & STOI & PESQ & SI-SDR (dB) \\
\midrule
SB-ECAPA     & 0.0007 & 1.0000 & 0.1743 & 14.407 & \multirow{3}{*}{0.9709} & \multirow{3}{*}{2.9218} & \multirow{3}{*}{16.2909} \\
Resemblyzer  & 0.0014 & 1.0000 & 0.1896 & 8.0265 & & & \\
WavLM        & 0.0016 & 1.0000 & 0.1963 & 3.1434 & & & \\
\bottomrule
\end{tabular}
\end{table}

Two further observations follow from the full table
(Table~\ref{tab:results} shows the pooled figures; per-verifier $d'$
values are reported individually because, unlike EER/TAR@1\%FAR, $d'$
does not saturate and is therefore the more informative
verifier-to-verifier comparison here). First, $d'$ decreases
monotonically from SB-ECAPA (14.41) to Resemblyzer (8.03) to WavLM
(3.14) even though all three verifiers agree on EER and TAR@1\%FAR --
WavLM's genuine/impostor score distributions are closer together in
absolute terms, it is simply not close enough for this to change the
verification decision at any practical operating point. Second, the
near-ceiling STOI/PESQ/SI-SDR values corroborate that POP achieves
this near-perfect verification-score preservation without a
correspondingly large perceptible change to the waveform -- consistent
with an $\varepsilon=8/255$ perturbation budget designed specifically
to stay below the threshold of casual audibility.

\subsection{Illustrative single-trial example}
To make the trial-construction protocol concrete before the aggregate
statistics above are finalized, Table~\ref{tab:example} reports one
real, representative genuine/impostor trial set from the in-progress
run: the ASVspoof~2021 recording \texttt{asvspoof\_0000}, cloaked by
POP, scored by SB-ECAPA against its own genuine original (the genuine
trial) and ten other, randomly drawn real speakers (the impostor
trials).

\begin{table}[H]
\centering
\caption{One real genuine/impostor trial set (SB-ECAPA, POP-cloaked \texttt{asvspoof\_0000}).}
\label{tab:example}
\begin{tabular}{lcc}
\toprule
Trial & Score & Label \\
\midrule
Genuine (own speaker)      & 0.937 & 1 \\
Impostor 1 & 0.163 & 0 \\
Impostor 2 & 0.102 & 0 \\
Impostor 3 & 0.097 & 0 \\
Impostor 4 & 0.167 & 0 \\
Impostor 5 & 0.017 & 0 \\
Impostor 6 & 0.131 & 0 \\
Impostor 7 & 0.312 & 0 \\
Impostor 8 & 0.125 & 0 \\
Impostor 9 & 0.006 & 0 \\
Impostor 10 & 0.008 & 0 \\
\bottomrule
\end{tabular}
\end{table}

For this individual trial, the genuine score (0.937) is well separated
from all ten impostor scores (0.006--0.312), consistent with POP's
design goal of imperceptibility: SB-ECAPA still confidently recognizes
the cloaked recording as belonging to its true speaker. Whether this
pattern holds in aggregate across all 4,864 recordings and all three
verifiers, and how it compares across the ten FakeAVCeleb demographic
groups, is reported in Table~\ref{tab:results} and the per-group
breakdown of Table~\ref{tab:ablation} below.

\subsection{Ablation: verification leakage by descent $\times$ gender}
\label{sec:ablation}
Table~\ref{tab:ablation} restricts the same trials to the
FakeAVCeleb-sourced subset only (the ten descent$\times$gender groups
of Table~\ref{tab:dataset}; LibriSpeech and ASVspoof~2021 entries carry
no such metadata and are excluded from this breakdown, not folded into
a fake ``unknown'' group) and reports each group's own EER and
TAR@1\%FAR per verifier.

\begin{table}[H]
\centering
\caption{clean\_bonafide vs.\ POP-cloaked: EER / TAR@1\%FAR by descent
$\times$ gender (FakeAVCeleb subset).}
\label{tab:ablation}
\begin{tabular}{lcccccc}
\toprule
& \multicolumn{2}{c}{SB-ECAPA} & \multicolumn{2}{c}{Resemblyzer} & \multicolumn{2}{c}{WavLM} \\
Group & EER & TAR@1\% & EER & TAR@1\% & EER & TAR@1\% \\
\midrule
African (men)               & 0.0009 & 1.0000 & 0.0010 & 1.0000 & 0.0006 & 1.0000 \\
African (women)              & 0.0010 & 1.0000 & 0.0004 & 1.0000 & 0.0023 & 1.0000 \\
Asian, East (men)            & 0.0010 & 1.0000 & 0.0011 & 1.0000 & 0.0003 & 0.9466 \\
Asian, East (women)          & 0.0008 & 1.0000 & 0.0008 & 1.0000 & 0.0007 & 1.0000 \\
Asian, South (men)           & 0.0004 & 1.0000 & 0.0005 & 1.0000 & 0.0028 & 1.0000 \\
Asian, South (women)         & 0.0004 & 1.0000 & 0.0011 & 1.0000 & 0.0006 & 1.0000 \\
Caucasian, American (men)    & 0.0004 & 1.0000 & 0.0008 & 1.0000 & 0.0007 & 1.0000 \\
Caucasian, American (women)  & 0.0004 & 1.0000 & 0.0022 & 1.0000 & 0.0003 & 1.0000 \\
Caucasian, European (men)    & 0.0003 & 0.9650 & 0.0012 & 1.0000 & 0.0017 & 1.0000 \\
Caucasian, European (women)  & 0.0004 & 1.0000 & 0.0005 & 0.9988 & 0.0003 & 1.0000 \\
\bottomrule
\end{tabular}
\end{table}

Every group, verifier, and gender combination lands in the same
sub-0.3\% EER / near-100\% TAR@1\%FAR regime as the pooled result --
POP's imperceptibility holds uniformly across descent and gender in
this dataset, with no demographic group standing out as
disproportionately easier or harder to re-identify after cloaking. The
only two departures from a perfect TAR@1\%FAR of 100\% are
Caucasian(European)-men under SB-ECAPA (96.50\%) and
Asian(East)-men under WavLM (94.66\%); both remain well above the
FakeAVCeleb calibration baseline's TAR@1\%FAR for the same verifiers
(11--15\%, Table~\ref{tab:calibration} being the LibriSpeech/ASVspoof/
overall-FakeAVCeleb calibration and \S\ref{sec:calibration}'s
discussion applying equally to these two per-group figures), so neither
represents a meaningful protection effect -- both are still far closer
to ``trivially re-identifiable'' than to the uncloaked FakeAVCeleb
ceiling.

\section{Results: attack-vc}
\label{sec:results-attackvc}

Table~\ref{tab:results-attackvc} reports the same verification-leakage
and signal-quality metrics as Table~\ref{tab:results}, this time for
clean\_bonafide vs.\ attack-vc-cloaked (4,863 cloaked recordings --
4,863 of the 4,867-recording corpus succeeded; the remaining 4 are
permanently excluded, analogous to POP's 3 Whisper-hallucination
exclusions -- 53,493 SB-ECAPA/Resemblyzer trials and 53,471 WavLM
trials). As with POP, attack-vc's embedding attack is designed for
imperceptibility rather than resynthesis-time protection
(\S\ref{sec:limitations} applies here identically), and the results
bear that out: EER is below 2.1\% and TAR@1\%FAR above 96.9\% for
every verifier.

\begin{table}[H]
\centering
\caption{clean\_bonafide vs.\ attack-vc-cloaked: verification leakage
and signal-quality metrics, pooled across all three source corpora.}
\label{tab:results-attackvc}
\begin{tabular}{lccccccc}
\toprule
SV & EER & TAR@1\% & minDCF & $d'$ & STOI & PESQ & SI-SDR (dB) \\
\midrule
SB-ECAPA     & 0.0010 & 0.9996 & 0.2235 & 12.3957 & \multirow{3}{*}{0.5472} & \multirow{3}{*}{3.5803} & \multirow{3}{*}{$-$33.9665} \\
Resemblyzer  & 0.0008 & 0.9996 & 0.2075 & 7.8385  & & & \\
WavLM        & 0.0208 & 0.9688 & 0.3464 & 2.9521  & & & \\
\bottomrule
\end{tabular}
\end{table}

Two things distinguish attack-vc from POP here. First, WavLM's EER
(2.08\%) is roughly $13\times$ higher than POP's WavLM EER (0.16\%),
and its TAR@1\%FAR (96.88\%) is the only sub-99\% pooled figure across
both cloaking methods -- attack-vc's imperceptibility is real but
measurably less complete than POP's, at least against WavLM
specifically. Second, and more strikingly, the signal-quality metrics
diverge sharply from POP's: STOI drops from POP's 0.971 to 0.547, and
SI-SDR collapses from $+16.29$~dB to $-33.97$~dB, even though PESQ
actually \emph{improves} slightly (2.92 $\rightarrow$ 3.58) and every
verifier still confirms near-perfect speaker-identity preservation.
We read this apparent contradiction as a measurement artifact of
attack-vc's synthesis path rather than evidence of a grossly degraded
signal: unlike POP, which perturbs the waveform directly, attack-vc
perturbs a mel-spectrogram representation and reconstructs the
waveform through AdaIN-VC's vocoder, which does not guarantee
sample-level phase or gain alignment with the original waveform.
SI-SDR is defined via a sample-aligned projection onto the reference
signal (\S\ref{sec:quality-metrics}) and is accordingly acutely
sensitive to exactly this kind of phase/scale mismatch, whereas PESQ
and STOI are perceptually motivated and comparatively tolerant of
it -- consistent with PESQ actually scoring attack-vc \emph{higher}
than POP despite SI-SDR scoring it far lower. We have not yet
confirmed this hypothesis by direct waveform inspection; we flag it
here as the leading explanation rather than a verified fact, and
recommend it as a target for follow-up verification before drawing
conclusions from the SI-SDR figure alone.

\subsection{Ablation: verification leakage by descent $\times$ gender}
\label{sec:ablation-attackvc}
Table~\ref{tab:ablation-attackvc} reports the same per-group breakdown
as Table~\ref{tab:ablation}, restricted to the FakeAVCeleb subset,
for attack-vc.

\begin{table}[H]
\centering
\caption{clean\_bonafide vs.\ attack-vc-cloaked: EER / TAR@1\%FAR by
descent $\times$ gender (FakeAVCeleb subset).}
\label{tab:ablation-attackvc}
\begin{tabular}{lcccccc}
\toprule
& \multicolumn{2}{c}{SB-ECAPA} & \multicolumn{2}{c}{Resemblyzer} & \multicolumn{2}{c}{WavLM} \\
Group & EER & TAR@1\% & EER & TAR@1\% & EER & TAR@1\% \\
\midrule
African (men)               & 0.0006 & 0.9977 & 0.0006 & 1.0000 & 0.0033 & 1.0000 \\
African (women)              & 0.0009 & 1.0000 & 0.0006 & 1.0000 & 0.0013 & 0.9964 \\
Asian, East (men)            & 0.0010 & 1.0000 & 0.0006 & 1.0000 & 0.0028 & 1.0000 \\
Asian, East (women)          & 0.0009 & 0.9908 & 0.0003 & 1.0000 & 0.0046 & 0.9975 \\
Asian, South (men)           & 0.0012 & 1.0000 & 0.0009 & 1.0000 & 0.0072 & 0.9919 \\
Asian, South (women)         & 0.0004 & 0.9623 & 0.0003 & 1.0000 & 0.0057 & 0.9975 \\
Caucasian, American (men)    & 0.0004 & 0.9976 & 0.0006 & 0.9929 & 0.0064 & 0.9906 \\
Caucasian, American (women)  & 0.0004 & 1.0000 & 0.0012 & 1.0000 & 0.0031 & 1.0000 \\
Caucasian, European (men)    & 0.0007 & 0.9950 & 0.0023 & 1.0000 & 0.0023 & 0.9975 \\
Caucasian, European (women)  & 0.0006 & 1.0000 & 0.0006 & 1.0000 & 0.0013 & 1.0000 \\
\bottomrule
\end{tabular}
\end{table}

The comparison to POP's per-group breakdown (Table~\ref{tab:ablation})
is instructive. POP's TAR@1\%FAR is exactly 100\% in 28 of its 30
group$\times$verifier cells, with only two groups dropping (both to
$\geq$94.7\%). attack-vc's per-group TAR@1\%FAR is more scattered: 12
of its 30 cells fall below 100\% (range 96.23--99.77\%), spread
unevenly across verifiers -- Resemblyzer drops in only 1 of its 10
groups, SB-ECAPA in 5, and WavLM in 6 (WavLM still reaches exactly
100\% in the remaining 4 groups: African-men, Asian(East)-men,
Caucasian(American)-women, and Caucasian(European)-women). This is
consistent with the pooled finding above: attack-vc's imperceptibility,
while still strong in absolute terms, is measurably less uniform
across both verifiers and demographic groups than POP's. As with
POP's two outlier groups (\S\ref{sec:ablation}), every attack-vc
group$\times$verifier figure remains far closer to POP's near-ceiling
regime than to the uncloaked FakeAVCeleb calibration baseline
(TAR@1\%FAR 11--15\%, Table~\ref{tab:calibration}), so none of this
scatter constitutes a meaningful protective effect either -- it is
noise within an already-near-perfect imperceptibility result, not a
sign that any particular group or verifier is disproportionately
protected.

\section{Results: AntiFake}
\label{sec:results-antifake}

Table~\ref{tab:results-antifake} reports the same verification-leakage
and signal-quality metrics as Tables~\ref{tab:results}/\ref{tab:results-attackvc},
this time for clean\_bonafide vs.\ AntiFake-cloaked, pooled across all
three source corpora (4,844 of the 4,867-recording corpus succeeded --
21 permanent per-clip failures and 2 recordings never attempted;
53,284 trials per verifier, uniformly across all three, with none of
POP's WavLM rescoring-driven shortfall). Unlike POP and
attack-vc, AntiFake's decoy-target attack does \emph{not} achieve
imperceptibility on this corpus: EER is 22--66\% and TAR@1\%FAR is only
5--25\%, roughly two orders of magnitude worse than either of the other
two techniques.

\begin{table}[H]
\centering
\caption{clean\_bonafide vs.\ AntiFake-cloaked: verification leakage and
signal-quality metrics, pooled across all three source corpora.}
\label{tab:results-antifake}
\begin{tabular}{lccccccc}
\toprule
SV & EER & TAR@1\% & minDCF & $d'$ & STOI & PESQ & SI-SDR (dB) \\
\midrule
SB-ECAPA     & 0.4781 & 0.2120 & 0.9422 & 0.3613 & \multirow{3}{*}{0.7031} & \multirow{3}{*}{1.4642} & \multirow{3}{*}{1.9004} \\
Resemblyzer  & 0.6587 & 0.0508 & 0.9833 & $-$0.5421 & & & \\
WavLM        & 0.2229 & 0.2519 & 0.9717 & 1.5477 & & & \\
\bottomrule
\end{tabular}
\end{table}

The signal-quality metrics corroborate this directly, and -- unlike
attack-vc's STOI/PESQ contradiction (\S\ref{sec:results-attackvc}) --
they agree with each other and with the verification scores. STOI
(0.703) sits between POP's 0.971 and attack-vc's 0.547; PESQ (1.46) is
the \emph{lowest} of all three techniques, below both POP's 2.92 and
attack-vc's 3.58; SI-SDR (1.9~dB) is far below POP's 16.3~dB, though
nowhere near attack-vc's $-$34.0~dB anomaly. Every metric here points
the same direction: AntiFake's perturbation is genuinely more audible
and more disruptive to the raw waveform than either other technique,
consistent with a decoy-target attack that optimizes directly for
embedding displacement rather than for staying below a
perceptibility budget the way POP's fixed $\ell_\infty$ patch or
attack-vc's mel-domain perturbation do. Resemblyzer -- the encoder
family AntiFake's own ensemble does \emph{not} include as a surrogate
(RTVC, AutoVC, and Coqui YourTTS are the three targets;
\S\ref{sec:cloaking}) -- shows the worst imperceptibility of the three
verifiers (EER 65.9\%), while WavLM, architecturally unrelated to all
three surrogates, is comparatively the least affected (EER 22.3\%) --
the opposite ordering from POP and attack-vc, where WavLM was
consistently the \emph{hardest} to fool and Resemblyzer among the
easiest (\S\ref{sec:verifiers}).

\subsection{Ablation: verification leakage by descent $\times$ gender}
\label{sec:ablation-antifake}
Table~\ref{tab:ablation-antifake} reports the same per-group breakdown
as Tables~\ref{tab:ablation}/\ref{tab:ablation-attackvc}, restricted to
the FakeAVCeleb subset, for AntiFake.

\begin{table}[H]
\centering
\caption{clean\_bonafide vs.\ AntiFake-cloaked: EER / TAR@1\%FAR by
descent $\times$ gender (FakeAVCeleb subset).}
\label{tab:ablation-antifake}
\begin{tabular}{lcccccc}
\toprule
& \multicolumn{2}{c}{SB-ECAPA} & \multicolumn{2}{c}{Resemblyzer} & \multicolumn{2}{c}{WavLM} \\
Group & EER & TAR@1\% & EER & TAR@1\% & EER & TAR@1\% \\
\midrule
African (men)               & 0.4553 & 0.1804 & 0.6794 & 0.0796 & 0.1994 & 0.3259 \\
African (women)             & 0.5199 & 0.1371 & 0.6831 & 0.0373 & 0.2305 & 0.2030 \\
Asian, East (men)           & 0.4537 & 0.1885 & 0.6671 & 0.0413 & 0.2332 & 0.2181 \\
Asian, East (women)         & 0.5203 & 0.1413 & 0.6599 & 0.0817 & 0.2768 & 0.1424 \\
Asian, South (men)          & 0.4827 & 0.1617 & 0.6497 & 0.0593 & 0.1907 & 0.2857 \\
Asian, South (women)        & 0.5443 & 0.1570 & 0.6276 & 0.0595 & 0.2177 & 0.2380 \\
Caucasian, American (men)   & 0.5662 & 0.0969 & 0.7244 & 0.0281 & 0.2515 & 0.1546 \\
Caucasian, American (women) & 0.6036 & 0.0844 & 0.7364 & 0.0230 & 0.3070 & 0.0537 \\
Caucasian, European (men)   & 0.5897 & 0.0955 & 0.6714 & 0.0477 & 0.2115 & 0.2804 \\
Caucasian, European (women) & 0.6575 & 0.0409 & 0.7412 & 0.0102 & 0.2788 & 0.1317 \\
\bottomrule
\end{tabular}
\end{table}

Averaged across the three verifiers, women have a higher mean EER than
men (0.5070 vs.\ 0.4684) and the two Caucasian descent groups are the
highest-EER (worst-imperceptibility) group -- the same gender and
descent direction found throughout this paper, but here it means
something different: since AntiFake's design goal is a \emph{low} EER
(the cloak remaining verifiable, matching POP's and attack-vc's
near-ceiling regime), a higher EER indicates the cloak is failing its
own imperceptibility objective \emph{more badly} for women and
Caucasian speakers, not that they are better protected. We return to
this same demographic ordering in
\S\ref{sec:protective-efficacy-demographics-antifake} and show it
predicts protective efficacy in the same direction, consistent with a
single underlying mechanism: however far the decoy attack displaces a
given speaker's embedding, that displacement shows up as \emph{both} a
larger imperceptibility failure (measured here) \emph{and} a larger
protective effect downstream (measured in
\S\ref{sec:protective-efficacy-antifake}).

\section{Results: ProtectYourAudio}
\label{sec:results-protectyouraudio}

Table~\ref{tab:results-protectyouraudio} reports the same
verification-leakage and signal-quality metrics as
Tables~\ref{tab:results}/\ref{tab:results-attackvc}/\ref{tab:results-antifake},
this time for clean\_bonafide vs.\ ProtectYourAudio-cloaked (4,850 of
the 4,867-recording corpus succeeded, all on this evaluation machine
plus a second machine cloaking in parallel against the same shared
output directory -- see \S\ref{sec:cloaking}'s engineering-obstacles
discussion for this project's general approach to partial-corpus
runs; 48,500 trials per verifier). Like POP and attack-vc,
ProtectYourAudio's greedy frequency-band-masking search
(\S\ref{sec:related}) is explicitly budgeted against a
perceptual-similarity constraint rather than optimized to push an
embedding as far as possible, and the results bear that out: EER is
below 2.1\% and TAR@1\%FAR above 96.2\% for every verifier -- an
imperceptibility result essentially indistinguishable from attack-vc's
(Table~\ref{tab:results-attackvc}) and far stronger than AntiFake's
(Table~\ref{tab:results-antifake}).

\begin{table}[H]
\centering
\caption{clean\_bonafide vs.\ ProtectYourAudio-cloaked: verification
leakage and signal-quality metrics, pooled across all three source
corpora.}
\label{tab:results-protectyouraudio}
\begin{tabular}{lccccccc}
\toprule
SV & EER & TAR@1\% & minDCF & $d'$ & STOI & PESQ & SI-SDR (dB) \\
\midrule
SB-ECAPA     & 0.0010 & 1.0000 & 0.2054 & 11.295 & \multirow{3}{*}{0.9411} & \multirow{3}{*}{3.4764} & \multirow{3}{*}{$-$25.3909} \\
Resemblyzer  & 0.0007 & 1.0000 & 0.1351 & 7.2448 & & & \\
WavLM        & 0.0205 & 0.9619 & 0.4847 & 2.883  & & & \\
\bottomrule
\end{tabular}
\end{table}

ProtectYourAudio's EER/TAR@1\%FAR figures track attack-vc's almost
exactly, verifier by verifier (SB-ECAPA 0.0010 vs.\ 0.0010,
Resemblyzer 0.0007 vs.\ 0.0008, WavLM 0.0205 vs.\ 0.0208) -- two
mechanistically unrelated cloaking methods (discrete frequency-band
masking search vs.\ a continuous embedding-attack perturbation)
converging on essentially the same imperceptibility result. The
signal-quality metrics, however, do \emph{not} track attack-vc: STOI
(0.9411) sits close to POP's 0.9709, not attack-vc's degraded 0.5472,
while SI-SDR ($-$25.4~dB) is negative like attack-vc's ($-$34.0~dB,
\S\ref{sec:results-attackvc}) rather than POP's strongly positive
$+$16.3~dB, and PESQ (3.48) is closer to attack-vc's 3.58 than POP's
2.92. ProtectYourAudio is, in other words, the one technique in this
paper that does not fall cleanly into either the ``POP-like'' or
``attack-vc-like'' quality profile -- high perceptual intelligibility
(STOI) paired with a collapsed sample-alignment score (SI-SDR) is a
combination neither other technique produces. We read this the same
way as attack-vc's own STOI/SI-SDR divergence
(\S\ref{sec:results-attackvc}): ProtectYourAudio's masking search
operates in the frequency domain and its output undergoes a
sample-rate conversion (24kHz native vs.\ clean\_bonafide's 16kHz)
before scoring, either or both of which can perturb sample-level
phase/timing alignment without perceptibly degrading the signal --
SI-SDR is acutely sensitive to exactly that kind of mismatch, while
STOI and PESQ are comparatively tolerant of it
(\S\ref{sec:quality-metrics}). As with attack-vc, we have not
confirmed this by direct waveform inspection and flag it as the
leading explanation rather than a verified fact.

\subsection{Ablation: verification leakage by descent $\times$ gender}
\label{sec:ablation-protectyouraudio}
Table~\ref{tab:ablation-protectyouraudio} reports the same per-group
breakdown as Tables~\ref{tab:ablation}/\ref{tab:ablation-attackvc}/\ref{tab:ablation-antifake},
restricted to the FakeAVCeleb subset, for ProtectYourAudio.

\begin{table}[H]
\centering
\caption{clean\_bonafide vs.\ ProtectYourAudio-cloaked: EER /
TAR@1\%FAR by descent $\times$ gender (FakeAVCeleb subset).}
\label{tab:ablation-protectyouraudio}
\begin{tabular}{lcccccc}
\toprule
& \multicolumn{2}{c}{SB-ECAPA} & \multicolumn{2}{c}{Resemblyzer} & \multicolumn{2}{c}{WavLM} \\
Group & EER & TAR@1\% & EER & TAR@1\% & EER & TAR@1\% \\
\midrule
African (men)               & 0.0006 & 1.0000 & 0.0008 & 1.0000 & 0.0206 & 0.9610 \\
African (women)              & 0.0010 & 1.0000 & 0.0020 & 1.0000 & 0.0096 & 0.9898 \\
Asian, East (men)            & 0.0013 & 1.0000 & 0.0003 & 1.0000 & 0.0145 & 0.9773 \\
Asian, East (women)          & 0.0005 & 1.0000 & 0.0004 & 1.0000 & 0.0120 & 0.9849 \\
Asian, South (men)           & 0.0005 & 1.0000 & 0.0005 & 1.0000 & 0.0117 & 0.9668 \\
Asian, South (women)         & 0.0004 & 1.0000 & 0.0003 & 1.0000 & 0.0249 & 0.9291 \\
Caucasian, American (men)    & 0.0006 & 1.0000 & 0.0009 & 0.8162 & 0.0076 & 0.9975 \\
Caucasian, American (women)  & 0.0008 & 1.0000 & 0.0005 & 1.0000 & 0.0072 & 0.9974 \\
Caucasian, European (men)    & 0.0005 & 1.0000 & 0.0000 & 0.7563 & 0.0161 & 0.9825 \\
Caucasian, European (women)  & 0.0005 & 1.0000 & 0.0004 & 1.0000 & 0.0127 & 0.9821 \\
\bottomrule
\end{tabular}
\end{table}

The per-group pattern is the most uniform of any technique reported in
this paper: SB-ECAPA's TAR@1\%FAR is exactly 100\% in all 10 groups,
and WavLM's is never below 92.9\%. Resemblyzer is the one exception --
Caucasian(American)-men (81.6\%) and Caucasian(European)-men (75.6\%)
are the only two group$\times$verifier cells below 99\% anywhere in
this table, both concentrated in the same verifier and the same
gender, unlike attack-vc's more evenly scattered shortfall
(\S\ref{sec:ablation-attackvc}) or POP's two isolated outlier groups
(\S\ref{sec:ablation}). We do not have a confirmed explanation for why
Resemblyzer specifically -- the encoder family ProtectYourAudio's own
vendored pipeline uses as a surrogate (\S\ref{sec:verifiers}) -- shows
this concentrated, gender-specific shortfall while SB-ECAPA and WavLM
do not; the surrogate-encoder relationship makes it plausible that
this is where ProtectYourAudio's masking search is most directly
optimizing against, but we flag this as an open question rather than a
confirmed mechanism. As with every other technique in this paper, even
this most-scattered group is far closer to the near-ceiling
imperceptibility regime than to the uncloaked FakeAVCeleb calibration
baseline (TAR@1\%FAR 11--15\%, Table~\ref{tab:calibration}), so none of
this scatter constitutes a meaningful protective effect on its own.

\section{Protective-Efficacy Evaluation: Resynthesis from POP-Cloaked Audio}
\label{sec:protective-efficacy}

Every result reported so far in this paper (Tables~\ref{tab:results},
\ref{tab:ablation}, \ref{tab:results-attackvc},
\ref{tab:ablation-attackvc}) compares a cloaked recording
\emph{directly} against the genuine original. This measures
\emph{imperceptibility}, not protection, and the two must not be
conflated: if a cloak is well designed, the cloaked recording
\emph{should} still verify as the same speaker (the ``same'' outcome
above is success for imperceptibility). Protective efficacy is a
different, downstream question entirely -- does the cloak stop an
attacker who feeds the cloaked recording into a cloning system from
extracting the speaker's true identity? -- and it can only be
answered by taking the further step of actually cloning
\emph{from the cloaked audio} and comparing \emph{that} output against
the genuine original. For that comparison, the interpretation is the
opposite of the imperceptibility tables above: if the resulting clone
still verifies as the same speaker (same as the clean-audio ceiling of
Table~\ref{tab:synth-ceiling}), the cloak \emph{failed} to protect --
the attacker still obtained a usable clone of the true voice. Only if
the clone verifies as \emph{different} from the genuine speaker
(higher EER / lower TAR@FAR than the clean-audio ceiling) has the
cloak actually protected the speaker.

We ran this evaluation for POP: for every one of the 4,864
POP-cloaked recordings, we cloned it with all three systems from
Table~\ref{tab:synth} (SV2TTS, Seed-VC, and F5-TTS) used to build the
clean-audio ceiling of Table~\ref{tab:synth-ceiling}, following the
identical harmonized trial-construction protocol of
\S\ref{sec:threat}, then scored the resulting clone against the
genuine clean\_bonafide original with all three verifiers.
Table~\ref{tab:protective-efficacy} reports both the clean-audio
ceiling (repeated from Table~\ref{tab:synth-ceiling}) and the
POP-cloaked-source result side by side, so the shift between them can
be read directly as POP's protective effect (or lack thereof) against
each system.

\begin{table}[H]
\centering
\caption{Cloning from POP-cloaked source audio vs.\ the clean-audio
ceiling (Table~\ref{tab:synth-ceiling}): does the resulting clone
still verify as the genuine speaker?}
\label{tab:protective-efficacy}
\begin{tabular}{llcccc}
\toprule
System & SV & \multicolumn{2}{c}{Clean-audio ceiling} & \multicolumn{2}{c}{POP-cloaked source} \\
       &    & EER & TAR@1\% & EER & TAR@1\% \\
\midrule
SV2TTS   & SB-ECAPA    & 0.1415 & 0.6007 & 0.2065 & 0.3657 \\
SV2TTS   & Resemblyzer & 0.0547 & 0.8453 & 0.1018 & 0.6763 \\
SV2TTS   & WavLM       & 0.1667 & 0.3469 & 0.1918 & 0.2251 \\
\addlinespace
Seed-VC  & SB-ECAPA    & 0.0311 & 0.9316 & 0.0359 & 0.9167 \\
Seed-VC  & Resemblyzer & 0.0499 & 0.8436 & 0.0531 & 0.8227 \\
Seed-VC  & WavLM       & 0.1601 & 0.2509 & 0.1674 & 0.2461 \\
\addlinespace
F5-TTS   & SB-ECAPA    & 0.0166 & 0.9768 & 0.0256 & 0.9607 \\
F5-TTS   & Resemblyzer & 0.0223 & 0.9632 & 0.0251 & 0.9568 \\
F5-TTS   & WavLM       & 0.0924 & 0.6658 & 0.1092 & 0.6028 \\
\bottomrule
\end{tabular}
\end{table}

The three systems tell markedly different stories. Against
\textbf{SV2TTS}, POP shows a real, if moderate, protective effect: EER
rises from the clean-audio ceiling by 46\% (SB-ECAPA), 86\%
(Resemblyzer), and 15\% (WavLM) in relative terms, while TAR@1\%FAR
falls by 20--39\% relative across all three verifiers -- cloning from
POP-cloaked source audio measurably degrades the resulting clone's
resemblance to the true speaker, exactly the ``different from the
ceiling'' outcome that constitutes evidence of protection. Against
\textbf{Seed-VC}, by contrast, the shift is marginal: EER rises by
only 5--15\% relative and TAR@1\%FAR falls by only 2--3\% relative
across all three verifiers -- well within the kind of run-to-run noise
we would expect even absent any cloaking effect at all, and far short
of the shift seen with SV2TTS. \textbf{F5-TTS} again falls in between,
echoing the same pattern seen for attack-vc
(\S\ref{sec:protective-efficacy-attackvc}): EER rises by a moderate
13--54\% relative (SB-ECAPA 54\%, WavLM 18\%, Resemblyzer 13\%), but
TAR@1\%FAR barely moves (down just 0.7--9.5\% relative) -- the
distributions separate somewhat on average without meaningfully
changing how many clones would be accepted at a fixed, practical
false-accept-rate threshold. Read together with
Table~\ref{tab:results}/\ref{tab:ablation} (which confirm POP-cloaked
audio remains, on its own, essentially indistinguishable from genuine
audio), this indicates POP's protective effect is
\emph{system-dependent}: it meaningfully disrupts a zero-shot
TTS attacker (SV2TTS), has a real but more modest effect on a newer
zero-shot TTS attacker (F5-TTS) that barely moves the security-relevant
TAR@1\%FAR operating point, and has little to no effect against a
voice-conversion attacker (Seed-VC). We initially hypothesized that
POP's surrogate during training is VITS, itself a TTS architecture
(\S\ref{sec:related}), so the crafted perturbation might transfer
better to TTS attackers (SV2TTS, and to a lesser extent F5-TTS) than to
a voice-conversion attacker (Seed-VC) that does not share VITS's
sequence-to-sequence, text-conditioned synthesis mechanism. We tested
this directly in \S\ref{sec:ablation-seedvc-mechanism} and found it
does not hold: neither cloak's surrogate/target shares an encoder
architecture with Seed-VC's or SV2TTS's own conditioning mechanism at
all, and an embedding-level probe shows the perturbation reaches
Seed-VC's own identity representation \emph{at least as much as} it
reaches SV2TTS's -- Seed-VC's near-immunity is therefore better
explained by its decoder tolerating a perturbed embedding than by the
perturbation failing to transfer to that embedding in the first place.
This still does not explain why F5-TTS, also nominally a TTS system,
is affected noticeably less than SV2TTS; \S\ref{sec:ablation-seedvc-mechanism}
discusses this as an open refinement rather than a fully confirmed
mechanism.

\subsection{Does descent or gender predict protective efficacy?}
\label{sec:protective-efficacy-demographics}

Averaging each FakeAVCeleb group's EER across the three verifiers (as
in \S\ref{sec:ablation}) reveals a clear, consistent demographic
pattern in the resynthesis-based protective-efficacy result: across
all three systems, cloning from women's recordings yields a
consistently higher EER (harder to re-identify) than cloning from
men's, and the two Caucasian descent groups are consistently the
\emph{lowest}-EER (least-protected) group. Table~\ref{tab:protective-efficacy-demographics}
reports this breakdown alongside the equivalent breakdown for the
clean-audio ceiling (no cloaking at all), to test whether this pattern
is something POP introduces or merely inherits.

\begin{table}[H]
\centering
\caption{Mean EER by descent/gender group, clean-audio ceiling vs.\
POP-cloaked source (averaged across the three verifiers).}
\label{tab:protective-efficacy-demographics}
\begin{tabular}{lcccccc}
\toprule
& \multicolumn{2}{c}{SV2TTS} & \multicolumn{2}{c}{Seed-VC} & \multicolumn{2}{c}{F5-TTS} \\
Group & Ceiling & Cloaked & Ceiling & Cloaked & Ceiling & Cloaked \\
\midrule
Men                  & 0.1497 & 0.1510 & 0.0702 & 0.0759 & 0.0343 & 0.0499 \\
Women                & 0.1752 & 0.2002 & 0.0904 & 0.0925 & 0.0447 & 0.0552 \\
\addlinespace
African              & 0.1728 & 0.1790 & 0.0933 & 0.0920 & 0.0536 & 0.0649 \\
Asian (East)         & 0.1543 & 0.1694 & 0.0792 & 0.0796 & 0.0595 & 0.0670 \\
Asian (South)        & 0.2115 & 0.2137 & 0.1021 & 0.1088 & 0.0340 & 0.0608 \\
Caucasian (American)  & 0.1360 & 0.1547 & 0.0642 & 0.0706 & 0.0185 & 0.0250 \\
Caucasian (European)  & 0.1377 & 0.1611 & 0.0629 & 0.0701 & 0.0321 & 0.0451 \\
\bottomrule
\end{tabular}
\end{table}

The answer is that this disparity \emph{predates} POP entirely: the
gender and descent ordering in each cloaked column is essentially the
same ordering already present in that system's own clean-audio
ceiling column. Women are already harder than men for all three
systems to convincingly reproduce even from \emph{clean, uncloaked}
reference audio, and POP's cloaking shifts the absolute EER upward for
every group (its genuine protective effect, \S\ref{sec:protective-efficacy}),
without meaningfully changing \emph{which} groups that protection
favors. Concretely, the men/women EER gap is 2.55 points in SV2TTS's
clean ceiling and 4.92 points after POP-cloaking -- POP widens the
pre-existing gap somewhat rather than closing it, but it did not
create it. The same is true for descent: SV2TTS and Seed-VC both rank
Asian (South) as the highest-EER (most-protected) descent group, but
F5-TTS instead ranks Asian (East) highest, with Asian (South) actually
tied for second-lowest with Caucasian (European) in F5-TTS's own
clean-audio ceiling (0.0340 each) -- this system-specific reordering
already exists before any cloaking is applied, exactly mirroring what
we find for attack-vc
(\S\ref{sec:protective-efficacy-demographics-attackvc}), so it
reflects a property of the synthesizer/verifier pipelines rather than
anything specific to POP. We read the overall pattern as an important
caveat on the demographic-fairness framing of \S\ref{sec:dataset}: a
demographically balanced \emph{dataset} does not, by itself, guarantee
that a cloaking method's protective benefit is distributed evenly
across that dataset's demographic groups, since the downstream
cloning systems being defended against are themselves not
demographically neutral.

\section{Protective-Efficacy Evaluation: Resynthesis from attack-vc-Cloaked Audio}
\label{sec:protective-efficacy-attackvc}

We repeated the resynthesis-based protective-efficacy evaluation of
\S\ref{sec:protective-efficacy} for attack-vc, this time against all
three downstream systems from Table~\ref{tab:synth} (SV2TTS, Seed-VC,
\emph{and} F5-TTS -- unlike the POP evaluation, all three of
attack-vc's runs were complete at time of writing). For every one of
the 4,863 attack-vc-cloaked recordings, we cloned it with each system
using the identical harmonized protocol of \S\ref{sec:threat}, then
scored the resulting clone against the genuine clean\_bonafide
original with all three verifiers. As in \S\ref{sec:protective-efficacy},
the interpretation is: if the clone verifies as the same speaker (same
as the clean-audio ceiling of Table~\ref{tab:synth-ceiling}), attack-vc
\emph{failed} to protect; only a shift toward \emph{different} (higher
EER / lower TAR@FAR than the ceiling) constitutes evidence of
protection.

\begin{table}[H]
\centering
\caption{Cloning from attack-vc-cloaked source audio vs.\ the
clean-audio ceiling (Table~\ref{tab:synth-ceiling}): does the resulting
clone still verify as the genuine speaker?}
\label{tab:protective-efficacy-attackvc}
\begin{tabular}{llcccc}
\toprule
System & SV & \multicolumn{2}{c}{Clean-audio ceiling} & \multicolumn{2}{c}{attack-vc-cloaked source} \\
       &    & EER & TAR@1\% & EER & TAR@1\% \\
\midrule
SV2TTS   & SB-ECAPA    & 0.1415 & 0.6007 & 0.2112 & 0.3530 \\
SV2TTS   & Resemblyzer & 0.0547 & 0.8453 & 0.1064 & 0.6615 \\
SV2TTS   & WavLM       & 0.1667 & 0.3469 & 0.1914 & 0.2150 \\
\addlinespace
Seed-VC  & SB-ECAPA    & 0.0311 & 0.9316 & 0.0329 & 0.9241 \\
Seed-VC  & Resemblyzer & 0.0499 & 0.8436 & 0.0508 & 0.8327 \\
Seed-VC  & WavLM       & 0.1601 & 0.2509 & 0.1661 & 0.2499 \\
\addlinespace
F5-TTS   & SB-ECAPA    & 0.0166 & 0.9768 & 0.0212 & 0.9696 \\
F5-TTS   & Resemblyzer & 0.0223 & 0.9632 & 0.0269 & 0.9568 \\
F5-TTS   & WavLM       & 0.0924 & 0.6658 & 0.1014 & 0.6440 \\
\bottomrule
\end{tabular}
\end{table}

The three-system picture closely mirrors POP's
(\S\ref{sec:protective-efficacy}), both in which system is most
affected and by roughly how much. Against \textbf{SV2TTS}, attack-vc
shows a real, moderate protective effect essentially the same size as
POP's: EER rises by 49\% (SB-ECAPA), 95\% (Resemblyzer), and 15\%
(WavLM) in relative terms, and TAR@1\%FAR falls by 22--41\% relative
across all three verifiers -- despite attack-vc and POP being
mechanistically unrelated (an inference-time embedding attack against
a voice-conversion encoder vs.\ a training-time unlearnable-example
perturbation against a TTS surrogate), they converge on almost
identical SV2TTS outcomes. Against \textbf{Seed-VC}, the shift is
negligible (EER up 2--6\% relative, TAR@1\%FAR down under 1.3\%
relative) -- the same near-zero effect POP showed. \textbf{F5-TTS}
falls in between and reveals a nuance the EER numbers alone would
miss: EER rises by a moderate 10--28\% relative (comparable in
direction to SV2TTS, just smaller), but TAR@1\%FAR -- the
security-relevant, fixed-operating-point metric -- barely moves at all
(down 0.7--3.3\% relative, closer to Seed-VC's negligible shift than to
SV2TTS's 22--41\% drop). This is exactly why \S\ref{sec:metrics}
reports TAR@FAR alongside EER rather than EER alone: F5-TTS's higher
EER reflects the two score distributions moving slightly farther apart
on average, but not enough to change how many clones would actually be
accepted at a practical, fixed false-accept-rate threshold. Taken
together, both cloaking techniques evaluated so far transfer well to
a zero-shot TTS attacker (SV2TTS) and poorly to a voice-conversion
attacker (Seed-VC), with F5-TTS -- a newer, stronger zero-shot TTS
system -- sitting closer to the voice-conversion outcome than the
older SV2TTS outcome despite being the same broad system type,
suggesting the transfer pattern tracks something more specific than
just ``TTS vs.\ VC.''

\subsection{Does descent or gender predict protective efficacy?}
\label{sec:protective-efficacy-demographics-attackvc}

Repeating the demographic breakdown of
\S\ref{sec:protective-efficacy-demographics} for attack-vc's three
systems (Table~\ref{tab:protective-efficacy-demographics-attackvc})
reproduces the same qualitative pattern found for POP: cloning from
women's recordings yields a consistently higher mean EER than men's,
across all three systems and both the clean-audio ceiling and the
attack-vc-cloaked-source conditions; and the two Caucasian descent
groups are consistently the \emph{lowest}-EER (least-protected) group
in every system and condition. Both patterns already exist in each
system's own clean-audio ceiling, confirming again that this
disparity is inherited from the underlying synthesizer/verifier
pipeline, not introduced or meaningfully worsened by the cloak itself.

\begin{table}[H]
\centering
\caption{Mean EER by descent/gender group, clean-audio ceiling vs.\
attack-vc-cloaked source (averaged across the three verifiers).}
\label{tab:protective-efficacy-demographics-attackvc}
\begin{tabular}{lcccccc}
\toprule
& \multicolumn{2}{c}{SV2TTS} & \multicolumn{2}{c}{Seed-VC} & \multicolumn{2}{c}{F5-TTS} \\
Group & Ceiling & Cloaked & Ceiling & Cloaked & Ceiling & Cloaked \\
\midrule
Men                  & 0.1497 & 0.1568 & 0.0702 & 0.0757 & 0.0343 & 0.0396 \\
Women                & 0.1752 & 0.2051 & 0.0904 & 0.0890 & 0.0447 & 0.0473 \\
\addlinespace
African              & 0.1728 & 0.1824 & 0.0933 & 0.0930 & 0.0536 & 0.0577 \\
Asian (East)         & 0.1543 & 0.1713 & 0.0792 & 0.0805 & 0.0595 & 0.0614 \\
Asian (South)        & 0.2115 & 0.2216 & 0.1021 & 0.1114 & 0.0340 & 0.0433 \\
Caucasian (American)  & 0.1360 & 0.1577 & 0.0642 & 0.0654 & 0.0185 & 0.0226 \\
Caucasian (European)  & 0.1377 & 0.1717 & 0.0629 & 0.0614 & 0.0340 & 0.0324 \\
\bottomrule
\end{tabular}
\end{table}

One genuine difference from POP is worth flagging rather than
smoothing over: for SV2TTS and Seed-VC, \emph{Asian (South)} is the
highest-EER (most-protected) descent group, but for \emph{F5-TTS} it
is \emph{Asian (East)} instead, with Asian (South) dropping to the
middle of the ranking -- tied for second-lowest with Caucasian
(European) in the clean-audio ceiling (0.0340 each, both just above
Caucasian (American)'s 0.0185), then to third-lowest of five once
cloaked (0.0433, closer to the African/Asian (East) group than to
either Caucasian group). This ordering already differs between
systems in the clean-audio ceiling itself (i.e.\ before any cloaking),
so it reflects a property of the synthesizer/verifier pipelines rather
than anything attack-vc's perturbation introduces -- but it does mean
the ``which descent group benefits most'' answer is system-dependent,
not a fixed property of the dataset, and should not be generalized
from a single downstream system.

\section{Protective-Efficacy Evaluation: Resynthesis from AntiFake-Cloaked Audio}
\label{sec:protective-efficacy-antifake}

We repeated the resynthesis-based protective-efficacy evaluation of
\S\ref{sec:protective-efficacy}/\S\ref{sec:protective-efficacy-attackvc}
for AntiFake, against all three downstream systems from
Table~\ref{tab:synth}. For each of the 4,844 AntiFake-cloaked
recordings, we cloned it with each system using the identical
harmonized protocol of \S\ref{sec:threat}, then scored the resulting
clone against the genuine clean\_bonafide original with all three
verifiers. The interpretation is unchanged from the other two
techniques: a clone that verifies as \emph{different} from the genuine
speaker (higher EER / lower TAR@FAR than the clean-audio ceiling of
Table~\ref{tab:synth-ceiling}) is evidence of protection.

\begin{table}[H]
\centering
\caption{Cloning from AntiFake-cloaked source audio vs.\ the
clean-audio ceiling (Table~\ref{tab:synth-ceiling}): does the resulting
clone still verify as the genuine speaker?}
\label{tab:protective-efficacy-antifake}
\begin{tabular}{llcccc}
\toprule
System & SV & \multicolumn{2}{c}{Clean-audio ceiling} & \multicolumn{2}{c}{AntiFake-cloaked source} \\
       &    & EER & TAR@1\% & EER & TAR@1\% \\
\midrule
SV2TTS   & SB-ECAPA    & 0.1415 & 0.6007 & 0.5578 & 0.0117 \\
SV2TTS   & Resemblyzer & 0.0547 & 0.8453 & 0.4999 & 0.0246 \\
SV2TTS   & WavLM       & 0.1667 & 0.3469 & 0.4659 & 0.0077 \\
\addlinespace
Seed-VC  & SB-ECAPA    & 0.0311 & 0.9316 & 0.4931 & 0.0833 \\
Seed-VC  & Resemblyzer & 0.0499 & 0.8436 & 0.3249 & 0.2051 \\
Seed-VC  & WavLM       & 0.1601 & 0.2509 & 0.3454 & 0.0293 \\
\addlinespace
F5-TTS   & SB-ECAPA    & 0.0166 & 0.9768 & 0.3817 & 0.1661 \\
F5-TTS   & Resemblyzer & 0.0223 & 0.9632 & 0.2138 & 0.3750 \\
F5-TTS   & WavLM       & 0.0924 & 0.6658 & 0.3013 & 0.0992 \\
\bottomrule
\end{tabular}
\end{table}

AntiFake breaks the pattern both POP and attack-vc established.
Against \textbf{SV2TTS}, the protective effect is far stronger than
either other technique: EER rises 179--814\% relative over the
clean-audio ceiling (SB-ECAPA 294\%, Resemblyzer 814\%, WavLM 179\%),
and TAR@1\%FAR collapses by 97--98\% relative, down to 0.8--2.5\% in
absolute terms -- compare to POP's and attack-vc's 20--41\% relative
TAR@1\%FAR drops against the same system. Against \textbf{Seed-VC} --
the system POP and attack-vc both had \emph{negligible} effect against
(EER up only 2--15\% relative, \S\ref{sec:protective-efficacy}/\S\ref{sec:protective-efficacy-attackvc})
-- AntiFake instead shows a large effect here too: EER rises
116--1{,}485\% relative (SB-ECAPA's 0.0311 ceiling in particular means
even a modest absolute EER of 0.4931 is a nearly 15$\times$ relative
increase), and TAR@1\%FAR falls 76--91\% relative. Against
\textbf{F5-TTS} -- the system where both POP and attack-vc showed only
a moderate effect with a security-relevant TAR@1\%FAR that barely moved
(\S\ref{sec:protective-efficacy}/\S\ref{sec:protective-efficacy-attackvc})
-- AntiFake's effect is the largest of all nine (system, verifier)
cells reported across all three techniques: EER rises 226--2{,}199\%
relative (SB-ECAPA alone: from 0.0166 to 0.3817, a 23$\times$ increase),
and TAR@1\%FAR collapses 61--85\% relative, down to 9.9--37.5\% in
absolute terms. AntiFake is, uniquely among the three techniques
evaluated in this paper, effective against \emph{all three} downstream
systems -- two zero-shot TTS attackers and a voice-conversion attacker
-- rather than only one or two. Read together with
\S\ref{sec:results-antifake}, this is the same conclusion from a
different angle: AntiFake's decoy-target attack is not calibrated to a
fixed imperceptibility budget the way POP and attack-vc are, so it
pushes harder in embedding space, and that same aggressiveness that
makes it fail as an imperceptible cloak is exactly what makes its
protective effect transfer uniformly across systems where the other two
techniques' more conservative, budget-constrained perturbations do not.

\subsection{Does descent or gender predict protective efficacy?}
\label{sec:protective-efficacy-demographics-antifake}

Table~\ref{tab:protective-efficacy-demographics-antifake} reports the
same mean-EER-by-group breakdown as
Tables~\ref{tab:protective-efficacy-demographics}/\ref{tab:protective-efficacy-demographics-attackvc},
averaged across the three verifiers, for all three downstream systems.

\begin{table}[H]
\centering
\caption{Mean EER by descent/gender group, clean-audio ceiling vs.\
AntiFake-cloaked source (averaged across the three verifiers).}
\label{tab:protective-efficacy-demographics-antifake}
\begin{tabular}{lcccccc}
\toprule
& \multicolumn{2}{c}{SV2TTS} & \multicolumn{2}{c}{Seed-VC} & \multicolumn{2}{c}{F5-TTS} \\
Group & Ceiling & Cloaked & Ceiling & Cloaked & Ceiling & Cloaked \\
\midrule
Men                  & 0.1497 & 0.5159 & 0.0702 & 0.4128 & 0.0343 & 0.3282 \\
Women                & 0.1752 & 0.5538 & 0.0904 & 0.4345 & 0.0447 & 0.3295 \\
\addlinespace
African              & 0.1728 & 0.5255 & 0.0933 & 0.4146 & 0.0536 & 0.3609 \\
Asian (East)         & 0.1543 & 0.5179 & 0.0792 & 0.3966 & 0.0595 & 0.2870 \\
Asian (South)        & 0.2115 & 0.5260 & 0.1021 & 0.4048 & 0.0340 & 0.4008 \\
Caucasian (American)  & 0.1360 & 0.5542 & 0.0642 & 0.4662 & 0.0185 & 0.2626 \\
Caucasian (European)  & 0.1377 & 0.5506 & 0.0629 & 0.4361 & 0.0321 & 0.3329 \\
\bottomrule
\end{tabular}
\end{table}

For SV2TTS and Seed-VC, the gender and descent pattern matches the
direction found for POP and attack-vc -- women have a higher mean EER
than men (0.5538 vs.\ 0.5159 for SV2TTS, 0.4345 vs.\ 0.4128 for
Seed-VC), and the two Caucasian descent groups are the
\emph{highest}-EER group under AntiFake (the opposite of POP/attack-vc,
where Caucasian groups were consistently \emph{lowest}-EER, i.e.\ least
protected). This inversion is not a contradiction: it follows directly
from \S\ref{sec:results-antifake}, where Caucasian speakers already
showed the worst imperceptibility (highest EER, i.e.\ the cloak already
pushed their embeddings furthest from genuine) -- the same displacement
that hurts imperceptibility here helps protective efficacy, so the
group ordering \emph{by imperceptibility failure} carries over directly
to the group ordering \emph{by protection}. This is a genuinely
different mechanism from POP and attack-vc, where the demographic
pattern in protective efficacy was inherited from the downstream
synthesizer/verifier pipeline's own pre-existing bias
(\S\ref{sec:protective-efficacy-demographics}), not from the cloak's
own imperceptibility profile -- for AntiFake, the two are coupled by
construction, since the decoy-target attack optimizes the same
embedding-displacement objective that both metrics measure.

\textbf{F5-TTS breaks this pattern}, and we report the divergence
plainly rather than force a false consistency. The gender gap nearly
vanishes -- men 0.3282 vs.\ women 0.3295, a difference an order of
magnitude smaller than SV2TTS's or Seed-VC's own gender gaps under
AntiFake, and well within the kind of noise we would not read as
meaningful elsewhere in this paper. The descent ranking partially
agrees (Asian (South) is again the highest-EER group, 0.4008,
consistent with SV2TTS and Seed-VC) but partially inverts: Caucasian
(American) is the \emph{lowest}-EER group under F5-TTS (0.2626), the
opposite of its highest-EER standing under SV2TTS and Seed-VC.
Caucasian (American) also has, by a wide margin, the lowest clean-audio
ceiling of any group under F5-TTS (0.0185, versus 0.0629--0.0595 for
every other group) -- the same floor-effect caveat already raised in
\S\ref{sec:protective-efficacy-demographics}/\S\ref{sec:protective-efficacy-demographics-attackvc}
for F5-TTS specifically applies here: a group whose clean-audio ceiling
is already near-perfect leaves less room for the embedding-displacement
mechanism described above to manifest as a proportionally larger EER
increase, independent of how strongly AntiFake actually displaces that
group's embedding. We do not have a confirmed explanation for why this
floor effect dominates for F5-TTS specifically but not for SV2TTS or
Seed-VC; we report it as an open question rather than force it into the
same displacement-coupling account that fits the other two systems.

\section{Protective-Efficacy Evaluation: Resynthesis from ProtectYourAudio-Cloaked Audio}
\label{sec:protective-efficacy-protectyouraudio}

We repeated the resynthesis-based protective-efficacy evaluation of
\S\ref{sec:protective-efficacy}/\S\ref{sec:protective-efficacy-attackvc}/\S\ref{sec:protective-efficacy-antifake}
for ProtectYourAudio, against all three downstream systems from
Table~\ref{tab:synth}. For each of the 4,850 ProtectYourAudio-cloaked
recordings (4,793 for F5-TTS specifically, reflecting that system's own
small per-clip failure rate already documented elsewhere in this
paper), we cloned it with each system using the identical harmonized
protocol of \S\ref{sec:threat}, then scored the resulting clone
against the genuine clean\_bonafide original with all three verifiers.
The interpretation is unchanged: a clone that verifies as
\emph{different} from the genuine speaker (higher EER / lower
TAR@FAR than the clean-audio ceiling of Table~\ref{tab:synth-ceiling})
is evidence of protection.

\begin{table}[H]
\centering
\caption{Cloning from ProtectYourAudio-cloaked source audio vs.\ the
clean-audio ceiling (Table~\ref{tab:synth-ceiling}): does the resulting
clone still verify as the genuine speaker?}
\label{tab:protective-efficacy-protectyouraudio}
\begin{tabular}{llcccc}
\toprule
System & SV & \multicolumn{2}{c}{Clean-audio ceiling} & \multicolumn{2}{c}{ProtectYourAudio-cloaked source} \\
       &    & EER & TAR@1\% & EER & TAR@1\% \\
\midrule
SV2TTS   & SB-ECAPA    & 0.1415 & 0.6007 & 0.2399 & 0.3136 \\
SV2TTS   & Resemblyzer & 0.0547 & 0.8453 & 0.1268 & 0.5899 \\
SV2TTS   & WavLM       & 0.1667 & 0.3469 & 0.2199 & 0.1835 \\
\addlinespace
Seed-VC  & SB-ECAPA    & 0.0311 & 0.9316 & 0.0621 & 0.8210 \\
Seed-VC  & Resemblyzer & 0.0499 & 0.8436 & 0.0739 & 0.7210 \\
Seed-VC  & WavLM       & 0.1601 & 0.2509 & 0.1844 & 0.1990 \\
\addlinespace
F5-TTS   & SB-ECAPA    & 0.0166 & 0.9768 & 0.0416 & 0.9255 \\
F5-TTS   & Resemblyzer & 0.0223 & 0.9632 & 0.0384 & 0.9230 \\
F5-TTS   & WavLM       & 0.0924 & 0.6658 & 0.1496 & 0.4744 \\
\bottomrule
\end{tabular}
\end{table}

ProtectYourAudio is, among the three techniques evaluated with a full
imperceptibility budget (i.e.\ excluding AntiFake, whose decoy attack
is not budget-constrained, \S\ref{sec:results-antifake}), the one with
the broadest cross-system protective effect. Against \textbf{SV2TTS},
the effect is comparable to POP's and attack-vc's own strongest
results: EER rises 32--132\% relative over the clean-audio ceiling
(SB-ECAPA 70\%, Resemblyzer 132\%, WavLM 32\%), and TAR@1\%FAR falls
30--48\% relative. Against \textbf{Seed-VC} -- the system both POP and
attack-vc had \emph{negligible} effect against (EER up only 2--15\%
relative, \S\ref{sec:protective-efficacy}/\S\ref{sec:protective-efficacy-attackvc})
-- ProtectYourAudio instead shows a real, if more modest, effect: EER
rises 15--100\% relative (SB-ECAPA 100\%, Resemblyzer 48\%, WavLM
15\%), and TAR@1\%FAR falls 12--21\% relative. Against \textbf{F5-TTS}
-- the system where POP and attack-vc showed only a moderate effect
with TAR@1\%FAR barely moving
(\S\ref{sec:protective-efficacy}/\S\ref{sec:protective-efficacy-attackvc})
-- ProtectYourAudio's effect is, surprisingly, the \emph{largest} of
its three systems in relative-EER terms: EER rises 62--151\% relative
(SB-ECAPA 151\%, Resemblyzer 72\%, WavLM 62\%), and TAR@1\%FAR falls
4--29\% relative, with WavLM's 29\% drop the most security-relevant
shift against F5-TTS reported anywhere in this paper outside AntiFake.
Read together, ProtectYourAudio is the one technique in this study
(again, excluding AntiFake) whose protective effect meaningfully
transfers beyond the SV2TTS system its own frequency-masking search
was tuned against -- a genuinely different outcome from POP and
attack-vc, whose Seed-VC/F5-TTS results were dominated by each
system's own near-ceiling baseline rather than by any real cloak
effect (\S\ref{sec:protective-efficacy}, \S\ref{sec:protective-efficacy-attackvc}).
We do not have a confirmed mechanism for why a discrete,
budget-constrained frequency-masking search generalizes across
downstream systems better than either a waveform-level ($\ell_\infty$,
POP) or mel-spectrogram-level (attack-vc) continuous perturbation
does, and flag it as an open question for follow-up work rather than
force a single account onto all three techniques.

\subsection{Does descent or gender predict protective efficacy?}
\label{sec:protective-efficacy-demographics-protectyouraudio}

Table~\ref{tab:protective-efficacy-demographics-protectyouraudio}
reports the same mean-EER-by-group breakdown as
Tables~\ref{tab:protective-efficacy-demographics}/\ref{tab:protective-efficacy-demographics-attackvc}/\ref{tab:protective-efficacy-demographics-antifake},
averaged across the three verifiers, for all three downstream systems.

\begin{table}[H]
\centering
\caption{Mean EER by descent/gender group, clean-audio ceiling vs.\
ProtectYourAudio-cloaked source (averaged across the three verifiers).}
\label{tab:protective-efficacy-demographics-protectyouraudio}
\begin{tabular}{lcccccc}
\toprule
& \multicolumn{2}{c}{SV2TTS} & \multicolumn{2}{c}{Seed-VC} & \multicolumn{2}{c}{F5-TTS} \\
Group & Ceiling & Cloaked & Ceiling & Cloaked & Ceiling & Cloaked \\
\midrule
Men                  & 0.1497 & 0.1769 & 0.0702 & 0.0923 & 0.0343 & 0.0664 \\
Women                & 0.1752 & 0.2398 & 0.0904 & 0.1272 & 0.0447 & 0.0866 \\
\addlinespace
African              & 0.1728 & 0.2105 & 0.0933 & 0.1205 & 0.0536 & 0.0840 \\
Asian (East)         & 0.1543 & 0.2017 & 0.0792 & 0.1120 & 0.0595 & 0.0929 \\
Asian (South)        & 0.2115 & 0.2456 & 0.1021 & 0.1405 & 0.0340 & 0.1139 \\
Caucasian (American)  & 0.1360 & 0.1942 & 0.0642 & 0.0883 & 0.0185 & 0.0323 \\
Caucasian (European)  & 0.1377 & 0.1899 & 0.0629 & 0.0875 & 0.0321 & 0.0592 \\
\bottomrule
\end{tabular}
\end{table}

The gender and descent pattern matches POP's and attack-vc's direction
throughout, not AntiFake's inverted one: women have a higher mean EER
than men in every system (SV2TTS 0.2398 vs.\ 0.1769; Seed-VC 0.1272
vs.\ 0.0923; F5-TTS 0.0866 vs.\ 0.0664), and the two Caucasian descent
groups remain among the \emph{lowest}-EER (least-protected) groups
under every system, the same ordering already present in each system's
own clean-audio ceiling (\S\ref{sec:protective-efficacy-demographics}).
This confirms, for a third mechanistically distinct cloaking method,
that this demographic pattern in protective efficacy is inherited from
the downstream synthesizer/verifier pipeline's own pre-existing bias
rather than introduced by any particular cloak's optimization target --
the same conclusion \S\ref{sec:protective-efficacy-demographics} and
\S\ref{sec:protective-efficacy-demographics-attackvc} already reached
for POP and attack-vc, and the same contrast
\S\ref{sec:protective-efficacy-demographics-antifake} draws against
AntiFake's inverted, decoy-attack-specific pattern. One partial
exception: Asian (South) is the highest-EER group under F5-TTS
(0.1139) by a wide margin over every other group (0.0323--0.0929),
consistent with that group's already-documented status as the
highest-EER group under F5-TTS's own clean-audio ceiling pattern
elsewhere in this paper's restoration-study demographics
(\S\ref{sec:restoration-seedvc-f5tts}) -- a recurring group-specific
sensitivity under F5-TTS specifically, not unique to ProtectYourAudio.

\section{Ablation: Why Is Seed-VC So Much Less Affected?}
\label{sec:ablation-seedvc-mechanism}

Both protective-efficacy evaluations above (\S\ref{sec:protective-efficacy},
\S\ref{sec:protective-efficacy-attackvc}) end on the same open question:
both cloaks meaningfully degrade SV2TTS but barely move Seed-VC, and we
offered a hypothesis -- that POP's and attack-vc's perturbations
transfer better to systems whose speaker conditioning resembles the
surrogate/target used to craft the attack -- without testing it
directly. Note that this Seed-VC-resistance pattern is specific to POP
and attack-vc: AntiFake, evaluated in
\S\ref{sec:protective-efficacy-antifake}, affects Seed-VC substantially
(EER up 116--1{,}485\% relative). This is not a contradiction --
\S\ref{sec:ablation-embedding-probe} shows it is the same
decoder-tolerance mechanism at a larger dose -- but it means Seed-VC's
apparent resistance should be read as a magnitude-dependent tolerance,
not a fixed property of the system. This section reports two checks
against the original hypothesis:
one architectural (does either cloak share an encoder with Seed-VC's
or SV2TTS's own identity representation?) and one empirical (does the
perturbation actually move each system's own embedding, independent of
whether the final synthesized clone is affected?).

\subsection{Encoder architecture identity}
\label{sec:ablation-encoder-lineage}
We inspected the actual model code behind each component to determine
whether any of them share an architecture family or a literal
checkpoint, since a shared encoder would imply near-white-box transfer
rather than genuine cross-architecture generalization.
Table~\ref{tab:encoder-lineage} summarizes what each component actually
uses.

\begin{table}[H]
\centering
\caption{Speaker/identity encoder architecture used by each cloak's
surrogate/target and by each downstream system's own conditioning
mechanism.}
\label{tab:encoder-lineage}
\begin{tabular}{p{3.4cm}p{4.4cm}p{5.6cm}}
\toprule
Component & Architecture & Evidence \\
\midrule
SV2TTS's own speaker encoder & GE2E LSTM d-vector (256-d, 3-layer LSTM) & \texttt{synth/sv2tts/encoder/model.py} \\
attack-vc's target encoder (embedding-attack variant) & AdaIN-VC's bespoke conv-bank speaker encoder & \texttt{cloak/attack\_vc\_vendored/models.py:129--215} \\
POP's VITS surrogate's speaker conditioning & No encoder network: a plain learned per-speaker embedding lookup table over the surrogate's 50 fixed training speakers & \texttt{cloak/pop\_vendored/vits/models.py:456--457} \\
Seed-VC's own speaker/style encoder & CAMPPlus (D-TDNN, 80-dim fbank in, 192-d embedding out; pretrained FunASR checkpoint) & \texttt{synth/seedvc\_vendored/seed\_vc\_wrapper.py:128--131} \\
\bottomrule
\end{tabular}
\end{table}

None of the four share an architecture family with each other. This
already weakens the strongest form of a ``shared encoder'' explanation
for SV2TTS being the most affected system: neither attack-vc's AdaIN-VC
conv-bank encoder nor POP's VITS surrogate is architecturally related
to SV2TTS's own GE2E encoder, so SV2TTS's larger protective-efficacy
shift is not simple white-box transfer. More notably, POP's surrogate
has no generalizable speaker encoder at all -- VITS conditions on
identity via a closed lookup table indexed by a fixed set of 50
training-time speaker IDs, not a network that extracts identity from
arbitrary reference audio. POP's perturbation is therefore not crafted
against a transferable identity-embedding space in the first place,
which means any cross-system transfer it does exhibit must operate
through some other channel (e.g. generic acoustic/spectral disruption)
rather than embedding-space proximity to a downstream system's own
encoder.

\subsection{Embedding-level transfer probe}
\label{sec:ablation-embedding-probe}
The architectural check above rules out literal encoder reuse but does
not tell us whether the perturbation nonetheless moves each downstream
system's own identity representation. We tested this directly,
bypassing full synthesis entirely: for a random sample of 500 cloaked
clips per cloaking method, we computed each system's own embedding
(SV2TTS's GE2E encoder; Seed-VC's CAMPPlus encoder) on both the clean
original and the cloaked recording, and measured the cosine similarity
between the two. A value near 1.0 means the cloak barely moved that
system's own embedding; a lower value means it did, independent of how
robust that system's downstream decoder turns out to be to such a
shift.

\begin{table}[H]
\centering
\caption{Cosine similarity between each system's own embedding of the
clean original vs.\ the cloaked recording (n=500 per cell, no synthesis
involved). Lower = the cloak moved that system's own identity
representation more.}
\label{tab:embedding-probe}
\begin{tabular}{lcccc}
\toprule
System & Cloak & Mean & Median & Min \\
\midrule
SV2TTS (GE2E)     & POP       & 0.9917 & 0.9949 & 0.7875 \\
SV2TTS (GE2E)     & attack-vc & 0.9836 & 0.9886 & 0.8744 \\
SV2TTS (GE2E)     & AntiFake  & 0.4051 & 0.3864 & 0.1447 \\
Seed-VC (CAMPPlus) & POP      & 0.9787 & 0.9825 & 0.8038 \\
Seed-VC (CAMPPlus) & attack-vc & 0.8955 & 0.9153 & 0.5605 \\
Seed-VC (CAMPPlus) & AntiFake  & 0.4763 & 0.4765 & 0.0478 \\
\bottomrule
\end{tabular}
\end{table}

The result contradicts the simplest version of our hypothesis. If
Seed-VC's near-immunity in Tables~\ref{tab:protective-efficacy} and
\ref{tab:protective-efficacy-attackvc} were caused by the perturbation
failing to transfer to its own encoder, we would expect Seed-VC's
cosine similarity to be closer to 1.0 than SV2TTS's. Instead the
opposite holds, most sharply for attack-vc: Seed-VC's own CAMPPlus
embedding shifts substantially \emph{more} than SV2TTS's GE2E embedding
does (mean cosine similarity 0.8955 vs.\ 0.9836; Seed-VC's standard
deviation is also over $4\times$ larger, 0.0665 vs.\ 0.0160, with a
minimum as low as 0.56 -- some individual clips move Seed-VC's
embedding drastically). Yet this same cloak barely moves Seed-VC's
end-to-end EER (Table~\ref{tab:protective-efficacy-attackvc}: up only
2--6\% relative, the smallest shift of the three downstream systems).
POP shows the same direction of effect, though smaller in magnitude
(Seed-VC mean 0.9787 vs.\ SV2TTS 0.9917).

Taken together with \S\ref{sec:ablation-encoder-lineage}, this points
away from an encoder-transfer explanation and toward a
\emph{decoder-robustness} explanation: the perturbation does reach
Seed-VC's own identity representation, in some cases more than it
reaches SV2TTS's, but Seed-VC's DiT-based conversion decoder appears
substantially more tolerant of a given style-embedding perturbation
than SV2TTS's synthesizer is, so the shifted embedding does not
translate into a correspondingly degraded clone. We had originally
framed the "why is Seed-VC resistant" question in
\S\ref{sec:protective-efficacy} around whether the attack's surrogate
architecture (TTS vs.\ voice conversion) transfers to the downstream
system's \emph{conditioning mechanism}; this ablation instead locates
the effect downstream of conditioning, in how the decoder \emph{uses}
that (measurably perturbed) conditioning signal.

This decoder-robustness hypothesis predicts a specific, testable
consequence: if Seed-VC's resistance to POP and attack-vc is a
tolerance \emph{threshold} rather than blanket immunity, a cloak that
displaces the embedding far enough should eventually break through it.
AntiFake is exactly that test, and it confirms the prediction. Repeating
the same probe on AntiFake-cloaked audio (n=500, same protocol) gives a
mean cosine similarity of 0.4051 for SV2TTS's GE2E embedding and 0.4763
for Seed-VC's CAMPPlus embedding -- roughly $2$--$4\times$ further from
1.0 than POP's or attack-vc's worst case (0.5605, attack-vc's Seed-VC
minimum) on either encoder. Both systems' embeddings are displaced far
more by AntiFake than by either other cloak, and this time Seed-VC's
end-to-end EER moves substantially too
(Table~\ref{tab:protective-efficacy-antifake}: up 116--1{,}485\%
relative, versus 2--15\% for POP/attack-vc). This is a dose-response
result, not a new mechanism: POP's and attack-vc's embedding
displacements (cosine similarity 0.88--0.99) stayed within whatever
tolerance Seed-VC's decoder has, while AntiFake's substantially larger
displacement (cosine similarity 0.41--0.48) exceeds it. Seed-VC's
apparent resistance in \S\ref{sec:protective-efficacy}/\S\ref{sec:protective-efficacy-attackvc}
was therefore never blanket immunity to any perturbation reaching its
encoder -- it was tolerance up to a magnitude that POP and attack-vc
happened not to cross and that AntiFake does. We still have not
isolated the decoder-robustness mechanism directly (e.g.\ feeding
Seed-VC's decoder a synthetically perturbed embedding of matched
magnitude, independent of any actual cloak, to confirm the threshold
account in isolation), so this remains a well-supported hypothesis
rather than a fully confirmed mechanism.

\subsection{F5-TTS's reference transcription is part of the attack surface, not a confound to remove}
\label{sec:ablation-f5tts-transcript}

Every F5-TTS clone reported in this paper (\S\ref{sec:synth}) is
produced with \texttt{ref\_text=""}: an empty reference transcript,
which tells F5-TTS's own inference pipeline to transcribe the
reference recording itself, via an internal Whisper pass, rather than
being handed a transcript. We considered, and rejected, ``fixing''
this by instead supplying F5-TTS the externally-known ground-truth
transcript (\texttt{Dataset/clean\_bonafide/transcripts.json}) for
every clip, regardless of which condition (clean, cloaked, or
restored) is being cloned. We report the reasoning here because the
two options answer genuinely different questions, and conflating them
would have been a real methodological error in the other direction.

A real attacker in this paper's own threat model
(\S\ref{sec:threat}, \S\ref{sec:protective-efficacy}) has access only
to the cloaked (or, for the adaptive-attacker case of
\S\ref{sec:restoration}, restored) recording -- not to a
privately-held ground-truth transcript of what it says. Such an
attacker, using F5-TTS as intended, gets exactly the behavior
\texttt{ref\_text=""} produces: F5-TTS transcribes whatever audio it
is actually handed. If a cloaking or restoration perturbation degrades
that transcription -- and we confirmed it can:
\texttt{finish\_antifake\_f5tts\_stragglers.py} documents 23
AntiFake-cloaked clips whose Whisper transcription returned a
near-empty string (one clip's transcript was literally
\texttt{"..."}), driving F5-TTS's duration heuristic far enough past
its internal buffer size to crash outright -- that degradation is not
an artifact to be corrected out of the measurement. It is a real
consequence a real attacker would experience when pointing an
off-the-shelf cloning system at cloaked audio, and any resulting drop
in clone quality is legitimately part of the cloak's protective
effect, exactly as a filter that happens to also disrupt a
downstream system's unrelated preprocessing step is still a real
disruption. Supplying the correct transcript externally would instead
measure a strictly easier, unrealistic scenario -- an attacker with
privileged knowledge no real attacker in this threat model has -- and
we do not use it for any of this paper's primary results.

The one place we did use it is the 23 crashed AntiFake clips above:
without some intervention, F5-TTS produces no clone at all for those
clips, which is a missing data point rather than a scored trial one
way or the other. We regenerated only those 23 with the ground-truth
transcript so that a trial exists to score, and flag this explicitly
as a deviation specific to those 23 clips (out of 4,844 AntiFake/F5-TTS
combinations) rather than a general practice; every other F5-TTS clone
in this paper, cloaked, restored, or clean, uses F5-TTS's own default
transcription behavior unchanged. We also built (but, per the
reasoning above, do not use to alter any reported result) a
transcript-supplied variant of the synthesis pipeline
(\texttt{synthesize\_f5tts.py --ref-transcripts}) as a diagnostic tool,
should a future study want to separately quantify how much of F5-TTS's
apparent sensitivity to a given cloak is mediated through its
transcription step specifically versus its speaker-identity
conditioning; we did not run it to completion and report no numbers
from it here.

This consideration is specific to F5-TTS among the three downstream
systems in this paper: SV2TTS's speaker encoder extracts an identity
embedding from the reference audio without transcribing it, and
Seed-VC performs voice conversion directly on the reference's acoustic
content, so neither requires or produces a transcript of the reference
clip at all, and neither has an analogous attack surface to consider.

\subsection{Trial-level identifiers and matched impostor sampling}
\label{sec:ablation-trial-ids}

Two further limitations of the original scoring pipeline, unrelated to
the F5-TTS-specific issue above, affect every genuine/impostor trial
reported anywhere in this paper (\S\ref{sec:threat}), and we document
them here because correcting them is an active, in-progress part of
finalizing this paper rather than something we can yet present as
uniformly resolved throughout.

\paragraph{Missing clip identity.} The original score files recorded
only a similarity score, a genuine/impostor label, and a demographic
group per trial -- never which clip, or which specific impostor
recording, produced that row. This made two things impossible after
the fact: auditing which clips had silently failed and been dropped
(\S\ref{sec:cloaking}'s engineering-obstacles discussion applies
identically to scoring, not only cloak generation), and constructing a
\emph{matched} before/after comparison on the same clip across two
conditions (e.g.\ a cloak's cloaked-source result against the
clean-audio ceiling, or a restoration technique's result against its
cloak's own no-restoration baseline) rather than comparing two
independently-estimated pooled statistics. We now record, per trial,
the clip's own filename, whether it is the genuine trial or which
numbered impostor trial, and the specific impostor recording used, so
every reported number can be traced back to and reconstructed from the
individual trials underlying it.

\paragraph{Impostor sampling not shared across verifiers or
conditions.} For each clip, \S\ref{sec:threat} draws $k=10$ impostor
recordings uniformly at random. The original implementation drew all
of a condition's impostors from a single random-number generator,
seeded once and never reset, that ran continuously across every clip
and -- critically -- across all three verifiers' passes over those
same clips, one verifier's full pass at a time. Two consequences
follow, neither of which is a sampling-bias problem in expectation but
both of which undermine the kind of matched, clip-level comparison
\S\ref{sec:ablation-trial-ids} above and the restoration study's
reversal criterion (\S\ref{sec:restoration-summary}) depend on. First,
because each verifier's pass continues the same generator rather than
restarting it, a given clip's ten impostors differed from one verifier
to the next: SB-ECAPA, Resemblyzer, and WavLM were, in effect,
answering three different questions about the same clip rather than
the same question three times. Second, because a clip whose genuine
trial failed contributed no draws to the generator at all, every clip
processed after any failure was silently shifted onto a different
random draw than it would otherwise have received -- so two runs of
the same condition, or two conditions that happened to lose different
clips to transient failures, were not guaranteed to compare a given
clip against the same impostor set even when both intended to. We
corrected this by seeding a fresh generator per clip, keyed on the
clip's own filename (and a fixed global seed), so a given clip is
compared against the identical ten impostor recordings in every
condition, by every verifier, regardless of processing order or which
other clips succeeded or failed elsewhere in the run. We validated
this directly: the same clip, scored by all three verifiers, and
scored again inside an entirely different, non-overlapping batch of
clips, produces byte-identical impostor selections in every case.

\paragraph{Four evaluation objectives and what each is compared
against.} The dataset produces four distinct kinds of audio --
\texttt{cloaked\_bonafide} (each cloak's direct output),
\texttt{cloaked\_synthesize} (a cloaked clip re-cloned by a TTS/VC
system), \texttt{restoration techniques} (a cloaked clip after a
restoration/attack technique has tried to strip the cloak, prior to
any re-cloning), and \texttt{restored\_synthesize} (that restored
audio re-cloned) -- and it is easy to conflate what each is measured
against, so we state it explicitly:
\begin{enumerate}
  \item \textbf{Imperceptibility} (\texttt{cloaked\_bonafide}, \S\ref{sec:results-antifake}
    and its counterparts): each cloak's output is compared against the
    \emph{clean original recording it was cloaked from} (same
    recording, same words) -- the easiest possible case for a
    verifier. We additionally report a stricter, cross-recording
    variant below in which the genuine reference is instead a
    \emph{different} recording of the same speaker.
  \item \textbf{Protective efficacy} (\texttt{cloaked\_synthesize},
    \S\ref{sec:protective-efficacy-demographics} and its
    counterparts): a cloaked clip, re-cloned by SV2TTS, Seed-VC, or
    F5-TTS, is compared against \emph{that same synthesizer's
    clean-audio ceiling} (Table~\ref{tab:synth-ceiling}: the identical
    synthesizer cloning the \emph{uncloaked} original) -- the
    best-case, no-protection outcome a real attacker would otherwise
    get. A cloak is protective if the cloaked-source clone's EER is
    significantly \emph{higher} (TAR significantly \emph{lower}) than
    this ceiling, paired on matched clips
    (\texttt{protection\_verdict}, \S\ref{sec:ablation-trial-ids}
    above).
  \item \textbf{Restoration audio quality} (\texttt{restoration
    techniques}): the raw restored signal, \emph{before} any
    re-cloning, is not itself put through a speaker verifier -- it is
    scored only on signal-quality metrics (STOI, PESQ, delay-corrected
    SI-SDR, \S\ref{sec:quality-metrics}) against the cloaked audio it
    was restored from, to characterize how much a restoration
    technique itself degrades the signal. It is the \emph{input} to
    the next objective, not an endpoint in itself.
  \item \textbf{Restoration reversal} (\texttt{restored\_synthesize},
    \S\ref{sec:restoration}): the restored clip, re-cloned by the same
    synthesizer, is compared against \emph{that cloak's own
    no-restoration baseline} -- i.e.\ objective 2's cloaked-source
    clone, \emph{not} the clean-audio ceiling. This is deliberate: the
    question a restoration technique must answer is whether it gives
    an attacker who already possesses only the cloaked audio (our
    threat model, \S\ref{sec:threat}) any advantage over simply
    cloning that cloaked audio directly, not whether it recovers
    unprotected, never-cloaked performance. Reversal is significant
    (\texttt{reversal\_verdict}) only when the restored clone's EER is
    significantly \emph{lower} than this no-restoration baseline,
    paired on matched clips, on every verifier.
\end{enumerate}
Objectives 2 and 4 both report EER/TAR numbers that are, as
throughout this paper, computed against the clean original as the
genuine reference in every individual trial (\S\ref{sec:threat}); what
differs between them is only which \emph{other} condition each is
statistically paired against for its protection/reversal verdict.

\paragraph{Cross-recording imperceptibility: a stricter test.} The
imperceptibility results in
\S\ref{sec:results}, \S\ref{sec:results-attackvc},
\S\ref{sec:results-antifake}, and \S\ref{sec:results-protectyouraudio}
compare each
cloaked clip against the exact recording it was cloaked from -- same
room, same microphone, same words -- which is the easiest case for any
verifier and does not reflect how a verifier is actually used (against
an enrolled reference recorded at another time). We therefore also
score every cloak's output against a \emph{different} recording of the
same speaker as the genuine trial (impostor trials unchanged, same
per-clip seeded sample as everywhere else in this paper), on the 4,801
of 4,867 clips (437 LibriSpeech, 498 ASVspoof2021, 3,866 FakeAVCeleb)
whose speaker has a second recording available, with the identical
trials scored against the clean original as a baseline. A cloak is
imperceptible by this stricter test only if its cross-recording EER is
statistically indistinguishable from (not higher than) this clean
baseline. Table~\ref{tab:crossrecording} reports the paired
cluster-bootstrap difference (cloaked minus clean baseline) in EER,
per cloak, source corpus, and verifier -- these are the same four
cloaks as \S\ref{sec:results} (POP), \S\ref{sec:results-attackvc},
\S\ref{sec:results-antifake}, and \S\ref{sec:results-protectyouraudio};
FakeAVCeleb's baseline EER is already near chance ($\approx$0.27--0.42)
on all three verifiers, so it has little room to show a difference
even when one exists, while LibriSpeech and ASVspoof2021 have low
clean-baseline EER and can. AntiFake fails this test badly and
consistently: EER rises by 0.06--0.58 (all differences significant at
$p<0.001$) on every corpus and every verifier, including the
otherwise-insensitive FakeAVCeleb baseline -- an order of magnitude
larger than any other cloak. attack-vc and ProtectYourAudio show
smaller but still real and significant increases concentrated on
WavLM and on the cleaner LibriSpeech/ASVspoof2021 baselines (e.g.\
attack-vc on ASVspoof2021/WavLM: $+0.0435$, CI $[0.0297, 0.0575]$;
ProtectYourAudio on ASVspoof2021/WavLM: $+0.0354$, CI $[0.0246,
0.0479]$). POP is the only cloak with no consistent significant
degradation across corpora and verifiers under this stricter test.
This suggests that AntiFake's same-recording imperceptibility numbers
above substantially overstate how imperceptible the cloak would be
against a verifier's actual enrolled reference audio, while POP's
same-recording numbers are a fair approximation of the harder,
more realistic case.

\begin{table}[H]
\centering
\small
\caption{Cross-recording imperceptibility: paired cluster-bootstrap
difference in EER (cloaked minus clean baseline, same clips), genuine
trial = a different recording of the same speaker. Bold: $p<0.05$.}
\label{tab:crossrecording}
\begin{tabular}{llrrrr}
\toprule
Cloak & Source & SV & $\Delta$EER & 95\% CI & $p$ \\
\midrule
POP & LibriSpeech & SB-ECAPA & 0.0006 & [$-$0.0003, 0.0026] & 0.948 \\
POP & LibriSpeech & Resemblyzer & $-$0.0005 & [$-$0.0042, 0.0022] & 0.796 \\
POP & LibriSpeech & WavLM & 0.0059 & [$-$0.0025, 0.0080] & 0.376 \\
POP & ASVspoof2021 & SB-ECAPA & 0.0057 & [$-$0.0019, 0.0093] & 0.136 \\
POP & ASVspoof2021 & Resemblyzer & $-$0.0001 & [$-$0.0048, 0.0087] & 0.572 \\
POP & ASVspoof2021 & WavLM & $-$0.0032 & [$-$0.0117, 0.0040] & 0.352 \\
POP & FakeAVCeleb & SB-ECAPA & 0.0007 & [$-$0.0034, 0.0053] & 0.656 \\
POP & FakeAVCeleb & Resemblyzer & \textbf{$-$0.0038} & [$-$0.0064, $-$0.0005] & \textbf{0.024} \\
POP & FakeAVCeleb & WavLM & 0.0005 & [$-$0.0023, 0.0033] & 0.848 \\
\addlinespace
attack-vc & LibriSpeech & SB-ECAPA & 0.0010 & [$-$0.0005, 0.0035] & 0.888 \\
attack-vc & LibriSpeech & Resemblyzer & $-$0.0006 & [$-$0.0041, 0.0026] & 0.864 \\
attack-vc & LibriSpeech & WavLM & \textbf{0.0126} & [0.0041, 0.0201] & \textbf{0.008} \\
attack-vc & ASVspoof2021 & SB-ECAPA & \textbf{0.0159} & [0.0043, 0.0264] & \textbf{0.000} \\
attack-vc & ASVspoof2021 & Resemblyzer & \textbf{0.0142} & [0.0052, 0.0271] & \textbf{0.004} \\
attack-vc & ASVspoof2021 & WavLM & \textbf{0.0435} & [0.0297, 0.0575] & \textbf{0.000} \\
attack-vc & FakeAVCeleb & SB-ECAPA & $-$0.0025 & [$-$0.0080, 0.0020] & 0.244 \\
attack-vc & FakeAVCeleb & Resemblyzer & $-$0.0034 & [$-$0.0061, 0.0008] & 0.112 \\
attack-vc & FakeAVCeleb & WavLM & \textbf{0.0047} & [0.0011, 0.0081] & \textbf{0.012} \\
\addlinespace
AntiFake & LibriSpeech & SB-ECAPA & \textbf{0.2977} & [0.2668, 0.3304] & \textbf{0.000} \\
AntiFake & LibriSpeech & Resemblyzer & \textbf{0.4209} & [0.4007, 0.4465] & \textbf{0.000} \\
AntiFake & LibriSpeech & WavLM & \textbf{0.1183} & [0.0983, 0.1326] & \textbf{0.000} \\
AntiFake & ASVspoof2021 & SB-ECAPA & \textbf{0.1965} & [0.1733, 0.2136] & \textbf{0.000} \\
AntiFake & ASVspoof2021 & Resemblyzer & \textbf{0.5828} & [0.5608, 0.6106] & \textbf{0.000} \\
AntiFake & ASVspoof2021 & WavLM & \textbf{0.0794} & [0.0633, 0.0930] & \textbf{0.000} \\
AntiFake & FakeAVCeleb & SB-ECAPA & \textbf{0.0991} & [0.0857, 0.1118] & \textbf{0.000} \\
AntiFake & FakeAVCeleb & Resemblyzer & \textbf{0.1717} & [0.1604, 0.1851] & \textbf{0.000} \\
AntiFake & FakeAVCeleb & WavLM & \textbf{0.0582} & [0.0504, 0.0673] & \textbf{0.000} \\
\addlinespace
ProtectYourAudio & LibriSpeech & SB-ECAPA & 0.0031 & [$-$0.0002, 0.0063] & 0.336 \\
ProtectYourAudio & LibriSpeech & Resemblyzer & 0.0043 & [0.0000, 0.0096] & 0.064 \\
ProtectYourAudio & LibriSpeech & WavLM & \textbf{0.0152} & [0.0069, 0.0230] & \textbf{0.000} \\
ProtectYourAudio & ASVspoof2021 & SB-ECAPA & 0.0060 & [$-$0.0016, 0.0112] & 0.136 \\
ProtectYourAudio & ASVspoof2021 & Resemblyzer & 0.0062 & [$-$0.0019, 0.0181] & 0.104 \\
ProtectYourAudio & ASVspoof2021 & WavLM & \textbf{0.0354} & [0.0246, 0.0479] & \textbf{0.000} \\
ProtectYourAudio & FakeAVCeleb & SB-ECAPA & 0.0008 & [$-$0.0054, 0.0077] & 0.816 \\
ProtectYourAudio & FakeAVCeleb & Resemblyzer & \textbf{0.0063} & [0.0016, 0.0114] & \textbf{0.008} \\
ProtectYourAudio & FakeAVCeleb & WavLM & \textbf{0.0073} & [0.0026, 0.0117] & \textbf{0.000} \\
\bottomrule
\end{tabular}
\end{table}

\paragraph{Status at time of writing.} Both fixes (trial identifiers
and per-clip matched impostor sampling) are implemented, validated,
and now applied throughout: all 18 primary conditions (the four
cloaks' own imperceptibility results, plus the 12 cloaked-source
protective-efficacy results and the three clean-audio ceilings they
are compared against) and all 143 restoration-study conditions
(\S\ref{sec:restoration}; ProtectYourAudio is not part of the
restoration study -- see \S\ref{sec:restoration}'s scope note) have
completed rescoring. The clean-audio SV2TTS ceiling was also extended
from its original 1,500-clip subset to the full 4,867-clip corpus
(matching F5-TTS's and Seed-VC's coverage, so every protective-efficacy
comparison is paired on matched clips regardless of synthesizer); an
earlier synthesis pass fell 1,015 clips short of this after transient
GPU-memory contention, which we caught before any number derived from
it was reported and corrected by resynthesizing the missing clips. The
cross-recording imperceptibility results above are likewise complete
for all four cloaks. The paired protection/reversal verdicts
throughout this paper are being recomputed one final time against this
now-complete, fully matched dataset as this version is finalized.

\section{Theoretical Framework: Cloaking and Restoration as Constrained Optimization}
\label{sec:framework}

The sections above describe each cloaking technique's mechanism in
prose (\S\ref{sec:cloaking}) and report its empirical effect. This
section makes the underlying optimization problem explicit and common
to all three gradient-based techniques, so that (i) POP, attack-vc,
and AntiFake can be read as three points in one parameter space rather
than three unrelated algorithms, and (ii) the restoration study's
uniform failure (\S\ref{sec:restoration}) can be explained
\emph{predictively} from the structure of that space, rather than only
described post hoc. The purpose of formalizing this is to turn two
empirical observations already established in this paper -- that
restoration never reversed either cloak, and that Seed-VC's resistance
is a magnitude threshold rather than blanket immunity
(\S\ref{sec:ablation-embedding-probe}) -- into concrete, falsifiable
directions for improving both the attacker's (cloaking) and the
defender's (restoration) side, which we give in
\S\ref{sec:framework-improve-cloaking} and
\S\ref{sec:framework-improve-restoration}.

\subsection{Notation}
\label{sec:framework-notation}
Let $x \in \mathbb{R}^T$ be a clean recording of speaker $s$, and
$\mathrm{Enc}_k(\cdot)$ the speaker-embedding function of encoder $k$
(a verifier from Table~\ref{tab:verifiers}, or a surrogate/downstream
encoder from \S\ref{sec:synth}). A cloak produces
\[
\hat{x} = x + \delta, \qquad \|\delta\|_\infty \leq \varepsilon,
\]
an additive perturbation bounded in a fixed $\ell_\infty$ ball (POP's
own stated budget is $\varepsilon = 8/255$; attack-vc and AntiFake use
comparably small budgets in the same spirit, chosen so that $\hat{x}$
remains close to $x$ in the perceptual metrics of
\S\ref{sec:metrics}). A downstream cloning system with encoder
$\mathrm{Enc}$ and decoder $\mathrm{Dec}$ produces a clone
$c = \mathrm{Dec}(\mathrm{Enc}(\hat{x}), \cdot)$; a verifier scores
identity via cosine similarity of $\mathrm{Enc}_v(\cdot)$ embeddings.
Imperceptibility (\S\ref{sec:results}) asks whether
$\mathrm{Enc}_v(\hat{x}) \approx \mathrm{Enc}_v(x)$; protective
efficacy (\S\ref{sec:protective-efficacy}) asks whether
$\mathrm{Enc}_v(c) \not\approx \mathrm{Enc}_v(x)$.

\subsection{The three cloaks as one constrained optimization problem}
\label{sec:framework-cloaks}
Every gradient-based technique in this study solves a variant of
\[
\delta^\ast = \argmin_{\|\delta\|_\infty \leq \varepsilon}
\; \lambda_{\mathrm{recon}}\, \mathcal{L}_{\mathrm{recon}}(x+\delta)
\;+\; \lambda_{\mathrm{emb}} \sum_{k \in K} w_k \,
\big\| \mathrm{Enc}_k(x+\delta) - \mathrm{Enc}_k(x_{\mathrm{decoy}}) \big\|_2^2,
\]
via projected gradient descent, differing only in which term is
active, over how many surrogate encoders $K$, and whether a decoy
target $x_{\mathrm{decoy}}$ (a real recording of a different, randomly
chosen speaker) is used at all. Table~\ref{tab:framework-cloaks}
locates each technique in this space.

\begin{table}[H]
\centering
\caption{The three gradient-based cloaking techniques as instances of
the same constrained optimization problem.}
\label{tab:framework-cloaks}
\begin{tabular}{lccc p{4.3cm}}
\toprule
Technique & $\lambda_{\mathrm{recon}}$ & $\lambda_{\mathrm{emb}}$ & $|K|$ & Consequence \\
\midrule
POP       & $>0$ & $0$ & --- & Minimizes a VITS surrogate's own
  reconstruction loss (an error-minimizing, ``unlearnable-example''
  objective) with no explicit embedding-displacement term at all. \\
attack-vc & $0$   & $>0$ & 1 & Single-encoder decoy-target embedding
  attack (AdaIN-VC's own encoder); no reconstruction term, so nothing
  constrains the objective to also minimize signal distortion beyond
  the $\varepsilon$ budget itself. \\
AntiFake  & $0$   & $>0$ & 3 & Same decoy-target objective as
  attack-vc, but jointly over an ensemble of three encoders (RTVC,
  AutoVC, YourTTS) for black-box transferability. \\
\bottomrule
\end{tabular}
\end{table}

This single equation explains three results already reported
separately. \textbf{First}, POP's near-perfect imperceptibility
(EER~$\leq$~0.16\%, \S\ref{sec:results}) alongside its
system-dependent protective effect (strong against SV2TTS, marginal
against Seed-VC, \S\ref{sec:protective-efficacy}) follows directly
from $\lambda_{\mathrm{emb}}=0$: nothing in POP's objective explicitly
displaces \emph{any} speaker embedding, so whatever protective effect
it has is an incidental side effect of degrading the VITS surrogate's
reconstruction, which transfers unevenly to encoders it was never
optimized against. \textbf{Second}, attack-vc and AntiFake's
imperceptibility failure risk follows from $\lambda_{\mathrm{recon}}=0$:
neither objective contains any term that penalizes the perturbation
for being audible, so imperceptibility is only as good as the
$\varepsilon$ budget happens to allow, with no explicit pressure toward
the low end of that budget. attack-vc's budget lands close enough to
transparent (EER~$\leq$~2.1\%) that this risk is not realized in
practice; AntiFake's does not (EER~22--66\%,
\S\ref{sec:results-antifake}) -- the same objective family, at a
different point in the same space, produces qualitatively different
imperceptibility outcomes. \textbf{Third}, the dose-response finding of
\S\ref{sec:ablation-embedding-probe} is exactly the statement that
$|K|=3$ (AntiFake's ensemble) pushes $\|\delta\|$'s effective embedding
displacement past a magnitude that $|K|=1$ (attack-vc) and
$\lambda_{\mathrm{emb}}=0$ (POP) do not reach for Seed-VC's decoder --
the three techniques are not qualitatively different mechanisms so much
as three settings of $(\lambda_{\mathrm{recon}}, \lambda_{\mathrm{emb}}, |K|)$
that happen to land on different sides of that threshold.

\subsection{Restoration as approximate inversion}
\label{sec:framework-restoration}
Every technique in \S\ref{sec:restoration} is an attempt to construct
an operator $R$ such that $R(\hat{x}) = R(x+\delta) \approx x$: recover
something closer to the clean signal than the cloaked one, without
access to $x$ or $\delta$ individually. The ten techniques fall into a
small number of operator classes, summarized in
Table~\ref{tab:framework-restoration-classes}.

\begin{table}[H]
\centering
\caption{The ten restoration techniques grouped by operator class.}
\label{tab:framework-restoration-classes}
\begin{tabular}{p{3.6cm}p{8.4cm}}
\toprule
Operator class & Techniques \\
\midrule
Linear time-invariant filtering & Low-pass+gain, high-pass,
  adaptive-filter-centroid, resampling (down/up-sampling's implicit
  anti-alias filter) \\
Lossy-basis projection & Mel-spectrogram inversion (project onto an
  80-mel, Griffin-Lim-invertible basis) \\
Nonlinear scalar quantization & Quantization-based inversion (8-bit
  round-trip) \\
Nonlinear spectral masking & Spectral subtraction (magnitude-domain
  floor/subtraction against a self-estimated noise profile) \\
Stochastic averaging & Ensemble averaging of perturbed variants
  ($N=5$ jittered copies) \\
Physical-channel simulation & Simulated re-recording (synthetic RIR
  convolution + mic bandpass + noise floor) \\
Generic additive noise & Second-cloak application (cloak-agnostic
  30~dB-SNR Gaussian noise, single pass, no averaging) \\
\bottomrule
\end{tabular}
\end{table}

All ten of these techniques are \emph{$\delta$-agnostic}: their form was
chosen from a generic prior about where perturbations \emph{typically}
live (high
frequency, fine amplitude precision, phase instability under jitter),
not derived from $\delta$ itself or from the attack's own gradient.
This matters because $\delta$ is not an arbitrary noise signal: by
construction, every $\delta$ that satisfies its technique's
imperceptibility goal must already resemble natural speech statistics
closely enough to survive a human listener (or, for AntiFake, at least
survive the $\varepsilon$ budget). Writing $R$'s own distortion of
clean speech as $\varepsilon_R = \|R(x) - x\|$ (nonzero for every
non-identity filter -- exactly what
\S\ref{sec:restoration-lowpass-gain} and
\S\ref{sec:restoration-adaptive-filter} measure implicitly whenever
they report the no-restoration baseline is a better outcome for the
attacker than the restored one), a $\delta$-agnostic, non-adaptive $R$
has no mechanism to preferentially attenuate $\delta$'s energy over
$x$'s: for an approximately linear $R$,
\[
R(\hat{x}) - x \;=\; \underbrace{\big(R(x) - x\big)}_{\varepsilon_R,\ \text{always present}}
\;+\; R(\delta),
\]
so $R$ helps only if $R$ attenuates $\delta$ enough to more than offset
the distortion $\varepsilon_R$ it adds to $x$ regardless of $\delta$'s
presence. Because $\delta$ was optimized to look like speech, generic
$R$ has no reason to attenuate it more than it attenuates $x$, and
every restoration result reported so far (\S\ref{sec:restoration-lowpass-gain},
\S\ref{sec:restoration-adaptive-filter}) shows EER moving \emph{further}
from the clean-audio ceiling after restoration, not closer -- consistent
with $\varepsilon_R$ dominating rather than being offset. This is a
predictive account, not only a post hoc one: it implies technique-by-technique
that the size of $\varepsilon_R$, which can be measured directly (apply
$R$ to \emph{clean, uncloaked} $x$ and measure the embedding-cosine
shift or STOI/PESQ drop it causes on its own, with no cloak involved at
all), should predict how much worse that technique makes protective
efficacy -- a testable, cheap-to-check consequence of the framework we
did not previously have a reason to check for
(\S\ref{sec:framework-improve-restoration} makes this concrete).

Second-cloak application (\S\ref{sec:restoration-secondcloak}) is, on
reflection, not an exception to this argument but a direct test of it.
An earlier implementation of this technique reused a cloak's own attack
machinery against its own already-cloaked output -- which would indeed
have broken the $\delta$-agnostic pattern above, since it required
knowing which specific cloak produced $\delta$ in the first place. We
corrected this: a restoration-side attacker has no way to know or
reproduce the specific cloak used, so the technique we actually report
is cloak-agnostic additive Gaussian noise, identical in spirit to
ensemble averaging's per-variant noise
(\S\ref{sec:restoration-ensemble}) but applied once rather than
averaged across several jittered copies. This makes it, like every
other technique in this study, a direct empirical test of whether a
generic, non-adversarial $\varepsilon_R$ can be made small enough to
avoid dominating $R(\delta)$'s attenuation, at a single, fixed noise
level (30~dB SNR) rather than the $N=5$-variant averaging of
\S\ref{sec:restoration-ensemble}.

\subsection{Improving the cloaking techniques}
\label{sec:framework-improve-cloaking}
The decomposition in \S\ref{sec:framework-cloaks} points to a specific,
testable change for each technique, aimed at the specific weakness the
data already shows.

\paragraph{POP: add an embedding term.} POP's near-perfect
imperceptibility (EER~$\leq$~0.16\%) at $\varepsilon=8/255$ indicates
its current objective is far from exhausting its perceptual budget --
there is room, in principle, to add a small
$\lambda_{\mathrm{emb}} \sum_k \|\mathrm{Enc}_k(x+\delta) -
\mathrm{Enc}_k(x_{\mathrm{decoy}})\|_2^2$ term (attack-vc/AntiFake's
own term) alongside POP's existing reconstruction term, within the
\emph{same} $\varepsilon$ budget, rather than replacing it. This
directly targets POP's demonstrated weakness -- a negligible effect
against Seed-VC (\S\ref{sec:protective-efficacy}) explained by
$\lambda_{\mathrm{emb}}=0$ -- while the reconstruction term is retained
specifically because it is what keeps POP's imperceptibility
near-perfect, something AntiFake's objective has no equivalent
safeguard for. The concrete, falsifiable prediction: a
$\lambda_{\mathrm{recon}} + \lambda_{\mathrm{emb}}$ hybrid should raise
POP's protective efficacy against Seed-VC substantially above its
current 5--15\% relative EER increase, without materially moving its
imperceptibility EER away from its current $\leq$~0.16\%, since the
budget itself is unchanged and only the loss landscape within it is.

\paragraph{attack-vc: ensemble the encoder.} attack-vc's protective
effect is strong against SV2TTS but negligible against Seed-VC and only
moderate against F5-TTS (\S\ref{sec:protective-efficacy-attackvc}) --
exactly the single-encoder ($|K|=1$) row of
Table~\ref{tab:framework-cloaks}. AntiFake's ensemble
($|K|=3$) shows this is not a fundamental limitation of the decoy-target
objective itself but of optimizing it against only one surrogate.
Extending attack-vc's own embedding attack to a small ensemble (e.g.\
$|K|=2$, adding one encoder architecturally distinct from AdaIN-VC's,
to control the added compute against AntiFake's full three-encoder
cost) is a direct test of whether $|K|$, independent of AntiFake's
other objective differences, is what drives cross-system
generalization.

\paragraph{AntiFake: add a regularizer, or find the minimum effective
dose.} AntiFake's protective effect is the largest of any technique
tested (EER up 116--2,199\% relative,
\S\ref{sec:protective-efficacy-antifake}) but its imperceptibility
failure (EER 22--66\%) is the clearest weakness of any technique in
this study, and Table~\ref{tab:framework-cloaks} identifies exactly why:
$\lambda_{\mathrm{recon}}=0$, so nothing in the objective is rewarded
for staying close to $x$ beyond the hard $\varepsilon$ constraint. Two
independent directions follow. \emph{(i)} Add POP's own reconstruction
term back in as a regularizer, trading some embedding-displacement
magnitude for imperceptibility -- the natural dual of the POP proposal
above. \emph{(ii)} More directly, \S\ref{sec:ablation-embedding-probe}
shows the embedding-cosine-similarity values that separate ``protective''
from ``not'' are not close to zero: POP/attack-vc's displacements
(0.88--0.99, insufficient) and AntiFake's (0.41--0.48, sufficient)
leave a wide unexplored middle. Since protective efficacy is
apparently governed by crossing a decoder-tolerance \emph{threshold}
rather than by the raw magnitude of displacement past it, AntiFake's
current objective is not solving the problem the evaluation actually
cares about (maximize $\|\delta\|$'s embedding effect) -- it should
instead solve a feasibility problem: find the minimum $\varepsilon$ (or
minimum PGD iteration count) such that
$\cos(\mathrm{Enc}_k(x+\delta), \mathrm{Enc}_k(x))$ falls below the
empirically observed critical threshold for every downstream system's
encoder simultaneously, and stop there. Both proposals predict the same
outcome -- imperceptibility recovers substantially while protective
efficacy is retained -- via different mechanisms (an explicit penalty
vs.\ an implicit budget reduction), so running both independently and
comparing which recovers more imperceptibility per unit of protective
efficacy lost would itself be informative about which term is doing
the damage.

\subsection{Improving the restoration techniques}
\label{sec:framework-improve-restoration}
\S\ref{sec:framework-restoration}'s argument that $\varepsilon_R$
(a restoration technique's own distortion of \emph{clean}, uncloaked
speech) should predict how much worse that technique makes protective
efficacy suggests a concrete, low-cost addition to the restoration
study's own protocol, plus two structurally different directions for
new restoration candidates.

\paragraph{Add a clean-audio control channel.} None of the ten
restoration results reported in \S\ref{sec:restoration} include the
one control the framework identifies as decisive: applying each
technique to \emph{clean\_bonafide} audio (no cloak at all) and
measuring the same embedding-cosine shift or STOI/PESQ drop used
elsewhere in this paper. This is inexpensive (no cloak generation or
downstream cloning required, just the restoration operator itself
applied once per technique to the existing 4,867-clip clean corpus)
and directly falsifiable: if $\varepsilon_R$ predicts the ranking of
how much each technique \emph{worsens} protective efficacy (as
\S\ref{sec:framework-restoration} argues it should), this control would
let future restoration candidates be screened \emph{before} the
expensive cloak-then-clone-then-verify pipeline is run on them at all,
rather than discovering after the fact -- as happened for
low-pass+gain and adaptive-filter-centroid -- that a technique made
things worse.

\paragraph{Move from $\delta$-agnostic to learned, $\delta$-aware
purification.} Every technique tested so far is fixed-form and
independent of any actual cloak's optimization process (Table
\ref{tab:framework-restoration-classes}). The framework's account of
why they failed --- no mechanism to attenuate $\delta$ preferentially
over $x$ --- applies to any fixed-form $R$, however cleverly
parameterized, and predicts that this whole class of technique is
unlikely to succeed regardless of which one is tried next. A
qualitatively different candidate is a \emph{learned} purification
network $R_\theta$, trained by minimizing
$\mathbb{E}_{x}\big[\|R_\theta(x+\delta(x)) - x\|\big]$ over many
$(x, \delta(x))$ pairs generated by running a cloak's own attack
pipeline against a training corpus -- i.e.\ training $R_\theta$
specifically against the distribution of perturbations a given cloak
produces, rather than against a generic prior about where
perturbations live, analogous to learned adversarial-purification
approaches in the image domain. This is a substantially larger
undertaking than any of the ten techniques implemented so far (it
requires its own training corpus and optimization, not just a
signal-processing pass), but it is the direction the framework
predicts is structurally capable of succeeding where fixed-form
filtering cannot.

\paragraph{A genuinely gradient-informed adaptive attacker remains
untested.} All ten restoration techniques evaluated in this study,
including the corrected second-cloak application
(\S\ref{sec:restoration-secondcloak}), are $\delta$-agnostic: none of
them use any gradient information about the specific cloak that
produced $\hat{x}$, because a real restoration-side attacker generally
does not know which cloak (if any) was applied. A qualitatively
different, \emph{white-box} adaptive attacker -- one who somehow knows
or correctly guesses the cloak and has access to its surrogate model --
could instead run a gray-box ``denoising PGD'': gradient
\emph{ascent} on the same surrogate loss the cloak minimized, i.e.\
$\delta_2^\ast = \argmax_{\|\delta_2\|_\infty \leq \varepsilon}
\mathcal{L}_{\mathrm{recon\ or\ emb}}(\hat{x}+\delta_2)$, an explicit
attempt to push the perturbation back toward the loss landscape's own
optimum at $\delta=0$. This is a stronger and less realistic threat
model than every technique in \S\ref{sec:restoration} (it presumes
knowledge no restoration-side attacker in our threat model has), so we
do not implement it here, but it remains the natural next step for
testing whether \emph{any} restoration strategy -- not just the
generic ones available to a cloak-agnostic attacker -- can reverse
these cloaks' protective effect.

\section{Adaptive-Attacker Restoration: Can Protective Efficacy Be Reversed?}
\label{sec:restoration}

Sections~\ref{sec:protective-efficacy} and
\ref{sec:protective-efficacy-attackvc} treat the downstream cloning
attacker as \emph{naive}: it clones directly from the cloaked
recording, with no knowledge of, or countermeasure against, the cloak
itself. A stronger, \emph{adaptive} attacker who knows a recording may
be cloaked could instead apply a signal-domain restoration technique
to the cloaked audio \emph{before} cloning, attempting to strip out or
average away the cloaking perturbation while leaving enough of the
speaker's true acoustic identity intact for cloning to succeed. This
section describes the ten restoration techniques we implemented to
test that hypothesis, applied to POP-cloaked, attack-vc-cloaked, and
AntiFake-cloaked source audio (4,864, 4,863, and 4,844 clips
respectively) for all ten techniques against all three cloak methods.
SV2TTS cloning-and-verification (the step that actually establishes
whether restoration reverses a cloak's protective effect) is complete
for all 30 (technique, cloak-method) combinations
(Table~\ref{tab:restoration-status} reports per-technique detail), and
we additionally extended cloning-and-verification to Seed-VC and
F5-TTS for all 30 combinations as well (\S\ref{sec:restoration-seedvc-f5tts}) --
90 (technique, cloak-method, synthesizer) combinations in total. We
state explicitly, per technique below, what each SV2TTS result shows,
and \S\ref{sec:restoration-summary}/\S\ref{sec:restoration-seedvc-f5tts}
draw the cross-technique conclusions these 90 combinations support:
restoration never reverses either POP or attack-vc, on any of the three
downstream synthesizers. For AntiFake, on SV2TTS, five of the ten
techniques (low-pass filtering -- now confirmed across its full
$M=1.0$--$2.0$ gain sweep, not just the single point originally
reported -- resampling's downsampling variant, ensemble averaging,
spectral subtraction, and simulated re-recording) produce a real,
statistically-supported -- though still partial, not complete --
movement toward the clean-audio ceiling, while the other five do not
(three show no effect, two are ambiguous across verifiers); the same
five techniques also show this effect, more weakly and less
consistently across verifiers, on Seed-VC and F5-TTS
(\S\ref{sec:restoration-seedvc-f5tts}).

\subsection{Common protocol}
Each technique is applied once to every cloaked recording, producing
one restored candidate per (cloaking method, recording). Testing
whether a candidate actually reverses the cloak's protective effect
then requires the same resynthesis-based protocol used in
Section~\ref{sec:protective-efficacy}: clone the restored candidate
with each of the three downstream systems (SV2TTS, Seed-VC, F5-TTS,
Section~\ref{sec:synth}) and score the result against the genuine
clean\_bonafide original with all three verifiers
(Section~\ref{sec:verifiers}). A restoration technique would count as
having reversed the cloak's protective effect if the resulting clone's
EER/TAR@FAR returns toward the clean-audio ceiling of
Table~\ref{tab:synth-ceiling}, i.e.\ back toward the ``clone succeeds''
regime, undoing the shift documented in
Section~\ref{sec:protective-efficacy}.

\paragraph{Two deliberately different comparisons, against two
deliberately different files.} Each restoration-technique cell in this
section reports two kinds of measurement, and they are computed against
two different points in the pipeline on purpose. EER/TAR@FAR (the
protective-efficacy question above) compares the genuine
clean\_bonafide original against the \emph{final synthesized clone} --
i.e.\ cloak $\to$ restoration technique $\to$ downstream
SV2TTS/Seed-VC/F5-TTS clone -- and asks whether a verifier can still
tell the clone apart from the real speaker. Signal-quality metrics
(STOI, PESQ, SI-SDR, Section~\ref{sec:quality-metrics}) instead compare
the genuine clean\_bonafide original against the \emph{restored audio
itself}, one stage earlier, \emph{before} it is fed into any downstream
cloner, and ask how much the restoration technique's own filtering or
perturbation distorted the signal.

These two comparisons cannot be collapsed into one because they are not
comparable pairs of recordings. STOI, PESQ, and SI-SDR are
sample-aligned metrics that assume both signals say the same words at
the same time (Section~\ref{sec:quality-metrics}); every downstream
synthesizer in this study uses a fixed reference script or source
clip rather than the original recording's own words (SV2TTS's and
F5-TTS's \texttt{SYNTH\_TEXT}, Seed-VC's fixed
\texttt{SEEDVC\_SOURCE\_CONTENT} -- Section~\ref{sec:synth}), so a
final clone never says the same thing as the clean\_bonafide original
it is scored against. Computing STOI/PESQ/SI-SDR between clean\_bonafide
and a final clone would therefore not measure signal quality at all --
it would return a near-floor score driven entirely by the linguistic
content mismatch, regardless of how faithfully the clone reproduced the
speaker's voice. The restored audio, by contrast, is still the
\emph{same} recording the restoration technique filtered or perturbed
in place; content is preserved, so the comparison is well-defined. This
mirrors exactly how imperceptibility itself is measured elsewhere in
this paper (Section~\ref{sec:quality-metrics}, Table~\ref{tab:verifiers}):
STOI/PESQ/SI-SDR are always computed against a content-preserving stage
of the pipeline, never against a resynthesized one. Quality results per
restoration technique are reported alongside each technique's EER/TAR
table below where scoring is complete; this two-comparison protocol is
implemented once, centrally, in the shared scoring driver
(\texttt{analysis/restoration\_protective\_efficacy.py}) rather than
per technique.

\subsection{Low-pass filtering with gain-multiplier search}
\label{sec:restoration-lowpass-gain}

\paragraph{Rationale.} Motivated by the hypothesis that adversarial
perturbations concentrate in high-frequency content and that residual
identity cues survive in the low-frequency band, this technique
low-pass filters the cloaked signal at a per-clip adaptive cutoff (the
signal's own spectral centroid $\times$ 1.5, order-10 filter design,
matching the cutoff rule already used elsewhere in this codebase), then
multiplies the filtered signal by a gain $M$ swept over the paper's
stated range $M \in \{1.2, 1.4, 1.6, 1.8, 2.0\}$, to test whether
amplifying the retained low-frequency content increases recoverability:
\[
y_{\mathrm{LP}}(t) = x(t) * h_{\mathrm{LP}}(t), \qquad
y_{\mathrm{mult}}(t) = M \cdot y_{\mathrm{LP}}(t).
\]
We compare three filter designs (Butterworth, Chebyshev type~I with
1~dB passband ripple, and Bessel), with Butterworth as the primary
configuration reported; $M=1.0$ under the Butterworth design reproduces
a plain, unamplified low-pass filter exactly, so we include it below as
an unamplified reference point in addition to the paper's own stated
$M \in \{1.2,\ldots,2.0\}$ sweep. \textbf{Status:} restored audio and
downstream SV2TTS cloning are complete for all three cloak methods
across all 6 gain values, including the full sweep for AntiFake (not
just the unamplified $M=1.0$ reference point originally reported).
Seed-VC/F5-TTS cloning is also complete for all three cloak methods
across the full gain sweep; see \S\ref{sec:restoration-seedvc-f5tts}.

\paragraph{Results.} Table~\ref{tab:restoration-lowpass-results} reports
full verification-leakage metrics for the attack-vc-cloaked source
(all 6 gains); Table~\ref{tab:restoration-lowpass-results-pop} reports
the same for the POP-cloaked source (all 6 gains). Both are
cloned with SV2TTS and scored against the genuine clean\_bonafide
original (identical harmonized trial-construction protocol of
\S\ref{sec:threat}). As throughout this paper, a clone that verifies as
\emph{different} from the genuine speaker (higher EER / lower TAR@FAR
than the reference points below) indicates the cloak's protection
survived restoration; a clone that verifies as \emph{same} (EER/TAR@FAR
moving back toward the clean-audio ceiling) would indicate restoration
reversed it.

\begin{table}[H]
\centering
\caption{Low-pass+gain restoration, cloned with SV2TTS from
attack-vc-cloaked source audio, full metrics per gain value and
verifier.}
\label{tab:restoration-lowpass-results}
\begin{tabular}{llccccccc}
\toprule
Gain & SV & EER & TAR@1\% & TAR@0.1\% & TAR@0.01\% & minDCF & $d'$ & ROC-AUC \\
\midrule
1.0 & SB-ECAPA    & 0.3310 & 0.1452 & 0.0345 & 0.0079 & 0.9957 & 0.9054 & 0.7338 \\
1.0 & Resemblyzer & 0.1993 & 0.2893 & 0.0810 & 0.0174 & 0.9934 & 1.7044 & 0.8853 \\
1.0 & WavLM       & 0.2745 & 0.1043 & 0.0241 & 0.0051 & 0.9971 & 1.2502 & 0.8085 \\
\addlinespace
1.2 & SB-ECAPA    & 0.3262 & 0.1429 & 0.0394 & 0.0147 & 0.9975 & 0.9229 & 0.7389 \\
1.2 & Resemblyzer & 0.1914 & 0.3027 & 0.0796 & 0.0138 & 0.9936 & 1.7400 & 0.8896 \\
1.2 & WavLM       & 0.2768 & 0.1073 & 0.0269 & 0.0043 & 1.0000 & 1.2471 & 0.8079 \\
\addlinespace
1.4 & SB-ECAPA    & 0.3273 & 0.1398 & 0.0365 & 0.0059 & 0.9975 & 0.9226 & 0.7389 \\
1.4 & Resemblyzer & 0.1927 & 0.2931 & 0.0954 & 0.0204 & 0.9920 & 1.7583 & 0.8918 \\
1.4 & WavLM       & 0.2741 & 0.1006 & 0.0224 & 0.0049 & 0.9979 & 1.2705 & 0.8119 \\
\addlinespace
1.6 & SB-ECAPA    & 0.3210 & 0.1345 & 0.0341 & 0.0112 & 0.9975 & 0.9358 & 0.7434 \\
1.6 & Resemblyzer & 0.1869 & 0.2958 & 0.0794 & 0.0177 & 0.9877 & 1.7810 & 0.8947 \\
1.6 & WavLM       & 0.2719 & 0.0994 & 0.0190 & 0.0043 & 0.9994 & 1.2816 & 0.8134 \\
\addlinespace
1.8 & SB-ECAPA    & 0.3204 & 0.1316 & 0.0308 & 0.0109 & 0.9975 & 0.9390 & 0.7436 \\
1.8 & Resemblyzer & 0.1869 & 0.2857 & 0.0831 & 0.0225 & 0.9877 & 1.7980 & 0.8967 \\
1.8 & WavLM       & 0.2712 & 0.0965 & 0.0156 & 0.0032 & 0.9992 & 1.2998 & 0.8160 \\
\addlinespace
2.0 & SB-ECAPA    & 0.3202 & 0.1349 & 0.0299 & 0.0053 & 0.9979 & 0.9477 & 0.7461 \\
2.0 & Resemblyzer & 0.1874 & 0.2999 & 0.0814 & 0.0162 & 0.9953 & 1.7875 & 0.8955 \\
2.0 & WavLM       & 0.2685 & 0.0973 & 0.0219 & 0.0041 & 0.9992 & 1.2978 & 0.8155 \\
\bottomrule
\end{tabular}
\end{table}

\begin{table}[H]
\centering
\caption{Low-pass+gain restoration, cloned with SV2TTS from
POP-cloaked source audio, full metrics per gain value and verifier.}
\label{tab:restoration-lowpass-results-pop}
\begin{tabular}{llccccccc}
\toprule
Gain & SV & EER & TAR@1\% & TAR@0.1\% & TAR@0.01\% & minDCF & $d'$ & ROC-AUC \\
\midrule
1.0 & SB-ECAPA    & 0.2978 & 0.1834 & 0.0584 & 0.0120 & 0.9994 & 1.0575 & 0.7709 \\
1.0 & Resemblyzer & 0.1662 & 0.3940 & 0.1057 & 0.0331 & 0.9825 & 1.9366 & 0.9115 \\
1.0 & WavLM       & 0.2668 & 0.0971 & 0.0214 & 0.0036 & 0.9993 & 1.2828 & 0.8148 \\
\addlinespace
1.2 & SB-ECAPA    & 0.2966 & 0.1819 & 0.0549 & 0.0111 & 0.9934 & 1.0750 & 0.7748 \\
1.2 & Resemblyzer & 0.1618 & 0.4183 & 0.1296 & 0.0324 & 0.9790 & 1.9838 & 0.9163 \\
1.2 & WavLM       & 0.2564 & 0.1235 & 0.0347 & 0.0088 & 0.9965 & 1.3447 & 0.8278 \\
\addlinespace
1.4 & SB-ECAPA    & 0.2949 & 0.1826 & 0.0542 & 0.0195 & 0.9940 & 1.0876 & 0.7774 \\
1.4 & Resemblyzer & 0.1633 & 0.4102 & 0.1212 & 0.0459 & 0.9718 & 1.9820 & 0.9168 \\
1.4 & WavLM       & 0.2558 & 0.1252 & 0.0318 & 0.0087 & 0.9979 & 1.3504 & 0.8291 \\
\addlinespace
1.6 & SB-ECAPA    & 0.2893 & 0.1794 & 0.0491 & 0.0072 & 0.9963 & 1.0946 & 0.7802 \\
1.6 & Resemblyzer & 0.1610 & 0.4132 & 0.1442 & 0.0269 & 0.9764 & 1.9925 & 0.9175 \\
1.6 & WavLM       & 0.2521 & 0.1223 & 0.0300 & 0.0056 & 0.9973 & 1.3781 & 0.8327 \\
\addlinespace
1.8 & SB-ECAPA    & 0.2905 & 0.1750 & 0.0514 & 0.0148 & 0.9951 & 1.1073 & 0.7825 \\
1.8 & Resemblyzer & 0.1596 & 0.3934 & 0.1129 & 0.0197 & 0.9866 & 1.9937 & 0.9178 \\
1.8 & WavLM       & 0.2528 & 0.1162 & 0.0275 & 0.0057 & 0.9967 & 1.3804 & 0.8330 \\
\addlinespace
2.0 & SB-ECAPA    & 0.2914 & 0.1816 & 0.0480 & 0.0083 & 0.9961 & 1.0905 & 0.7783 \\
2.0 & Resemblyzer & 0.1623 & 0.3878 & 0.1345 & 0.0392 & 0.9732 & 1.9603 & 0.9152 \\
2.0 & WavLM       & 0.2535 & 0.1161 & 0.0285 & 0.0081 & 0.9977 & 1.3658 & 0.8307 \\
\bottomrule
\end{tabular}
\end{table}

\begin{table}[H]
\centering
\caption{Low-pass restoration (unamplified, $M=1.0$ only), cloned with
SV2TTS from AntiFake-cloaked source audio, full metrics per verifier.}
\label{tab:restoration-lowpass-antifake-results}
\begin{tabular}{lccccccc}
\toprule
SV & EER & TAR@1\% & TAR@0.1\% & TAR@0.01\% & minDCF & $d'$ & ROC-AUC \\
\midrule
SB-ECAPA    & 0.5115 & 0.0367 & 0.0065 & 0.0012 & 0.9994 & 0.0209 & 0.4909 \\
Resemblyzer & 0.4038 & 0.0586 & 0.0134 & 0.0023 & 1.0000 & 0.5300 & 0.6414 \\
WavLM       & 0.4108 & 0.0198 & 0.0028 & 0.0008 & 0.9996 & 0.5123 & 0.6333 \\
\bottomrule
\end{tabular}
\end{table}

Neither cloak method's gain-multiplier sweep reverses the cloak, and
the gain itself barely matters for either. Against attack-vc-cloaked
source, EER sits in a narrow band across all 6 gains -- 0.320--0.331
(SB-ECAPA), 0.187--0.199 (Resemblyzer), and 0.269--0.277 (WavLM) --
substantially \emph{above} both the clean-audio ceiling (SV2TTS row of
Table~\ref{tab:synth-ceiling}: 0.1415/0.0547/0.1667) and the
no-restoration attack-vc-cloaked baseline (SV2TTS row of
Table~\ref{tab:protective-efficacy-attackvc}: 0.2112/0.1064/0.1914).
Against POP-cloaked source, the pattern is identical in kind though
the absolute EER is somewhat lower: 0.289--0.298 (SB-ECAPA, now
confirmed across the full 6-gain sweep), 0.160--0.166 (Resemblyzer),
0.252--0.267 (WavLM), again clearly above
both the clean-audio ceiling and the no-restoration POP-cloaked
baseline (SV2TTS row of Table~\ref{tab:protective-efficacy}:
0.2065/0.1018/0.1918). Since $M=2.0$ (attack-vc) is not meaningfully
better or worse for the attacker than the unamplified $M=1.0$ baseline
under either cloak method, amplifying the retained low-frequency
content does not help recover identity, contradicting this technique's
motivating hypothesis that residual identity cues concentrate in the
low-frequency band and simply need boosting. The restoration technique
is itself a signal-degrading filter, so applying it to already-cloaked
audio appears to compound distortion rather than remove it, for both
cloak methods symmetrically -- on this evidence, an adaptive attacker
using low-pass+gain restoration would be actively \emph{worse} off than
one who simply cloned the cloaked audio directly, regardless of which
of the two cloaks produced the source recording.

Against \textbf{AntiFake-cloaked} source, only the unamplified $M=1.0$
reference point was run (Table~\ref{tab:restoration-lowpass-antifake-results}),
but it already shows the pattern first documented in
\S\ref{sec:restoration-spectral}: EER moves \emph{down} from the
no-restoration AntiFake baseline (SV2TTS row of
Table~\ref{tab:protective-efficacy-antifake}: 0.5578/0.4999/0.4659) to
0.5115/0.4038/0.4108 -- a 0.0463--0.0961 absolute drop, largest for
Resemblyzer. This is a genuine but partial shift back toward the clean
audio ceiling (0.1415/0.0547/0.1667), not a reversal (TAR@1\%FAR stays
at 2.0--5.9\%, versus the ceiling's 84--98\%), consistent with the same
low-pass filter behaving differently against AntiFake's broader-band
perturbation than against POP's or attack-vc's more
perceptually-constrained one (\S\ref{sec:framework-restoration}).

\paragraph{Does descent or gender predict the outcome?}
Table~\ref{tab:restoration-lowpass-demographics} reports mean EER by
descent and gender group, averaged across the three verifiers (as in
\S\ref{sec:protective-efficacy-demographics}), for each gain value and
cloak method.

\begin{table}[H]
\centering
\caption{Low-pass+gain restoration: mean EER by descent/gender group
and gain value, averaged across the three verifiers, attack-vc-cloaked
vs.\ POP-cloaked source (all 6 gains, both cloaks), plus AntiFake-cloaked
source ($M=1.0$ only).}
\label{tab:restoration-lowpass-demographics}
\begin{tabular}{lcccccc}
\toprule
& \multicolumn{6}{c}{attack-vc-cloaked source} \\
Group & 1.0 & 1.2 & 1.4 & 1.6 & 1.8 & 2.0 \\
\midrule
African (men)               & 0.2629 & 0.2585 & 0.2599 & 0.2584 & 0.2625 & 0.2472 \\
African (women)             & 0.3089 & 0.3211 & 0.3106 & 0.3010 & 0.3054 & 0.3132 \\
Asian, East (men)           & 0.2089 & 0.2077 & 0.2067 & 0.2112 & 0.2096 & 0.2170 \\
Asian, East (women)         & 0.3338 & 0.3278 & 0.3187 & 0.3237 & 0.3109 & 0.3069 \\
Asian, South (men)          & 0.2898 & 0.2899 & 0.2948 & 0.2932 & 0.2909 & 0.2915 \\
Asian, South (women)        & 0.3369 & 0.3301 & 0.3267 & 0.3162 & 0.3182 & 0.3143 \\
Caucasian, American (men)   & 0.2237 & 0.2303 & 0.2118 & 0.2082 & 0.2081 & 0.2060 \\
Caucasian, American (women) & 0.3337 & 0.3097 & 0.3185 & 0.3104 & 0.3057 & 0.3010 \\
Caucasian, European (men)   & 0.2564 & 0.2539 & 0.2551 & 0.2471 & 0.2417 & 0.2393 \\
Caucasian, European (women) & 0.2875 & 0.2908 & 0.2819 & 0.2788 & 0.2804 & 0.2709 \\
\midrule
& \multicolumn{6}{c}{POP-cloaked source} \\
Group & 1.0 & 1.2 & 1.4 & 1.6 & 1.8 & 2.0 \\
\midrule
African (men)               & 0.2419 & 0.2370 & 0.2283 & 0.2238 & 0.2307 & 0.2247 \\
African (women)              & 0.2908 & 0.2899 & 0.2757 & 0.2769 & 0.2703 & 0.2736 \\
Asian, East (men)            & 0.1986 & 0.1937 & 0.1930 & 0.1997 & 0.1935 & 0.2016 \\
Asian, East (women)          & 0.3117 & 0.3088 & 0.3014 & 0.2986 & 0.2898 & 0.2978 \\
Asian, South (men)           & 0.2313 & 0.2293 & 0.2326 & 0.2385 & 0.2378 & 0.2364 \\
Asian, South (women)         & 0.3195 & 0.3155 & 0.3111 & 0.3058 & 0.3071 & 0.3159 \\
Caucasian, American (men)    & 0.2034 & 0.1929 & 0.2023 & 0.1948 & 0.1889 & 0.1867 \\
Caucasian, American (women)  & 0.2982 & 0.2862 & 0.2845 & 0.2695 & 0.2728 & 0.2680 \\
Caucasian, European (men)    & 0.2093 & 0.2019 & 0.2165 & 0.2069 & 0.2019 & 0.2066 \\
Caucasian, European (women)  & 0.2514 & 0.2408 & 0.2461 & 0.2345 & 0.2264 & 0.2321 \\
\midrule
& \multicolumn{6}{c}{AntiFake-cloaked source ($M=1.0$ only)} \\
Group & \multicolumn{6}{c}{Mean EER} \\
\midrule
African (men)               & \multicolumn{6}{c}{0.4490} \\
African (women)              & \multicolumn{6}{c}{0.4761} \\
Asian, East (men)            & \multicolumn{6}{c}{0.4197} \\
Asian, East (women)          & \multicolumn{6}{c}{0.4987} \\
Asian, South (men)           & \multicolumn{6}{c}{0.4482} \\
Asian, South (women)         & \multicolumn{6}{c}{0.4530} \\
Caucasian, American (men)    & \multicolumn{6}{c}{0.4578} \\
Caucasian, American (women)  & \multicolumn{6}{c}{0.5198} \\
Caucasian, European (men)    & \multicolumn{6}{c}{0.4563} \\
Caucasian, European (women)  & \multicolumn{6}{c}{0.5070} \\
\bottomrule
\end{tabular}
\end{table}

The gender gap already documented for the no-restoration case
(\S\ref{sec:protective-efficacy-demographics}) survives restoration
essentially unchanged in direction \emph{for both cloak methods}:
averaged over all ten groups, women have a higher mean EER than men at
every single gain value under both attack-vc-cloaked source
(0.3201~vs.\ 0.2483 at $M=1.0$) and POP-cloaked source (0.2943~vs.\
0.2169 at $M=1.0$, narrowing slightly to 0.2771~vs.\ 0.2127 at
$M=1.6$). Since higher EER here means \emph{harder to re-identify},
women's clones remain harder for an adaptive attacker to recover
regardless of which cloak produced the source audio. The two cloak
methods agree on the highest-EER (least-recoverable) group --
Asian(South)-women, at every gain value under both POP and attack-vc --
but diverge slightly on the lowest: attack-vc's low-pass+gain
consistently favors Caucasian(American)-men as most-recoverable, while
POP's alternates between Asian(East)-men and Caucasian(American)-men
depending on the gain value, both consistently among the two
lowest-EER groups either way.

The same direction holds under AntiFake-cloaked source despite the
partial reversal: women's mean EER (0.4909) exceeds men's (0.4462), and
the highest/lowest-EER groups (Caucasian(American)-women highest at
0.5198, Asian(East)-men lowest at 0.4197) match the ordering seen under
POP and attack-vc -- the partial reversal documented above lowers every
group's EER by a broadly similar amount rather than disproportionately
helping any one demographic recover their identity.

\subsection{Adaptive frequency filtering (centroid-based)}
\label{sec:restoration-adaptive-filter}

\paragraph{Rationale.} Rather than a fixed cutoff, this technique
derives both a high-pass and a low-pass cutoff from the same per-clip
spectral centroid $C$, band-passing the signal around it:
\[
f_{\mathrm{HP}} = \alpha \cdot C, \qquad f_{\mathrm{LP}} = \beta \cdot C,
\]
under the hypothesis that cloaking redistributes energy across
frequency bands and that restoring the spectral balance around the
centroid may recover identity cues. The paper excerpt available to us
does not specify $\alpha$/$\beta$; we use $\alpha=0.5$, $\beta=1.5$
(order-10 Butterworth band-pass by default, Chebyshev/Bessel also
implemented), so the retained passband straddles the centroid from
half of it up to 1.5$\times$ it -- a default chosen for consistency
with the 1.5$\times$ multiplier already used in
\S\ref{sec:restoration-lowpass-gain}, not a value taken from the paper
itself. \textbf{Status:} restored audio and downstream SV2TTS cloning
are complete for all three cloak methods, including AntiFake (single
$\alpha$/$\beta$ configuration, no sweep); Seed-VC/F5-TTS cloning is
also complete for all three, see \S\ref{sec:restoration-seedvc-f5tts}.

\paragraph{Results.} Table~\ref{tab:restoration-adaptive-results}
reports full verification-leakage metrics for POP and attack-vc;
Table~\ref{tab:restoration-adaptive-antifake-results} reports the same
for AntiFake. All are cloned with SV2TTS and scored against the genuine
clean\_bonafide original, following the same protocol as
\S\ref{sec:restoration-lowpass-gain}.

\begin{table}[H]
\centering
\caption{Adaptive-filter-centroid restoration, cloned with SV2TTS,
full metrics per verifier, POP-cloaked vs.\ attack-vc-cloaked source.}
\label{tab:restoration-adaptive-results}
\begin{tabular}{llccccccc}
\toprule
Source & SV & EER & TAR@1\% & TAR@0.1\% & TAR@0.01\% & minDCF & $d'$ & ROC-AUC \\
\midrule
POP        & SB-ECAPA    & 0.4937 & 0.0151 & 0.0019 & 0.0000 & 1.0000 & 0.0470 & 0.5112 \\
POP        & Resemblyzer & 0.4568 & 0.0113 & 0.0004 & 0.0002 & 1.0000 & 0.1691 & 0.5528 \\
POP        & WavLM       & 0.4455 & 0.0110 & 0.0010 & 0.0000 & 1.0000 & 0.2531 & 0.5712 \\
\addlinespace
attack-vc  & SB-ECAPA    & 0.5030 & 0.0160 & 0.0010 & 0.0000 & 1.0000 & 0.0253 & 0.5019 \\
attack-vc  & Resemblyzer & 0.4360 & 0.0253 & 0.0041 & 0.0014 & 0.9988 & 0.3039 & 0.5863 \\
attack-vc  & WavLM       & 0.4415 & 0.0141 & 0.0019 & 0.0000 & 1.0000 & 0.2864 & 0.5807 \\
\bottomrule
\end{tabular}
\end{table}

\begin{table}[H]
\centering
\caption{Adaptive-filter-centroid restoration, cloned with SV2TTS from
AntiFake-cloaked source audio, full metrics per verifier.}
\label{tab:restoration-adaptive-antifake-results}
\begin{tabular}{lccccccc}
\toprule
SV & EER & TAR@1\% & TAR@0.1\% & TAR@0.01\% & minDCF & $d'$ & ROC-AUC \\
\midrule
SB-ECAPA    & 0.5066 & 0.0103 & 0.0008 & 0.0000 & 1.0000 & $-$0.0194 & 0.4917 \\
Resemblyzer & 0.4913 & 0.0093 & 0.0000 & 0.0000 & 1.0000 & 0.0251 & 0.5098 \\
WavLM       & 0.4823 & 0.0100 & 0.0014 & 0.0004 & 0.9998 & 0.0750 & 0.5200 \\
\bottomrule
\end{tabular}
\end{table}

This technique's effect is far more dramatic than low-pass+gain's for
\emph{both} cloak methods: EER of 0.446--0.494 against POP-cloaked
source and 0.436--0.503 against attack-vc-cloaked source, roughly
\emph{three times} the clean-audio ceiling (SV2TTS row of
Table~\ref{tab:synth-ceiling}: 0.1415/0.0547/0.1667) and well above
either no-restoration baseline (POP: 0.2065/0.1018/0.1918; attack-vc:
0.2112/0.1064/0.1914), with TAR@1\%FAR collapsing to 1.1--2.5\% and
$d'$ near zero for SB-ECAPA under both cloaks (0.047 for POP, 0.025 for
attack-vc -- essentially no separation left between genuine and
impostor score distributions at all). Like low-pass+gain, this
restoration technique is itself a band-pass filter -- a
signal-degrading operation -- so applying it to already-cloaked audio
compounds distortion rather than restoring the spectral balance its
motivating hypothesis predicted would recover identity cues; this
remains, for both cloak methods, the clearest evidence in the study
that a restoration attempt can make an adaptive attacker's outcome
\emph{worse}, not better, than doing nothing.

Against \textbf{AntiFake-cloaked} source, the effect is smaller and
mixed in sign rather than uniformly compounding: EER moves to
0.5066/0.4913/0.4823 from the no-restoration baseline
(0.5578/0.4999/0.4659) -- SB-ECAPA and Resemblyzer drop slightly
(0.0512 and 0.0086 respectively) while WavLM rises slightly (0.0164).
Unlike spectral subtraction and simulated re-recording
(\S\ref{sec:restoration-spectral}, \S\ref{sec:restoration-rerecording}),
none of these shifts is large or consistent enough across verifiers to
count as a reversal by our criterion; adaptive-filter-centroid remains,
under all three cloak methods, a restoration technique that does not
help an adaptive attacker.

\paragraph{Does descent or gender predict the outcome?}
Table~\ref{tab:restoration-adaptive-demographics} reports mean EER by
descent and gender group, averaged across the three verifiers, for
all three cloak methods, restricted to the FakeAVCeleb subset (as in
\S\ref{sec:ablation}).

\begin{table}[H]
\centering
\caption{Adaptive-filter-centroid restoration: mean EER by descent
$\times$ gender (FakeAVCeleb subset), averaged across the three
verifiers, POP-cloaked vs.\ attack-vc-cloaked vs.\ AntiFake-cloaked
source.}
\label{tab:restoration-adaptive-demographics}
\begin{tabular}{lccc}
\toprule
Group & POP-cloaked & attack-vc-cloaked & AntiFake-cloaked \\
\midrule
African (men)               & 0.4797 & 0.5221 & 0.4864 \\
African (women)              & 0.4623 & 0.4385 & 0.5258 \\
Asian, East (men)            & 0.4310 & 0.4408 & 0.4118 \\
Asian, East (women)          & 0.4568 & 0.4167 & 0.5593 \\
Asian, South (men)           & 0.5033 & 0.5000 & 0.4860 \\
Asian, South (women)         & 0.4779 & 0.4660 & 0.5130 \\
Caucasian, American (men)    & 0.4859 & 0.5415 & 0.4754 \\
Caucasian, American (women)  & 0.5129 & 0.4445 & 0.5954 \\
Caucasian, European (men)    & 0.4333 & 0.4843 & 0.4527 \\
Caucasian, European (women)  & 0.4659 & 0.4377 & 0.5354 \\
\bottomrule
\end{tabular}
\end{table}

Averaged across the three verifiers, this is the only technique in this
study where the gender direction \emph{inverts} between the two cloak
methods, and we report it plainly rather than smooth it into the
pattern found everywhere else. Under \textbf{POP-cloaked} source, women
have the usual higher mean EER (0.4752~vs.\ 0.4666 for men), matching
every other result in this paper. Under \textbf{attack-vc-cloaked}
source, the direction flips: \emph{men} have the higher mean EER
(0.4977~vs.\ 0.4407 for women), the largest gender gap of any technique
in this section and the only case where men are the less-recoverable
group. The highest-EER cell also differs by cloak method --
Caucasian(American)-women under POP (0.513) vs.\
Caucasian(American)-\emph{men} under attack-vc (0.542, the single
highest cell in this entire table) -- while Asian(East) remains close
to the lowest-EER descent group under both (men 0.431 under POP,
women 0.417 under attack-vc). We do not have a confirmed mechanism for
why adaptive-filter-centroid specifically reverses the gender direction
under attack-vc when every other technique (including low-pass+gain,
built on a closely related filtering rationale) preserves it under both
cloaks; we flag it as an open question rather than force it into the
same account as the rest of this section. Under \textbf{AntiFake-cloaked}
source the direction does \emph{not} invert -- women's mean EER (0.5458)
again exceeds men's (0.4625), matching POP rather than attack-vc -- so
the attack-vc inversion documented above is specific to that
(technique, cloak-method) pair, not a general property of
adaptive-filter-centroid. The highest-EER cell under AntiFake is
Caucasian(American)-women (0.5954, the single highest cell in this
table), and the lowest is Asian(East)-men (0.4118), both consistent
with the ordering seen throughout this paper's AntiFake results.

\subsection{Quantization-based inversion}
\label{sec:restoration-quantization}

\paragraph{Rationale.} Reduces the cloaked signal's amplitude
resolution to 8-bit levels and immediately re-expresses it as floating
point, under the hypothesis that a cloaking perturbation relies on fine
amplitude precision that coarse quantization destroys while leaving
coarser (and, the hypothesis goes, more identity-relevant) amplitude
structure intact:
\[
q(t) = \mathrm{round}\!\left(\frac{x(t)+1}{2}\cdot(2^{B}-1)\right), \qquad
\hat{x}(t) = \frac{q(t)}{2^{B}-1}\cdot 2 - 1,
\]
with bit depth $B=8$ (256 levels), clipped to the valid range before
dequantization. This is the one technique among the ten reported here
that required no new implementation: an existing module in our
codebase already matched this exact quantize-then-dequantize procedure
byte-for-byte against the paper's description, which we confirmed with
a direct numeric trace before reusing it. \textbf{Status:} restored
audio and downstream SV2TTS cloning are complete for all three cloak
methods, including AntiFake; Seed-VC/F5-TTS cloning is also complete
for all three, see \S\ref{sec:restoration-seedvc-f5tts}.

\paragraph{Results.} Table~\ref{tab:restoration-quantization-results}
reports full verification-leakage metrics for POP and attack-vc;
Table~\ref{tab:restoration-quantization-antifake-results} reports the
same for AntiFake. All are cloned with SV2TTS and scored against the
genuine clean\_bonafide original.

\begin{table}[H]
\centering
\caption{Quantization-based inversion restoration, cloned with SV2TTS,
full metrics per verifier, POP-cloaked vs.\ attack-vc-cloaked source.}
\label{tab:restoration-quantization-results}
\begin{tabular}{llccccccc}
\toprule
Source & SV & EER & TAR@1\% & TAR@0.1\% & TAR@0.01\% & minDCF & $d'$ & ROC-AUC \\
\midrule
POP        & SB-ECAPA    & 0.2210 & 0.3174 & 0.1406 & 0.0537 & 0.9468 & 1.5135 & 0.8628 \\
POP        & Resemblyzer & 0.1205 & 0.5987 & 0.3265 & 0.1421 & 0.8355 & 2.3416 & 0.9485 \\
POP        & WavLM       & 0.2037 & 0.1970 & 0.0547 & 0.0181 & 0.9905 & 1.6054 & 0.8769 \\
\addlinespace
attack-vc  & SB-ECAPA    & 0.2299 & 0.3069 & 0.1400 & 0.0529 & 0.9566 & 1.4791 & 0.8552 \\
attack-vc  & Resemblyzer & 0.1307 & 0.5750 & 0.3169 & 0.1189 & 0.8382 & 2.2720 & 0.9432 \\
attack-vc  & WavLM       & 0.2082 & 0.1894 & 0.0549 & 0.0165 & 0.9858 & 1.5777 & 0.8720 \\
\bottomrule
\end{tabular}
\end{table}

\begin{table}[H]
\centering
\caption{Quantization-based inversion restoration, cloned with SV2TTS
from AntiFake-cloaked source audio, full metrics per verifier.}
\label{tab:restoration-quantization-antifake-results}
\begin{tabular}{lccccccc}
\toprule
SV & EER & TAR@1\% & TAR@0.1\% & TAR@0.01\% & minDCF & $d'$ & ROC-AUC \\
\midrule
SB-ECAPA    & 0.5601 & 0.0135 & 0.0021 & 0.0000 & 1.0000 & $-$0.2693 & 0.4160 \\
Resemblyzer & 0.4924 & 0.0246 & 0.0077 & 0.0019 & 0.9998 & 0.0553 & 0.5096 \\
WavLM       & 0.4659 & 0.0109 & 0.0023 & 0.0008 & 1.0000 & 0.2448 & 0.5563 \\
\bottomrule
\end{tabular}
\end{table}

Quantization does not reverse either cloak, but it is measurably milder
than low-pass+gain or adaptive-filter-centroid: EER (0.221/0.121/0.204
for POP; 0.230/0.131/0.208 for attack-vc) sits closer to, but still
clearly above, the respective no-restoration baselines
(0.2065/0.1018/0.1918 and 0.2112/0.1064/0.1914). This is consistent
with quantization's motivating hypothesis being partially, but not
fully, correct: coarsening amplitude resolution does not selectively
strip the perturbation, but it is a gentler transformation of the
signal overall than a full filtering pass, so it adds less of its own
distortion on top. As with every technique so far, both cloak methods
respond almost identically, reinforcing the theoretical account of
\S\ref{sec:framework-restoration} that a $\delta$-agnostic restoration
operator's effect is governed by its own distortion of the signal, not
by which cloak produced $\delta$.

Against \textbf{AntiFake-cloaked} source, quantization is the flattest
result of any technique tested against this cloak: EER
(0.5601/0.4924/0.4659) is essentially unchanged from the no-restoration
baseline (0.5578/0.4999/0.4659) -- SB-ECAPA rises by 0.0023,
Resemblyzer falls by 0.0075, and WavLM is identical to four decimal
places. Quantization's gentleness, which under POP/attack-vc kept it
milder than the filtering-based techniques without reversing anything,
here means it does not even produce the small partial-reversal effect
seen for the noisier low-pass and adaptive-filter-centroid results
against AntiFake -- consistent with the account that AntiFake's
broader-band perturbation needs a technique that actively attenuates
signal content (filtering, noise estimation) to be affected at all, not
merely one that coarsens amplitude resolution.

\paragraph{Does descent or gender predict the outcome?}
Table~\ref{tab:restoration-quantization-demographics} reports mean EER
by descent and gender group, averaged across the three verifiers, for
all three cloak methods.

\begin{table}[H]
\centering
\caption{Quantization-based inversion restoration: mean EER by descent
$\times$ gender (FakeAVCeleb subset), averaged across the three
verifiers, POP-cloaked vs.\ attack-vc-cloaked vs.\ AntiFake-cloaked
source.}
\label{tab:restoration-quantization-demographics}
\begin{tabular}{lccc}
\toprule
Group & POP-cloaked & attack-vc-cloaked & AntiFake-cloaked \\
\midrule
African (men)               & 0.1910 & 0.1956 & 0.5056 \\
African (women)              & 0.2040 & 0.1933 & 0.5420 \\
Asian, East (men)            & 0.1509 & 0.1501 & 0.4801 \\
Asian, East (women)          & 0.2005 & 0.2180 & 0.5534 \\
Asian, South (men)           & 0.2180 & 0.2371 & 0.5302 \\
Asian, South (women)         & 0.2541 & 0.2598 & 0.5166 \\
Caucasian, American (men)    & 0.1345 & 0.1420 & 0.5144 \\
Caucasian, American (women)  & 0.2037 & 0.2204 & 0.5735 \\
Caucasian, European (men)    & 0.1752 & 0.1706 & 0.5248 \\
Caucasian, European (women)  & 0.1923 & 0.2010 & 0.5637 \\
\bottomrule
\end{tabular}
\end{table}

The gender gap survives restoration for both cloaks (women's mean EER
exceeds men's in 9 of the 10 group pairs; the sole exception is
African/attack-vc, a near-exact tie at 0.1956~men vs.\ 0.1933~women),
and Caucasian(American)-men remains the single lowest-EER
(most-recoverable) group under both cloaks -- again the same ordering
found throughout this paper's protective-efficacy demographics. Under
\textbf{AntiFake-cloaked} source the gender gap also survives (women's
mean EER 0.5498 vs.\ men's 0.5110), but the lowest-EER descent group
shifts to Asian(East)-men (0.4801) rather than Caucasian(American)-men
-- consistent with the ordering seen elsewhere in this paper's AntiFake
results (\S\ref{sec:restoration-adaptive-filter}) rather than with the
POP/attack-vc ordering for this specific technique. As under several
other AntiFake-cloaked results in this section
(\S\ref{sec:restoration-summary}), Asian(South) is the one descent
group whose gender direction inverts here: men's mean EER (0.5302)
exceeds women's (0.5166), the opposite of the direction this same group
shows under both POP and attack-vc for quantization (and under every
other technique/cloak combination in this table).

\subsection{Ensemble averaging of perturbed variants}
\label{sec:restoration-ensemble}

\paragraph{Rationale.} Generates $N$ lightly perturbed variants of the
\emph{same} cloaked recording (a small random time shift plus
low-amplitude additive noise, independently per variant) and averages
them:
\[
Y_{\mathrm{avg}}(t) = \frac{1}{N}\sum_{i=1}^{N} Y_i(t),
\]
under the hypothesis that a cloaking perturbation is less stable under
small input perturbations than the underlying speech signal is, so
averaging several jittered copies together should attenuate the cloak
more than it attenuates genuine speech content. The paper specifies
the perturbation types (time shifts, additive noise) but not $N$ or
their magnitudes; we default to $N=5$ variants, each with an
independent random time shift of at most $\pm0.1$~ms and Gaussian noise
with standard deviation equal to 0.5\% of the signal's own RMS, with a
fixed seed for reproducibility. The shift magnitude was set empirically
rather than assumed: an initial $\pm2$~ms range destructively cancelled
broadband signal content under averaging (a directly measured
$\approx$60\% RMS reduction, consistent with phase rotation across
speech's dominant frequency content at that shift scale), which would
have made this technique behave like an aggressive implicit low-pass
filter rather than the light, stochastic perturbation the paper
describes; sweeping shift magnitudes empirically, $\pm0.1$~ms keeps the
RMS change to 5--7\%, which is the value we use throughout.
\textbf{Status:} restored audio and downstream SV2TTS cloning are
complete for all three cloak methods, including AntiFake; Seed-VC/F5-TTS
cloning is also complete for all three, see
\S\ref{sec:restoration-seedvc-f5tts}.

\paragraph{Results.} Table~\ref{tab:restoration-ensemble-results}
reports full verification-leakage metrics for POP and attack-vc;
Table~\ref{tab:restoration-ensemble-antifake-results} reports the same
for AntiFake.

\begin{table}[H]
\centering
\caption{Ensemble-averaging-of-perturbed-variants restoration, cloned
with SV2TTS, full metrics per verifier, POP-cloaked vs.\
attack-vc-cloaked source.}
\label{tab:restoration-ensemble-results}
\begin{tabular}{llccccccc}
\toprule
Source & SV & EER & TAR@1\% & TAR@0.1\% & TAR@0.01\% & minDCF & $d'$ & ROC-AUC \\
\midrule
POP        & SB-ECAPA    & 0.2345 & 0.3100 & 0.1388 & 0.0731 & 0.9435 & 1.4264 & 0.8443 \\
POP        & Resemblyzer & 0.1200 & 0.5923 & 0.3078 & 0.1102 & 0.8479 & 2.3487 & 0.9489 \\
POP        & WavLM       & 0.2087 & 0.1990 & 0.0664 & 0.0255 & 0.9893 & 1.5738 & 0.8710 \\
\addlinespace
attack-vc  & SB-ECAPA    & 0.2373 & 0.3070 & 0.1279 & 0.0440 & 0.9689 & 1.3987 & 0.8388 \\
attack-vc  & Resemblyzer & 0.1240 & 0.5649 & 0.3033 & 0.1357 & 0.8207 & 2.2734 & 0.9445 \\
attack-vc  & WavLM       & 0.2129 & 0.1832 & 0.0534 & 0.0142 & 0.9909 & 1.5660 & 0.8699 \\
\bottomrule
\end{tabular}
\end{table}

\begin{table}[H]
\centering
\caption{Ensemble-averaging-of-perturbed-variants restoration, cloned
with SV2TTS from AntiFake-cloaked source audio, full metrics per
verifier.}
\label{tab:restoration-ensemble-antifake-results}
\begin{tabular}{lccccccc}
\toprule
SV & EER & TAR@1\% & TAR@0.1\% & TAR@0.01\% & minDCF & $d'$ & ROC-AUC \\
\midrule
SB-ECAPA    & 0.5316 & 0.0225 & 0.0041 & 0.0004 & 1.0000 & $-$0.1473 & 0.4489 \\
Resemblyzer & 0.4366 & 0.0388 & 0.0112 & 0.0037 & 1.0000 & 0.3430 & 0.5909 \\
WavLM       & 0.4325 & 0.0101 & 0.0019 & 0.0008 & 0.9996 & 0.3983 & 0.5997 \\
\bottomrule
\end{tabular}
\end{table}

Ensemble averaging performs similarly to quantization -- mild relative
to the filtering-based techniques, but still failing to reverse either
cloak: EER (0.235/0.120/0.209 for POP; 0.237/0.124/0.213 for
attack-vc) again lands above both no-restoration baselines. This
partially supports the technique's motivating hypothesis in direction
(averaging is one of the gentler restoration operators tested, and its
EER is correspondingly closer to baseline than the filtering
techniques) but not in outcome: the cloak's protective effect is not
attenuated relative to the underlying speech content, it is instead
uniformly compounded by the small amount of added noise and
time-jitter, consistent with \S\ref{sec:framework-restoration}'s
account that averaging jittered copies of an already-cloaked signal is
itself a (mild) source of distortion, not a selective filter against
$\delta$.

Against \textbf{AntiFake-cloaked} source, ensemble averaging joins the
small set of techniques that produce a measurable, if partial,
reversal: EER (0.5316/0.4366/0.4325) drops from the no-restoration
baseline (0.5578/0.4999/0.4659) by 0.0262--0.0633, consistent in
direction across all three verifiers. The effect is smaller than
spectral subtraction's or simulated re-recording's
(\S\ref{sec:restoration-spectral}, \S\ref{sec:restoration-rerecording})
but in the same direction, adding a fourth data point to the pattern
that techniques which perturb or attenuate signal content -- rather
than merely re-encoding its amplitude resolution, as quantization does
-- are the ones capable of partially reversing AntiFake specifically.

\paragraph{Does descent or gender predict the outcome?}
Table~\ref{tab:restoration-ensemble-demographics} reports mean EER by
descent and gender group, averaged across the three verifiers, for all
three cloak methods.

\begin{table}[H]
\centering
\caption{Ensemble-averaging restoration: mean EER by descent $\times$
gender (FakeAVCeleb subset), averaged across the three verifiers,
POP-cloaked vs.\ attack-vc-cloaked vs.\ AntiFake-cloaked source.}
\label{tab:restoration-ensemble-demographics}
\begin{tabular}{lccc}
\toprule
Group & POP-cloaked & attack-vc-cloaked & AntiFake-cloaked \\
\midrule
African (men)               & 0.1978 & 0.1927 & 0.4813 \\
African (women)              & 0.2145 & 0.2315 & 0.4923 \\
Asian, East (men)            & 0.1437 & 0.1499 & 0.4376 \\
Asian, East (women)          & 0.2419 & 0.2464 & 0.5075 \\
Asian, South (men)           & 0.2007 & 0.2020 & 0.4946 \\
Asian, South (women)         & 0.2615 & 0.2602 & 0.4839 \\
Caucasian, American (men)    & 0.1586 & 0.1528 & 0.4844 \\
Caucasian, American (women)  & 0.2420 & 0.2286 & 0.5141 \\
Caucasian, European (men)    & 0.1620 & 0.1637 & 0.4924 \\
Caucasian, European (women)  & 0.1997 & 0.2076 & 0.5316 \\
\bottomrule
\end{tabular}
\end{table}

The same pattern holds again: women's mean EER exceeds men's in every
one of the ten group pairs under both cloak methods, and Asian(East)-men
is consistently the lowest-EER group while Asian(South)-women is
consistently among the highest. Under \textbf{AntiFake-cloaked} source
the gender gap survives in aggregate (women's mean EER 0.5059 vs.\
men's 0.4781) and Asian(East)-men remains the single lowest-EER group
(0.4376), but the pattern is not perfectly monotone at the individual
group level: Asian(South)-men (0.4946) exceeds Asian(South)-women
(0.4839) here -- the same AntiFake-specific Asian(South) inversion that
recurs across roughly half of the restoration techniques tested against
this cloak (\S\ref{sec:restoration-summary}), not a one-off specific to
ensemble averaging.

\subsection{Resampling (downsampling/upsampling round-trip)}
\label{sec:restoration-resampling}

\paragraph{Rationale.} Reduces the cloaked signal to a lower sample
rate and resamples it back to the original rate (downsampling variant:
8~kHz round-trip; upsampling variant: 48~kHz round-trip), under the
hypothesis that a resampling filter's implicit anti-aliasing low-pass
acts as an unintentional purifier, attenuating perturbation energy
introduced above the intermediate rate's Nyquist frequency. Unlike
\S\ref{sec:restoration-lowpass-gain}'s deliberately designed filter,
here the purification (if any) is a side effect of a standard,
non-adversarially-motivated signal-processing operation. \textbf{Status:}
restored audio and downstream SV2TTS cloning are complete for both
variants under all three cloak methods, including both variants under
AntiFake; Seed-VC/F5-TTS cloning is also complete for both variants
under all three cloak methods, see \S\ref{sec:restoration-seedvc-f5tts}.

\paragraph{Results.} Table~\ref{tab:restoration-resampling-results}
reports full verification-leakage metrics for the downsampling variant,
POP and attack-vc; Table~\ref{tab:restoration-downsampling-antifake-results}
reports the same downsampling variant for AntiFake.
Table~\ref{tab:restoration-upsampling-results} reports the upsampling
variant for POP and attack-vc, and
Table~\ref{tab:restoration-upsampling-antifake-results} reports the
same upsampling variant for AntiFake.

\begin{table}[H]
\centering
\caption{Resampling (downsampling round-trip) restoration, cloned with
SV2TTS, full metrics per verifier, POP-cloaked vs.\ attack-vc-cloaked
source.}
\label{tab:restoration-resampling-results}
\begin{tabular}{llccccccc}
\toprule
Source & SV & EER & TAR@1\% & TAR@0.1\% & TAR@0.01\% & minDCF & $d'$ & ROC-AUC \\
\midrule
POP        & SB-ECAPA    & 0.2722 & 0.2490 & 0.1054 & 0.0325 & 0.9653 & 1.2376 & 0.8081 \\
POP        & Resemblyzer & 0.1371 & 0.5084 & 0.2253 & 0.1042 & 0.9003 & 2.1520 & 0.9369 \\
POP        & WavLM       & 0.2321 & 0.1618 & 0.0484 & 0.0181 & 0.9949 & 1.4948 & 0.8528 \\
\addlinespace
attack-vc  & SB-ECAPA    & 0.2756 & 0.2422 & 0.0942 & 0.0425 & 0.9735 & 1.1953 & 0.7986 \\
attack-vc  & Resemblyzer & 0.1434 & 0.4812 & 0.2298 & 0.1125 & 0.8949 & 2.0825 & 0.9309 \\
attack-vc  & WavLM       & 0.2393 & 0.1609 & 0.0469 & 0.0170 & 0.9893 & 1.4634 & 0.8484 \\
\bottomrule
\end{tabular}
\end{table}

\begin{table}[H]
\centering
\caption{Resampling (downsampling round-trip) restoration, cloned with
SV2TTS from AntiFake-cloaked source audio, full metrics per verifier.}
\label{tab:restoration-downsampling-antifake-results}
\begin{tabular}{lccccccc}
\toprule
SV & EER & TAR@1\% & TAR@0.1\% & TAR@0.01\% & minDCF & $d'$ & ROC-AUC \\
\midrule
SB-ECAPA    & 0.5254 & 0.0264 & 0.0068 & 0.0017 & 0.9992 & $-$0.0566 & 0.4707 \\
Resemblyzer & 0.4106 & 0.0510 & 0.0116 & 0.0023 & 0.9988 & 0.4673 & 0.6268 \\
WavLM       & 0.4279 & 0.0108 & 0.0031 & 0.0008 & 0.9996 & 0.4229 & 0.6081 \\
\bottomrule
\end{tabular}
\end{table}

The implicit anti-aliasing low-pass does not reverse either cloak:
downsampling EER (0.272/0.137/0.232 for POP; 0.276/0.143/0.239 for
attack-vc) is clearly above both no-restoration baselines, in the same
range as the deliberately designed low-pass+gain filter
(\S\ref{sec:restoration-lowpass-gain}) despite resampling's anti-alias
filter never being tuned with cloaking in mind at all -- reinforcing
that the failure mode is generic to \emph{any} $\delta$-agnostic
filtering operation applied to already-cloaked audio
(\S\ref{sec:framework-restoration}), not specific to a filter designed
adversarially against the cloak.

Against \textbf{AntiFake-cloaked} source, downsampling instead joins
the set of techniques that partially reverse the cloak: EER
(0.5254/0.4106/0.4279) drops from the no-restoration baseline
(0.5578/0.4999/0.4659) by 0.0324--0.0893, in the same direction on all
three verifiers, a similar magnitude to ensemble averaging
(\S\ref{sec:restoration-ensemble}). The anti-aliasing low-pass built
into an 8~kHz round-trip discards everything above 4~kHz -- a
sufficiently aggressive attenuation, evidently, to meaningfully affect
AntiFake's broader-band perturbation even though the identical
attenuation compounds rather than reverses POP's and attack-vc's more
perceptually-constrained ones.

\begin{table}[H]
\centering
\caption{Resampling (upsampling round-trip) restoration, cloned with
SV2TTS, full metrics per verifier, POP-cloaked vs.\ attack-vc-cloaked
source.}
\label{tab:restoration-upsampling-results}
\begin{tabular}{llccccccc}
\toprule
Source & SV & EER & TAR@1\% & TAR@0.1\% & TAR@0.01\% & minDCF & $d'$ & ROC-AUC \\
\midrule
POP        & SB-ECAPA    & 0.2054 & 0.3621 & 0.1788 & 0.0598 & 0.9370 & 1.5969 & 0.8757 \\
POP        & Resemblyzer & 0.1011 & 0.6755 & 0.4233 & 0.2263 & 0.7311 & 2.5381 & 0.9611 \\
POP        & WavLM       & 0.1916 & 0.2164 & 0.0635 & 0.0318 & 0.9806 & 1.6958 & 0.8917 \\
\addlinespace
attack-vc  & SB-ECAPA    & 0.2171 & 0.3520 & 0.1568 & 0.0669 & 0.9414 & 1.5451 & 0.8660 \\
attack-vc  & Resemblyzer & 0.1048 & 0.6603 & 0.4093 & 0.1575 & 0.7549 & 2.4872 & 0.9586 \\
attack-vc  & WavLM       & 0.1975 & 0.2169 & 0.0827 & 0.0328 & 0.9821 & 1.6604 & 0.8861 \\
\bottomrule
\end{tabular}
\end{table}

\begin{table}[H]
\centering
\caption{Resampling (upsampling round-trip) restoration, cloned with
SV2TTS from AntiFake-cloaked source audio, full metrics per verifier.}
\label{tab:restoration-upsampling-antifake-results}
\begin{tabular}{lccccccc}
\toprule
SV & EER & TAR@1\% & TAR@0.1\% & TAR@0.01\% & minDCF & $d'$ & ROC-AUC \\
\midrule
SB-ECAPA    & 0.5651 & 0.0144 & 0.0016 & 0.0006 & 1.0000 & $-$0.2797 & 0.4123 \\
Resemblyzer & 0.4965 & 0.0231 & 0.0070 & 0.0017 & 0.9994 & 0.0335 & 0.5036 \\
WavLM       & 0.4673 & 0.0103 & 0.0019 & 0.0010 & 1.0000 & 0.2336 & 0.5532 \\
\bottomrule
\end{tabular}
\end{table}

The upsampling variant is measurably \emph{milder} than downsampling
despite sharing the same anti-aliasing-filter rationale: EER
(0.205/0.101/0.192 for POP; 0.217/0.105/0.198 for attack-vc) sits
closer to the no-restoration baselines than downsampling's own EER
does, in the same near-neutral tier as spectral subtraction and
high-pass filtering (\S\ref{sec:restoration-spectral},
\S\ref{sec:restoration-highpass}) rather than downsampling's
more-damaging tier. This is consistent with $\varepsilon_R$ scaling
with how aggressively a resampling round-trip changes the signal:
downsampling to 8~kHz discards everything above 4~kHz before
reconstructing it, while upsampling to 48~kHz only interpolates and
decimates within the original band, a gentler transformation that adds
less of its own distortion regardless of which cloak produced the
underlying $\delta$. Against \textbf{AntiFake-cloaked} source, the same
relative pattern holds: upsampling's EER (0.5651/0.4965/0.4673) sits
above the no-restoration baseline (0.5578/0.4999/0.4659) on SB-ECAPA and
WavLM but essentially ties it on Resemblyzer, i.e.\ unlike downsampling
(\S above), upsampling does \emph{not} satisfy the
all-three-verifiers-significant reversal criterion against AntiFake --
consistent with it being the milder of the two resampling variants
under every cloak method tested.

\paragraph{Does descent or gender predict the outcome?}
Table~\ref{tab:restoration-resampling-demographics} reports mean EER by
descent and gender group, averaged across the three verifiers, for
both variants and cloak methods.

\begin{table}[H]
\centering
\caption{Resampling restoration: mean EER by descent $\times$ gender
(FakeAVCeleb subset), averaged across the three verifiers, downsampling
vs.\ upsampling, POP-cloaked vs.\ attack-vc-cloaked vs.\
AntiFake-cloaked source.}
\label{tab:restoration-resampling-demographics}
\begin{tabular}{lcccccc}
\toprule
& \multicolumn{3}{c}{Downsampling} & \multicolumn{3}{c}{Upsampling} \\
Group & POP & attack-vc & AntiFake & POP & attack-vc & AntiFake \\
\midrule
African (men)               & 0.2095 & 0.2126 & 0.4619 & 0.1712 & 0.1883 & 0.5112 \\
African (women)              & 0.2529 & 0.2570 & 0.4698 & 0.1829 & 0.1782 & 0.5464 \\
Asian, East (men)            & 0.1693 & 0.1769 & 0.4464 & 0.1310 & 0.1396 & 0.4890 \\
Asian, East (women)          & 0.2829 & 0.2924 & 0.5154 & 0.1870 & 0.2051 & 0.5454 \\
Asian, South (men)           & 0.2099 & 0.2170 & 0.4647 & 0.1889 & 0.2056 & 0.5363 \\
Asian, South (women)         & 0.2816 & 0.2816 & 0.4866 & 0.2514 & 0.2506 & 0.5197 \\
Caucasian, American (men)    & 0.1910 & 0.1922 & 0.4704 & 0.1059 & 0.1184 & 0.5248 \\
Caucasian, American (women)  & 0.2649 & 0.2714 & 0.5253 & 0.1960 & 0.2036 & 0.5628 \\
Caucasian, European (men)    & 0.1897 & 0.1959 & 0.4522 & 0.1521 & 0.1446 & 0.5276 \\
Caucasian, European (women)  & 0.2206 & 0.2323 & 0.5310 & 0.1752 & 0.1942 & 0.5712 \\
\bottomrule
\end{tabular}
\end{table}

Women's mean EER exceeds men's in every group under both variants and
both cloaks, and Asian(East)-men or Caucasian(American)-men are
consistently among the lowest-EER groups -- for downsampling,
Asian(East)-men is lowest under both cloaks (identical to the ordering
found for ensemble averaging, \S\ref{sec:restoration-ensemble}); for
the milder upsampling variant, Caucasian(American)-men is lowest under
both cloaks instead, matching the pattern set by every near-neutral
technique in this section (quantization, spectral subtraction,
high-pass filtering). Asian(South)-women is the highest-EER group under
every combination of variant and cloak method. Under
\textbf{AntiFake-cloaked} downsampling, the gender gap again survives
(women's mean EER 0.5056 vs.\ men's 0.4591) and Asian(East)-men remains
the lowest-EER group (0.4464), but the highest-EER group shifts to
Caucasian(European)-women (0.5310) rather than Asian(South)-women --
consistent with the Caucasian-group-highest ordering seen throughout
this paper's other AntiFake restoration results rather than with
POP/attack-vc's own ordering for this technique. The AntiFake-cloaked
upsampling variant repeats the same pattern: the gender gap survives
(women's mean 0.5491 vs.\ men's 0.5178), Asian(East)-men is again the
lowest-EER group (0.4890), and Caucasian(European)-women is again the
highest (0.5712) -- the AntiFake-specific Caucasian-highest ordering
recurs across both resampling variants, not just downsampling.

\subsection{Mel-spectrogram-domain inversion}
\label{sec:restoration-mel}

\paragraph{Rationale.} Projects the cloaked signal onto an 80-band mel
spectrogram (a lossy, perceptually-scaled representation adversarial
perturbations are not optimized against) and reconstructs a waveform
with Griffin-Lim phase estimation (60 iterations), under the hypothesis
that whatever component of the perturbation does not survive the mel
projection is discarded for good, leaving primarily identity-relevant
structure behind. \textbf{Status:} restored audio and downstream SV2TTS
cloning are complete for all three cloak methods, including AntiFake;
Seed-VC/F5-TTS cloning is also complete for all three, see
\S\ref{sec:restoration-seedvc-f5tts}.

\paragraph{Results.} Table~\ref{tab:restoration-mel-results} reports
full verification-leakage metrics for POP and attack-vc;
Table~\ref{tab:restoration-mel-antifake-results} reports the same for
AntiFake.

\begin{table}[H]
\centering
\caption{Mel-spectrogram-domain-inversion restoration, cloned with
SV2TTS, full metrics per verifier, POP-cloaked vs.\ attack-vc-cloaked
source.}
\label{tab:restoration-mel-results}
\begin{tabular}{llccccccc}
\toprule
Source & SV & EER & TAR@1\% & TAR@0.1\% & TAR@0.01\% & minDCF & $d'$ & ROC-AUC \\
\midrule
POP        & SB-ECAPA    & 0.2880 & 0.2499 & 0.1236 & 0.0511 & 0.9622 & 1.1613 & 0.7906 \\
POP        & Resemblyzer & 0.1686 & 0.4662 & 0.2551 & 0.0620 & 0.8797 & 1.9160 & 0.9143 \\
POP        & WavLM       & 0.2315 & 0.1501 & 0.0413 & 0.0156 & 0.9971 & 1.4888 & 0.8534 \\
\addlinespace
attack-vc  & SB-ECAPA    & 0.2910 & 0.2472 & 0.1073 & 0.0388 & 0.9770 & 1.1424 & 0.7866 \\
attack-vc  & Resemblyzer & 0.1709 & 0.4529 & 0.2477 & 0.0813 & 0.8863 & 1.8768 & 0.9105 \\
attack-vc  & WavLM       & 0.2289 & 0.1497 & 0.0417 & 0.0193 & 0.9930 & 1.5096 & 0.8567 \\
\bottomrule
\end{tabular}
\end{table}

\begin{table}[H]
\centering
\caption{Mel-spectrogram-domain-inversion restoration, cloned with
SV2TTS from AntiFake-cloaked source audio, full metrics per verifier.}
\label{tab:restoration-mel-antifake-results}
\begin{tabular}{lccccccc}
\toprule
SV & EER & TAR@1\% & TAR@0.1\% & TAR@0.01\% & minDCF & $d'$ & ROC-AUC \\
\midrule
SB-ECAPA    & 0.5487 & 0.0182 & 0.0027 & 0.0002 & 1.0000 & $-$0.1889 & 0.4372 \\
Resemblyzer & 0.4531 & 0.0339 & 0.0085 & 0.0021 & 0.9998 & 0.2732 & 0.5706 \\
WavLM       & 0.4380 & 0.0120 & 0.0019 & 0.0004 & 0.9998 & 0.3667 & 0.5904 \\
\bottomrule
\end{tabular}
\end{table}

Mel-spectrogram inversion is among the most damaging restoration
techniques tested (EER 0.288/0.169/0.232 for POP, 0.291/0.171/0.229 for
attack-vc), on a par with re-recording
(\S\ref{sec:restoration-rerecording}) and clearly worse for the
attacker than resampling, quantization, or ensemble averaging. This is
notable given the technique's motivating hypothesis -- that projecting
onto a lossy, perceptually-scaled basis should discard the
perturbation, which is not optimized against that representation --
appears backwards in practice: the mel projection and Griffin-Lim phase
reconstruction discard real identity-relevant signal structure (visible
in the sizeable STOI-analogous information loss any Griffin-Lim
reconstruction incurs) at least as readily as they discard $\delta$,
producing one of the largest $\varepsilon_R$ values (per
\S\ref{sec:framework-restoration}'s notation) of any technique in this
study.

Against \textbf{AntiFake-cloaked} source, the same technique's effect
is much smaller and does not reverse: EER (0.5487/0.4531/0.4380) drops
only slightly from the no-restoration baseline (0.5578/0.4999/0.4659),
by 0.0091--0.0468, with SB-ECAPA's shift too small to count as
consistent evidence of reversal by the criterion used throughout this
section. This is a sharp contrast with mel-spectrogram inversion's own
behavior under POP/attack-vc, where it is among the \emph{most}
damaging techniques tested -- the large $\varepsilon_R$ that compounds
POP's and attack-vc's narrower-band perturbations evidently does not
translate into either compounding or reversing AntiFake's broader-band
one to the same degree, underscoring that a technique's effect size is
not a fixed property of the technique alone but depends on the specific
cloak's own spectral profile.

\paragraph{Does descent or gender predict the outcome?}
Table~\ref{tab:restoration-mel-demographics} reports mean EER by
descent and gender group, averaged across the three verifiers, for all
three cloak methods.

\begin{table}[H]
\centering
\caption{Mel-spectrogram-inversion restoration: mean EER by descent
$\times$ gender (FakeAVCeleb subset), averaged across the three
verifiers, POP-cloaked vs.\ attack-vc-cloaked vs.\ AntiFake-cloaked
source.}
\label{tab:restoration-mel-demographics}
\begin{tabular}{lccc}
\toprule
Group & POP-cloaked & attack-vc-cloaked & AntiFake-cloaked \\
\midrule
African (men)               & 0.1922 & 0.2029 & 0.4831 \\
African (women)              & 0.2432 & 0.2439 & 0.4996 \\
Asian, East (men)            & 0.1950 & 0.1876 & 0.4513 \\
Asian, East (women)          & 0.2739 & 0.2813 & 0.5254 \\
Asian, South (men)           & 0.2643 & 0.2540 & 0.5217 \\
Asian, South (women)         & 0.2872 & 0.2954 & 0.4890 \\
Caucasian, American (men)    & 0.1721 & 0.1672 & 0.4742 \\
Caucasian, American (women)  & 0.3107 & 0.3225 & 0.5422 \\
Caucasian, European (men)    & 0.2040 & 0.1940 & 0.5020 \\
Caucasian, European (women)  & 0.2728 & 0.2637 & 0.5358 \\
\bottomrule
\end{tabular}
\end{table}

The gender gap holds in every group pair under both cloaks, and
Caucasian(American)-women is the single highest-EER cell for both
cloak methods (POP: 0.311; attack-vc: 0.323) -- the same cell that
adaptive-filter-centroid identified as least-recoverable under
POP-cloaked source specifically
(\S\ref{sec:restoration-adaptive-filter}; under attack-vc-cloaked
source, adaptive-filter-centroid instead identifies
Caucasian(American)-\emph{men} as least-recoverable, part of that
technique's gender-direction inversion), while
Caucasian(American)-men is simultaneously among the lowest-EER cells
here,
the widest within-descent-group gender split of any technique reported
in this section. Under \textbf{AntiFake-cloaked} source the gender gap
survives in aggregate (women's mean EER 0.5184 vs.\ men's 0.4865),
Caucasian(American)-women remains the highest-EER cell (0.5422), and
Asian(East)-men remains the lowest (0.4513) -- but as with several other
AntiFake-cloaked results in this section (\S\ref{sec:restoration-summary}),
Asian(South) is the one descent group whose gender direction inverts:
men's mean EER (0.5217) exceeds women's (0.4890) here.

\subsection{Simulated re-recording}
\label{sec:restoration-rerecording}

\paragraph{Rationale.} Emulates physically playing the cloaked audio
through a loudspeaker and recapturing it with a microphone: convolution
with a synthetic room impulse response (RT60~$=0.25$~s), a 4th-order
Butterworth band-pass (120--7000~Hz) simulating combined
speaker/microphone frequency coloring, and additive noise at a $-35$~dB
floor. Physical re-recording is a classic real-world technique for
attempting to wash out digital adversarial perturbations, since the
analog playback/recapture path is not modeled by any of the cloak's
own optimization. \textbf{Status:} restored audio and downstream SV2TTS
cloning are complete for all three cloak methods, including AntiFake;
Seed-VC/F5-TTS cloning is also complete for all three, see
\S\ref{sec:restoration-seedvc-f5tts}.

\paragraph{Results.} Table~\ref{tab:restoration-rerecording-results}
reports full verification-leakage metrics for POP and attack-vc;
Table~\ref{tab:restoration-rerecording-antifake-results} reports the
same for AntiFake.

\begin{table}[H]
\centering
\caption{Simulated-re-recording restoration, cloned with SV2TTS, full
metrics per verifier, POP-cloaked vs.\ attack-vc-cloaked source.}
\label{tab:restoration-rerecording-results}
\begin{tabular}{llccccccc}
\toprule
Source & SV & EER & TAR@1\% & TAR@0.1\% & TAR@0.01\% & minDCF & $d'$ & ROC-AUC \\
\midrule
POP        & SB-ECAPA    & 0.2903 & 0.1615 & 0.0448 & 0.0061 & 0.9988 & 1.0948 & 0.7802 \\
POP        & Resemblyzer & 0.1833 & 0.3242 & 0.0794 & 0.0216 & 0.9924 & 1.8361 & 0.9004 \\
POP        & WavLM       & 0.2440 & 0.0868 & 0.0178 & 0.0008 & 1.0000 & 1.4492 & 0.8401 \\
\addlinespace
attack-vc  & SB-ECAPA    & 0.2902 & 0.1583 & 0.0290 & 0.0081 & 0.9979 & 1.0815 & 0.7777 \\
attack-vc  & Resemblyzer & 0.1863 & 0.3057 & 0.0971 & 0.0249 & 0.9870 & 1.8102 & 0.8972 \\
attack-vc  & WavLM       & 0.2431 & 0.0928 & 0.0138 & 0.0027 & 0.9994 & 1.4528 & 0.8396 \\
\bottomrule
\end{tabular}
\end{table}

Simulated re-recording is, together with mel-spectrogram inversion, one
of the two most damaging restoration techniques tested for both cloak
methods (EER 0.290/0.183/0.244 for POP, 0.290/0.186/0.243 for
attack-vc -- the two cloak methods land within 0.003 of each other on
every verifier, the closest agreement of any technique in this study).
The room-impulse-response convolution, band-pass coloring, and additive
noise floor together constitute a substantially larger $\varepsilon_R$
than any purely digital filtering technique, exactly as the physical
re-recording channel's own real-world purpose (removing information,
not selectively preserving identity) would predict for POP and
attack-vc: an attacker who physically re-records POP- or
attack-vc-cloaked audio through a speaker and microphone before cloning
would be doing themselves a clear disservice.

\begin{table}[H]
\centering
\caption{Simulated-re-recording restoration, cloned with SV2TTS from
AntiFake-cloaked source audio, full metrics per verifier.}
\label{tab:restoration-rerecording-antifake-results}
\begin{tabular}{lccccccc}
\toprule
SV & EER & TAR@1\% & TAR@0.1\% & TAR@0.01\% & minDCF & $d'$ & ROC-AUC \\
\midrule
SB-ECAPA    & 0.4781 & 0.0366 & 0.0060 & 0.0004 & 1.0000 & 0.1742 & 0.5380 \\
Resemblyzer & 0.3336 & 0.0596 & 0.0128 & 0.0012 & 0.9998 & 0.8525 & 0.7276 \\
WavLM       & 0.3537 & 0.0213 & 0.0029 & 0.0000 & 1.0000 & 0.8440 & 0.7071 \\
\bottomrule
\end{tabular}
\end{table}

Against \textbf{AntiFake-cloaked} source, the same reversal-favoring
direction found for spectral subtraction
(\S\ref{sec:restoration-spectral}) appears again, and here it is the
largest of any restoration result in this entire study. Relative to
AntiFake's no-restoration baseline (SV2TTS row of
Table~\ref{tab:protective-efficacy-antifake}: 0.5578/0.4999/0.4659),
re-recording drops EER to 0.4781/0.3336/0.3537 -- an absolute drop of
0.08--0.17, roughly $2$--$3\times$ the size of spectral subtraction's
drop on the same cloak. This is still well short of a full reversal
(the clean-audio ceiling remains 0.1415/0.0547/0.1667, and TAR@1\%FAR
stays at 2.1--6.0\%, far below the ceiling's 84--98\%), but it is the
single largest movement toward the ceiling found anywhere in this
restoration study, for either cloak-agnostic technique or cloak. The
physical playback/recapture channel's substantial $\varepsilon_R$ --
identified in \S\ref{sec:restoration-rerecording}'s own POP/attack-vc
results as the reason this technique is so damaging to those two cloaks
-- appears to cut the other way against AntiFake: whatever broadband,
higher-amplitude structure AntiFake's unconstrained decoy-target
objective (\S\ref{sec:framework-cloaks}) leaves in the signal is,
apparently, more vulnerable to a lossy analog channel than either
POP's or attack-vc's more perceptually-budgeted perturbation is. We do
not have direct evidence isolating why a physical-channel simulation in
particular is the most effective of the two partially-reversing
techniques found so far, and flag it as a natural target for follow-up
(e.g.\ comparing $\varepsilon_R$ measured directly against AntiFake's
own $\delta$, per \S\ref{sec:framework-improve-restoration}'s proposed
clean-audio control channel, extended to also measure attenuation of an
actual $\delta$ rather than only clean $x$).

\paragraph{Does descent or gender predict the outcome?}
Table~\ref{tab:restoration-rerecording-demographics} reports mean EER
by descent and gender group, averaged across the three verifiers, for
all three cloak methods.

\begin{table}[H]
\centering
\caption{Simulated-re-recording restoration: mean EER by descent
$\times$ gender (FakeAVCeleb subset), averaged across the three
verifiers, POP-cloaked vs.\ attack-vc-cloaked source.}
\label{tab:restoration-rerecording-demographics}
\begin{tabular}{lccc}
\toprule
Group & POP-cloaked & attack-vc-cloaked & AntiFake-cloaked \\
\midrule
African (men)               & 0.2379 & 0.2322 & 0.3992 \\
African (women)              & 0.2344 & 0.2359 & 0.3907 \\
Asian, East (men)            & 0.1846 & 0.1805 & 0.3494 \\
Asian, East (women)          & 0.2093 & 0.2144 & 0.3803 \\
Asian, South (men)           & 0.2853 & 0.2747 & 0.4241 \\
Asian, South (women)         & 0.3096 & 0.3071 & 0.4403 \\
Caucasian, American (men)    & 0.2400 & 0.2443 & 0.4084 \\
Caucasian, American (women)  & 0.2659 & 0.2692 & 0.4271 \\
Caucasian, European (men)    & 0.2597 & 0.2593 & 0.4229 \\
Caucasian, European (women)  & 0.2325 & 0.2413 & 0.4252 \\
\bottomrule
\end{tabular}
\end{table}

Re-recording is one of several techniques in this study where the
gender gap is not universal (see \S\ref{sec:restoration-summary} for
the full cross-technique picture, including a systematic
Asian(South)-specific inversion under AntiFake that re-recording itself
does \emph{not} show), though here the inversion is
partial rather than a full flip and involves different groups.
Caucasian(European)-women show
\emph{lower} mean EER than Caucasian(European)-men under both POP and
attack-vc (0.233~vs.\ 0.260 under POP, 0.241~vs.\ 0.259 under
attack-vc), but the usual women-higher direction returns under
AntiFake (0.4252~vs.\ 0.4229, a near-tie). African-women are lower than
African-men under \emph{both} POP (0.234~vs.\ 0.238) and AntiFake
(0.3907~vs.\ 0.3992) -- the only descent group that inverts under two
of the three cloak methods -- but not under attack-vc (0.236~vs.\ 0.232,
back to the usual direction). Every other descent group, including
Asian(East), keeps the same women-higher direction under all three
cloaks. Asian(South)-women remains the single highest-EER cell under
all three cloaks (POP: 0.310; attack-vc: 0.307; AntiFake: 0.440), and
Asian(East)-men the lowest under all three, so the overall extremes of
the ranking are unchanged even though two descent groups locally
invert. We do not have a confirmed explanation for why re-recording's
physical playback/recapture channel interacts with the
Caucasian(European)/African groups' acoustic characteristics
differently than every purely digital technique in this section does,
and report it as an open question.

\subsection{Spectral subtraction}
\label{sec:restoration-spectral}

\paragraph{Rationale.} Classic blind spectral subtraction: estimates a
noise profile from the cloaked signal's own quietest 10\% of
short-time-Fourier-transform frames (1024-point FFT, 256-sample hop)
and subtracts $1.2\times$ that profile's magnitude from every frame,
floored at 5\% of the original magnitude to avoid musical-noise
artifacts, reconstructing the waveform with the original phase. Unlike
a ground-truth-reference-based denoiser, a real attacker has no access
to the clean original, so the noise profile is estimated blindly from
the cloaked signal itself -- under the hypothesis that a cloaking
perturbation behaves enough like broadband noise for this classic
denoising technique to attenuate it. \textbf{Status:} restored audio
and downstream SV2TTS cloning are complete for all three cloak methods,
including AntiFake; Seed-VC/F5-TTS cloning is also complete for all
three, see \S\ref{sec:restoration-seedvc-f5tts}.

\paragraph{Results.} Table~\ref{tab:restoration-spectral-results}
reports full verification-leakage metrics for POP and attack-vc;
Table~\ref{tab:restoration-spectral-antifake-results} reports the same
for AntiFake.

\begin{table}[H]
\centering
\caption{Spectral-subtraction restoration, cloned with SV2TTS, full
metrics per verifier, POP-cloaked vs.\ attack-vc-cloaked source.}
\label{tab:restoration-spectral-results}
\begin{tabular}{llccccccc}
\toprule
Source & SV & EER & TAR@1\% & TAR@0.1\% & TAR@0.01\% & minDCF & $d'$ & ROC-AUC \\
\midrule
POP        & SB-ECAPA    & 0.2118 & 0.3797 & 0.1797 & 0.0947 & 0.9204 & 1.5843 & 0.8693 \\
POP        & Resemblyzer & 0.1036 & 0.6799 & 0.3808 & 0.1633 & 0.7722 & 2.5307 & 0.9603 \\
POP        & WavLM       & 0.1933 & 0.2256 & 0.0826 & 0.0368 & 0.9739 & 1.7039 & 0.8907 \\
\addlinespace
attack-vc  & SB-ECAPA    & 0.2101 & 0.3674 & 0.1651 & 0.0753 & 0.9332 & 1.5778 & 0.8698 \\
attack-vc  & Resemblyzer & 0.1024 & 0.6622 & 0.4126 & 0.1684 & 0.7532 & 2.5029 & 0.9599 \\
attack-vc  & WavLM       & 0.1966 & 0.2327 & 0.0899 & 0.0413 & 0.9743 & 1.6985 & 0.8907 \\
\bottomrule
\end{tabular}
\end{table}

Spectral subtraction is the technique that comes \emph{closest} to
having no effect of any of the ten tested. Against attack-vc-cloaked
source, EER (0.210/0.102/0.197) is essentially indistinguishable from
the no-restoration baseline (0.2112/0.1064/0.1914) -- SB-ECAPA and
Resemblyzer land marginally \emph{below} the baseline (by 0.0011 and
0.0040 respectively) while WavLM lands marginally above (by 0.0052),
differences well within the bootstrap confidence intervals reported in
the underlying results tables, not a meaningful shift in either
direction. Against POP-cloaked source the same near-parity holds
(0.2118/0.1036/0.1933 vs.\ 0.2065/0.1018/0.1918, all within 0.005 of
baseline). For POP and attack-vc, this is \emph{not} evidence of reversal -- at no
point does EER move back toward the clean-audio ceiling of
Table~\ref{tab:synth-ceiling} (0.1415/0.0547/0.1667) -- but spectral
subtraction is the one technique whose own distortion of the signal
($\varepsilon_R$ in the notation of \S\ref{sec:framework-restoration})
appears small enough that it neither meaningfully compounds nor
attenuates either cloak's protective effect, unlike every
filtering-based technique in this section.

\begin{table}[H]
\centering
\caption{Spectral-subtraction restoration, cloned with SV2TTS from
AntiFake-cloaked source audio, full metrics per verifier.}
\label{tab:restoration-spectral-antifake-results}
\begin{tabular}{lccccccc}
\toprule
SV & EER & TAR@1\% & TAR@0.1\% & TAR@0.01\% & minDCF & $d'$ & ROC-AUC \\
\midrule
SB-ECAPA    & 0.5410 & 0.0242 & 0.0044 & 0.0008 & 0.9998 & $-$0.1235 & 0.4508 \\
Resemblyzer & 0.4312 & 0.0489 & 0.0180 & 0.0075 & 0.9959 & 0.3971 & 0.6046 \\
WavLM       & 0.4347 & 0.0147 & 0.0029 & 0.0015 & 0.9990 & 0.4118 & 0.6018 \\
\bottomrule
\end{tabular}
\end{table}

Against \textbf{AntiFake-cloaked} source, the picture changes, and we
report it plainly since it breaks a pattern that held for every other
(technique, cloak-method) combination in this study. AntiFake's
no-restoration baseline (SV2TTS row of
Table~\ref{tab:protective-efficacy-antifake}: 0.5578/0.4999/0.4659) is
substantially higher than the restored result here
(0.5410/0.4312/0.4347) -- EER moves \emph{down}, toward the clean-audio
ceiling, on all three verifiers, most visibly for Resemblyzer (0.4999
$\to$ 0.4312, a 0.0687 absolute drop) and WavLM (0.4659 $\to$ 0.4347).
This is \emph{not} a full reversal -- every value remains far above the
ceiling of 0.1415/0.0547/0.1667, and TAR@1\%FAR stays in the single
digits (2.4--4.9\%, versus the ceiling's 84--98\%) -- but it is the
first genuine, meaningful movement in the reversing direction found
anywhere in this restoration study. We read this as consistent with
the framework of \S\ref{sec:framework-restoration}: AntiFake's
perturbation is not calibrated to an imperceptibility budget the way
POP's and attack-vc's are (\S\ref{sec:framework-cloaks}), so it likely
occupies a broader-band, higher-amplitude region of the spectrum than
either other cloak's more perceptually-constrained $\delta$ -- exactly
the profile a noise-estimation-based denoiser is best positioned to
attenuate, even though the underlying account (a $\delta$-agnostic $R$
has no mechanism to \emph{reliably} separate $\delta$ from $x$) does
not guarantee this will happen, only that it is not ruled out.

\paragraph{Does descent or gender predict the outcome?}
Table~\ref{tab:restoration-spectral-demographics} reports mean EER by
descent and gender group, averaged across the three verifiers, for all
three cloak methods.

\begin{table}[H]
\centering
\caption{Spectral-subtraction restoration: mean EER by descent
$\times$ gender (FakeAVCeleb subset), averaged across the three
verifiers, POP-cloaked vs.\ attack-vc-cloaked source.}
\label{tab:restoration-spectral-demographics}
\begin{tabular}{lccc}
\toprule
Group & POP-cloaked & attack-vc-cloaked & AntiFake-cloaked \\
\midrule
African (men)               & 0.1707 & 0.1676 & 0.5039 \\
African (women)              & 0.1956 & 0.1925 & 0.4827 \\
Asian, East (men)            & 0.1303 & 0.1299 & 0.4505 \\
Asian, East (women)          & 0.2257 & 0.2292 & 0.4931 \\
Asian, South (men)           & 0.1745 & 0.1837 & 0.5023 \\
Asian, South (women)         & 0.2375 & 0.2322 & 0.4616 \\
Caucasian, American (men)    & 0.1204 & 0.1213 & 0.5114 \\
Caucasian, American (women)  & 0.2079 & 0.1917 & 0.5244 \\
Caucasian, European (men)    & 0.1361 & 0.1406 & 0.4967 \\
Caucasian, European (women)  & 0.1914 & 0.1872 & 0.5282 \\
\bottomrule
\end{tabular}
\end{table}

Even at this near-baseline overall effect for POP and attack-vc, the
gender gap remains intact -- women's mean EER exceeds men's in every
group pair under both cloaks -- confirming again that this disparity is
a property of the downstream SV2TTS/verifier pipeline
(\S\ref{sec:protective-efficacy-demographics}), not of any particular
cloak or restoration technique's own behavior. Under AntiFake-cloaked
source, the same direction survives \emph{in aggregate} through the
partial-reversal effect documented above (women's mean EER 0.4980 vs.\
men's 0.4930, a narrower gap than either other cloak shows here) -- the
partial reversal lowers every group's EER by a broadly similar amount
rather than disproportionately helping one gender recover their
identity, and Asian(East)-men remains the lowest-EER group under all
three cloaks. At the individual-group level, though, Asian(South) is
the one descent group whose gender direction inverts here: men's mean
EER (0.5023) exceeds women's (0.4616), the same AntiFake-specific
Asian(South) inversion that recurs across roughly half of the
restoration techniques tested against this cloak
(\S\ref{sec:restoration-summary}).

\subsection{High-pass filtering}
\label{sec:restoration-highpass}

\paragraph{Rationale.} Removes low-frequency content below 80~Hz
(order-10 Butterworth), testing the hypothesis opposite to
\S\ref{sec:restoration-lowpass-gain}: that the protective perturbation
instead concentrates \emph{below} the speech band and can be removed
without materially affecting the speech content itself, all of which
lies well above 80~Hz. \textbf{Status:} restored audio and downstream
SV2TTS cloning are complete for all three cloak methods, including
AntiFake; Seed-VC/F5-TTS cloning is also complete for all three, see
\S\ref{sec:restoration-seedvc-f5tts}.

\paragraph{Results.} Table~\ref{tab:restoration-highpass-results}
reports full verification-leakage metrics for the POP- and
attack-vc-cloaked source; Table~\ref{tab:restoration-highpass-antifake-results}
reports the same for AntiFake.

\begin{table}[H]
\centering
\caption{High-pass-filtering restoration, cloned with SV2TTS, full
metrics per verifier, POP-cloaked vs.\ attack-vc-cloaked source.}
\label{tab:restoration-highpass-results}
\begin{tabular}{llccccccc}
\toprule
Source & SV & EER & TAR@1\% & TAR@0.1\% & TAR@0.01\% & minDCF & $d'$ & ROC-AUC \\
\midrule
POP        & SB-ECAPA    & 0.2084 & 0.3690 & 0.1797 & 0.0806 & 0.9204 & 1.5845 & 0.8719 \\
POP        & Resemblyzer & 0.1006 & 0.6900 & 0.3888 & 0.1238 & 0.7778 & 2.5603 & 0.9635 \\
POP        & WavLM       & 0.1890 & 0.2283 & 0.0843 & 0.0352 & 0.9805 & 1.7030 & 0.8936 \\
\addlinespace
attack-vc  & SB-ECAPA    & 0.2139 & 0.3573 & 0.1718 & 0.0722 & 0.9416 & 1.5600 & 0.8681 \\
attack-vc  & Resemblyzer & 0.1089 & 0.6609 & 0.4220 & 0.2094 & 0.7310 & 2.4883 & 0.9586 \\
attack-vc  & WavLM       & 0.1945 & 0.2184 & 0.0763 & 0.0339 & 0.9771 & 1.6803 & 0.8897 \\
\bottomrule
\end{tabular}
\end{table}

\begin{table}[H]
\centering
\caption{High-pass-filtering restoration, cloned with SV2TTS from
AntiFake-cloaked source audio, full metrics per verifier.}
\label{tab:restoration-highpass-antifake-results}
\begin{tabular}{lccccccc}
\toprule
SV & EER & TAR@1\% & TAR@0.1\% & TAR@0.01\% & minDCF & $d'$ & ROC-AUC \\
\midrule
SB-ECAPA    & 0.5599 & 0.0128 & 0.0015 & 0.0000 & 1.0000 & $-$0.2604 & 0.4196 \\
Resemblyzer & 0.4847 & 0.0217 & 0.0077 & 0.0014 & 1.0000 & 0.0820 & 0.5184 \\
WavLM       & 0.4627 & 0.0096 & 0.0019 & 0.0000 & 1.0000 & 0.2480 & 0.5579 \\
\bottomrule
\end{tabular}
\end{table}

Like spectral subtraction, high-pass filtering lands almost exactly on
the no-restoration baseline for \textbf{both POP and attack-vc}:
POP (0.2084/0.1006/0.1890 vs.\ 0.2065/0.1018/0.1918) and attack-vc
(0.2139/0.1089/0.1945 vs.\ 0.2112/0.1064/0.1914) both show only
marginal shifts, none exceeding 0.005 in either direction -- again a
near-null result, not a reversal, and the second of only two techniques
(alongside spectral subtraction, \S\ref{sec:restoration-spectral})
whose own signal distortion is small enough to avoid clearly compounding
either cloak's protective effect. This is consistent with the
technique's motivating hypothesis in a limited sense -- removing
content below 80~Hz, where essentially no speech energy lives, is close
to a no-op on the signal itself, so it neither helps nor meaningfully
hurts an adaptive attacker, for either cloak method.

Against \textbf{AntiFake-cloaked} source, high-pass filtering is
likewise near-null: EER (0.5599/0.4847/0.4627) is close to the
no-restoration baseline (0.5578/0.4999/0.4659), with SB-ECAPA rising
slightly (+0.0021) and Resemblyzer/WavLM falling slightly (-0.0152 and
-0.0032). Unlike spectral subtraction, whose analogous near-baseline
behavior under POP/attack-vc turned into a genuine partial reversal
under AntiFake (\S\ref{sec:restoration-spectral}), high-pass filtering
stays near-null under all three cloak methods -- removing content below
80~Hz evidently does not touch AntiFake's perturbation any more than it
touches POP's or attack-vc's, unlike the broader-band attenuation
low-pass filtering, downsampling, and ensemble averaging apply.

\paragraph{Does descent or gender predict the outcome?}
Table~\ref{tab:restoration-highpass-demographics} reports mean EER by
descent and gender group, averaged across the three verifiers, for all
three cloak methods.

\begin{table}[H]
\centering
\caption{High-pass-filtering restoration: mean EER by descent $\times$
gender (FakeAVCeleb subset), averaged across the three verifiers,
POP-cloaked vs.\ attack-vc-cloaked vs.\ AntiFake-cloaked source.}
\label{tab:restoration-highpass-demographics}
\begin{tabular}{lccc}
\toprule
Group & POP-cloaked & attack-vc-cloaked & AntiFake-cloaked \\
\midrule
African (men)               & 0.1711 & 0.1746 & 0.4976 \\
African (women)              & 0.1733 & 0.1794 & 0.5393 \\
Asian, East (men)            & 0.1235 & 0.1403 & 0.4671 \\
Asian, East (women)          & 0.2070 & 0.2191 & 0.5406 \\
Asian, South (men)           & 0.1917 & 0.1913 & 0.5363 \\
Asian, South (women)         & 0.2539 & 0.2529 & 0.5197 \\
Caucasian, American (men)    & 0.1079 & 0.1139 & 0.5124 \\
Caucasian, American (women)  & 0.1861 & 0.2004 & 0.5641 \\
Caucasian, European (men)    & 0.1326 & 0.1450 & 0.5195 \\
Caucasian, European (women)  & 0.1824 & 0.1954 & 0.5516 \\
\bottomrule
\end{tabular}
\end{table}

The gender gap holds in every group pair under both POP and attack-vc,
and Caucasian(American)-men is the single lowest-EER (most-recoverable)
cell under both cloaks (0.108 POP, 0.114 attack-vc), matching
quantization and spectral subtraction
(\S\ref{sec:restoration-quantization}, \S\ref{sec:restoration-spectral}) --
even the mildest restoration techniques do not disturb which groups are
easiest or hardest for SV2TTS to reproduce. Under
\textbf{AntiFake-cloaked} source the gender gap survives in aggregate
(women's mean EER 0.5431 vs.\ men's 0.5066) and Caucasian(American)-women
is again the highest-EER cell (0.5641), but Asian(South) is the one
descent group whose gender direction flips (men 0.5363 vs.\ women
0.5197) -- part of a systematic pattern specific to Asian(South) under
AntiFake-cloaked source that recurs across roughly half of the
restoration techniques tested against this cloak
(\S\ref{sec:restoration-summary}), not a one-off.

\subsection{Second-cloak application (corrected: cloak-agnostic additive noise)}
\label{sec:restoration-secondcloak}

\paragraph{Rationale, and a correction.} An earlier version of this
technique re-applied a cloak's own attack pipeline to its own
already-cloaked output, i.e.\ $\hat{x}'' = \hat{x} + \delta_2 =
x + \delta + \delta_2$ with $\delta_2$ a fresh perturbation optimized
against the same surrogate(s) $\delta$ was. On reflection this does
not fit this section's own threat model: every other restoration
technique in \S\ref{sec:restoration} is something a real attacker could
plausibly do \emph{without} knowing which cloak, if any, protected a
given recording -- that is what ``restoration'' means in an adaptive-
attacker threat model. Re-running a specific cloak's own attack
algorithm requires knowing that exact cloak and having its surrogate
model, which a restoration-side attacker in our threat model does not
have. We correct this: the technique we report here instead has the
attacker add their own \emph{generic}, cloak-agnostic noise on top of
whatever audio they downloaded, on the simple (unfounded, but
plausible-to-try) theory that adding noise might scramble an unknown
protective perturbation. This is single-pass additive white Gaussian
noise, $\hat{x}'' = \hat{x} + n$ -- no averaging (unlike
\S\ref{sec:restoration-ensemble}'s $N=5$-variant technique, which
this one is otherwise closest to in spirit), no knowledge of, or
access to, the specific cloak's surrogate model or gradient. The
underlying implementation that reused the cloak's own attack pipeline
(\texttt{protect\_already\_cloaked.py}) is left in place, not deleted,
but is no longer what this subsection reports.

\paragraph{Noise level.} We calibrate $n$ to a fixed signal-to-noise
ratio rather than an arbitrary amplitude, since SNR is the standard,
interpretable way to specify ``how much generic noise'' in the
speech-processing literature, and default to $\mathrm{SNR}=30$~dB. No
source paper specifies an exact level for this attacker action (unlike,
e.g., POP's own stated $\varepsilon=8/255$), so this is a deliberately
chosen default, justified as follows: 30~dB sits in the middle of the
range the ITU-T P.800-family perceptual-quality literature and the
broader speech-enhancement literature commonly treat as
\emph{mild-but-real} background noise -- clearly audible and measurably
present in the signal, but well short of the 0--10~dB range associated
with severely degraded, barely intelligible speech. This keeps the
technique consistent with the same design goal every other restoration
technique in this study was held to: a plausible action a real
attacker might actually take, not a maximally destructive one chosen to
force a particular result. We do not sweep SNR in this initial report;
Table~\ref{tab:restoration-status} below records this as a single-configuration
technique, the same status every other non-swept technique in this
section has.

\paragraph{Status.} Restored audio generation (single-pass, 30~dB SNR)
and downstream SV2TTS cloning are both complete for all three cloak
methods -- POP (4,777 of 4,864, a handful of clips permanently excluded
after two retry passes on transient failures), attack-vc (4,863 of
4,863), and AntiFake (4,844 of 4,844) -- extending this technique
beyond the POP/attack-vc scope of every other restoration technique in
\S\ref{sec:restoration}, since a cloak-agnostic noise-adding attacker
is equally well-defined against AntiFake-cloaked source audio.

\paragraph{Results.} Table~\ref{tab:restoration-secondcloak-results}
reports full verification-leakage metrics for all three cloak methods,
cloned with SV2TTS and scored against the genuine clean\_bonafide
original.

\begin{table}[H]
\centering
\caption{Second-cloak-application (corrected: cloak-agnostic additive
noise) restoration, cloned with SV2TTS, full metrics per verifier, all
three cloak methods.}
\label{tab:restoration-secondcloak-results}
\begin{tabular}{llccccccc}
\toprule
Source & SV & EER & TAR@1\% & TAR@0.1\% & TAR@0.01\% & minDCF & $d'$ & ROC-AUC \\
\midrule
POP        & SB-ECAPA    & 0.2092 & 0.3634 & 0.1666 & 0.0779 & 0.9297 & 1.5803 & 0.8725 \\
POP        & Resemblyzer & 0.1080 & 0.6602 & 0.3705 & 0.1545 & 0.7984 & 2.4797 & 0.9569 \\
POP        & WavLM       & 0.1922 & 0.2122 & 0.0638 & 0.0209 & 0.9931 & 1.6684 & 0.8880 \\
\addlinespace
attack-vc  & SB-ECAPA    & 0.2132 & 0.3422 & 0.1557 & 0.0536 & 0.9543 & 1.5584 & 0.8692 \\
attack-vc  & Resemblyzer & 0.1121 & 0.6396 & 0.3627 & 0.1010 & 0.8077 & 2.4593 & 0.9563 \\
attack-vc  & WavLM       & 0.2014 & 0.2079 & 0.0657 & 0.0222 & 0.9895 & 1.6492 & 0.8840 \\
\addlinespace
AntiFake   & SB-ECAPA    & 0.5568 & 0.0161 & 0.0017 & 0.0002 & 1.0000 & $-$0.2554 & 0.4198 \\
AntiFake   & Resemblyzer & 0.4960 & 0.0281 & 0.0080 & 0.0027 & 0.9994 & 0.0521 & 0.5077 \\
AntiFake   & WavLM       & 0.4668 & 0.0083 & 0.0019 & 0.0004 & 0.9996 & 0.2460 & 0.5562 \\
\bottomrule
\end{tabular}
\end{table}

For POP and attack-vc, this near-neutral, mild single-pass noise lands
almost exactly on the respective no-restoration baselines (POP:
0.2065/0.1018/0.1918; attack-vc: 0.2112/0.1064/0.1914) -- the same
near-parity regime spectral subtraction and high-pass filtering land in
(\S\ref{sec:restoration-spectral}, \S\ref{sec:restoration-highpass}),
consistent with 30~dB SNR noise being one of the gentlest perturbations
tested in this study. For AntiFake, the result is equally neutral
relative to that cloak's own no-restoration baseline (SV2TTS row of
Table~\ref{tab:protective-efficacy-antifake}: 0.5578/0.4999/0.4659 vs.\
0.5568/0.4960/0.4668 here) -- unlike spectral subtraction and
re-recording (\S\ref{sec:restoration-spectral},
\S\ref{sec:restoration-rerecording}), this technique does \emph{not}
show the AntiFake-specific partial-reversal effect, reinforcing that
that effect is tied to a restoration operator's own distortion
magnitude ($\varepsilon_R$) rather than being a general property of
every technique applied to AntiFake-cloaked audio: 30~dB single-pass
noise is simply too mild an operator to move AntiFake's result the way
spectral subtraction's noise-floor subtraction or re-recording's
physical-channel simulation do.

\paragraph{Does descent or gender predict the outcome?}
Table~\ref{tab:restoration-secondcloak-demographics} reports mean EER
by descent and gender group, averaged across the three verifiers, for
all three cloak methods.

\begin{table}[H]
\centering
\caption{Second-cloak-application (corrected) restoration: mean EER by
descent $\times$ gender (FakeAVCeleb subset), averaged across the three
verifiers, all three cloak methods.}
\label{tab:restoration-secondcloak-demographics}
\begin{tabular}{lccc}
\toprule
Group & POP-cloaked & attack-vc-cloaked & AntiFake-cloaked \\
\midrule
African (men)               & 0.1733 & 0.1864 & 0.4999 \\
African (women)              & 0.1923 & 0.1821 & 0.5480 \\
Asian, East (men)            & 0.1287 & 0.1367 & 0.4814 \\
Asian, East (women)          & 0.1967 & 0.2078 & 0.5463 \\
Asian, South (men)           & 0.1949 & 0.2151 & 0.5296 \\
Asian, South (women)         & 0.2571 & 0.2497 & 0.5151 \\
Caucasian, American (men)    & 0.1100 & 0.1166 & 0.5160 \\
Caucasian, American (women)  & 0.2066 & 0.2094 & 0.5669 \\
Caucasian, European (men)    & 0.1461 & 0.1497 & 0.5225 \\
Caucasian, European (women)  & 0.1829 & 0.1879 & 0.5727 \\
\bottomrule
\end{tabular}
\end{table}

The usual women-higher gender direction holds for POP and attack-vc,
and for AntiFake in 4 of 5 descent groups -- the exception is
Asian(South), where men's mean EER (0.5296) exceeds women's (0.5151),
part of the same systematic Asian(South)-under-AntiFake inversion
documented across roughly half of the restoration techniques tested
against this cloak (\S\ref{sec:restoration-summary}), not specific to
second-cloak-noise. Caucasian(American)-men is the lowest-EER group under
both POP and attack-vc, consistent with every mild technique in this
section, while under AntiFake, Asian(East)-men is lowest (0.4814) --
the same descent/gender combination that is lowest for AntiFake under
spectral subtraction and re-recording as well
(\S\ref{sec:restoration-spectral}, \S\ref{sec:restoration-rerecording}),
suggesting this group's relative recoverability under AntiFake is a
property of the cloak/verifier pipeline rather than of any specific
restoration technique.

\begin{table}[H]
\centering
\caption{Restoration-technique completion status at time of writing.
SV2TTS-clone completion reported separately for POP, attack-vc, and
AntiFake source (P / A / AF); ``---'' means that (technique,
cloak-method) combination was not attempted.}
\label{tab:restoration-status}
\begin{tabular}{lcc}
\toprule
Technique & Variants & SV2TTS clone (P / A / AF) \\
\midrule
Low-pass + gain search         & 6 gains $\times$ 3 methods & Complete$^\ddagger$ / Complete / Complete \\
Adaptive frequency filtering   & 1 config $\times$ 3 methods & Complete / Complete / Complete \\
Quantization-based inversion   & 1 config $\times$ 3 methods & Complete / Complete / Complete \\
Ensemble averaging (perturbed) & 1 config $\times$ 3 methods & Complete / Complete / Complete \\
Resampling (downsampling)      & 1 config $\times$ 3 methods & Complete / Complete / Complete \\
Resampling (upsampling)        & 1 config $\times$ 3 methods & Complete$^\ddagger$ / Complete$^\ddagger$ / Complete \\
Mel-spectrogram-domain inversion & 1 config $\times$ 3 methods & Complete / Complete / Complete \\
Simulated re-recording         & 1 config $\times$ 3 methods & Complete / Complete / Complete \\
Spectral subtraction           & 1 config $\times$ 3 methods & Complete / Complete / Complete \\
High-pass filtering            & 1 config $\times$ 3 methods & Complete / Complete$^\ddagger$ / Complete \\
Second-cloak application$^\dagger$ & 1 config $\times$ 3 methods & Complete / Complete / Complete \\
\bottomrule
\end{tabular}

\vspace{0.2em}
\footnotesize $^\dagger$ Corrected to cloak-agnostic additive noise
(\S\ref{sec:restoration-secondcloak}). $^\ddagger$ A small number of
clips (0.3--3\% of the relevant corpus) remained permanently
unattainable after two retry passes distinguishing transient GPU
contention from genuine per-clip failures (\S\ref{sec:cloaking}'s
engineering-obstacles discussion applies identically here); scored
against the achieved corpus size rather than excluded from this table.
\end{table}

\subsection{Cross-technique summary: does any restoration technique reverse either cloak?}
\label{sec:restoration-summary}

Table~\ref{tab:restoration-crosstechnique} summarizes SB-ECAPA EER
across all ten restoration techniques, now fully scored for all three
cloak methods (POP, attack-vc, and AntiFake), alongside the clean-audio
ceiling and each cloak's no-restoration baseline, ranked (for
POP/attack-vc) from most to least damaging to an adaptive attacker.

\begin{table}[H]
\centering
\caption{SB-ECAPA EER across all ten restoration techniques, ranked
(POP/attack-vc columns) from most to least damaging to an adaptive
attacker (SV2TTS clone vs.\ genuine). All three columns are now
complete; \textbf{bold} marks cells where EER drops \emph{below} the
no-restoration baseline on all three verifiers with non-overlapping
95\% bootstrap confidence intervals (statistically significant partial
reversal in direction, not merely numerically lower, and not a return
to the clean-audio ceiling -- see \S\ref{sec:restoration-seedvc-f5tts}).}
\label{tab:restoration-crosstechnique}
\begin{tabular}{lccc}
\toprule
Reference / Technique & POP-cloaked & attack-vc-cloaked & AntiFake-cloaked \\
\midrule
\emph{Clean-audio ceiling (Table~\ref{tab:synth-ceiling})} & \multicolumn{3}{c}{\emph{0.1415}} \\
\emph{No-restoration cloaked baseline} & \emph{0.2065} & \emph{0.2112} & \emph{0.5578} \\
\midrule
Adaptive-filter-centroid       & 0.4937 & 0.5030 & 0.5066 \\
Simulated re-recording         & 0.2903 & 0.2902 & \textbf{0.4781} \\
Mel-spectrogram inversion      & 0.2880 & 0.2910 & 0.5487 \\
Low-pass + gain ($M$=1.0--2.0)  & 0.289--0.298 & 0.320--0.331 & \textbf{0.504--0.513} \\
Resampling (downsampling)      & 0.2722 & 0.2756 & \textbf{0.5254} \\
Ensemble averaging (perturbed) & 0.2345 & 0.2373 & \textbf{0.5316} \\
Second-cloak-application (noise) & 0.2092 & 0.2132 & 0.5568 \\
Quantization                   & 0.2210 & 0.2299 & 0.5601 \\
Spectral subtraction           & 0.2118 & 0.2101 & \textbf{0.5410} \\
High-pass filtering            & 0.2084 & 0.2139 & 0.5599 \\
\bottomrule
\end{tabular}
\end{table}

Against \textbf{POP and attack-vc}, across all 20 (technique,
cloak-method) combinations, not one moves EER back toward the
clean-audio ceiling -- every cell in the first two columns of
Table~\ref{tab:restoration-crosstechnique} sits at or above its
respective no-restoration baseline. The ranking itself is informative
and consistent with the theoretical account of
\S\ref{sec:framework-restoration}: adaptive-filter-centroid, the most
aggressive band-pass operation tested, is by far the most damaging to
an attacker (EER nearly $3.5\times$ the no-restoration baseline);
simulated re-recording and mel-spectrogram inversion, which both
discard or reconstruct substantial signal content (a physical
playback/recapture channel and a lossy Griffin-Lim phase
reconstruction respectively), are the next most damaging; the gentler,
lower-$\varepsilon_R$ operators -- quantization, ensemble averaging,
upsampling, second-cloak-noise, and especially spectral subtraction and
high-pass filtering -- land closest to (though never below, in the
sense of reversing) the no-restoration baseline. \textbf{Every
technique ranks in essentially the same position for both cloak
methods}, strong evidence for the framework's claim that a
$\delta$-agnostic restoration operator's effect against a
perceptually-budgeted cloak is governed by its own distortion of the
signal ($\varepsilon_R$), not by which of POP's error-minimizing
reconstruction objective or attack-vc's decoy-target embedding
objective (\S\ref{sec:framework-cloaks}) produced the perturbation
being restored.

Against \textbf{AntiFake}, the picture is different, and this is the
one place in this restoration study where the $\delta$-agnostic
account of \S\ref{sec:framework-restoration} is not simply confirmed.
All ten techniques are now scored against AntiFake-cloaked source, on
SV2TTS; we classify each cell using a statistical, not merely
numerical, criterion: a technique \textbf{reverses} the cloak if EER
drops below the no-restoration baseline on \emph{all three} verifiers
with non-overlapping 95\% bootstrap confidence intervals (i.e., a
change too large to be sampling noise, in the reversing direction,
agreed on by every verifier). By this criterion, \textbf{five of ten}
techniques reverse: low-pass filtering (confirmed across the full
$M=1.0$--$2.0$ gain sweep, not just the single point originally
reported -- every one of the 6 gain values individually satisfies this
same all-three-verifiers-significant criterion), resampling
(downsampling), ensemble averaging, spectral subtraction, and simulated
re-recording (bolded in Table~\ref{tab:restoration-crosstechnique}).
\textbf{Three} show no
significant effect in either direction (quantization, high-pass
filtering, second-cloak-noise -- all three verifiers' confidence
intervals overlap the baseline). \textbf{One}, mel-spectrogram
inversion, is borderline: Resemblyzer and WavLM drop significantly but
SB-ECAPA does not, so we do not count it as reversing by our
all-three-verifiers criterion. \textbf{One}, adaptive-filter-centroid,
is genuinely mixed: SB-ECAPA drops significantly while WavLM
\emph{rises} significantly, the two verifiers disagreeing on direction,
not just magnitude. None of the five reversing techniques approaches
the clean-audio ceiling of 0.1415 -- re-recording, the largest effect,
remains 0.34 above it, and TAR@1\%FAR stays in the single digits for
every technique tested against AntiFake, far short of the ceiling's
84--98\% -- so this is not a reversal in any practical sense, but it is
a genuine, reproducible, statistically-supported partial effect for
half the techniques tested, not an isolated pair of outliers as
initially observed with only three techniques scored.
\S\ref{sec:restoration-spectral} and \S\ref{sec:restoration-rerecording}
discuss our leading explanation for why some techniques reverse and
others do not: AntiFake's unconstrained decoy-target objective
($\lambda_{\mathrm{recon}}=0$ in the notation of
\S\ref{sec:framework-cloaks}) leaves broader-band, higher-amplitude
signal structure than either budget-constrained cloak, and the five
techniques that reverse it are exactly the ones that actively attenuate
or perturb signal content (filtering, anti-aliasing, jittered
averaging, noise-floor subtraction, physical-channel simulation) rather
than merely re-encoding it (quantization) or leaving it structurally
unchanged in the frequency ranges that matter (high-pass, mild additive
noise) -- with mel-spectrogram inversion's Griffin-Lim reconstruction
and adaptive-filter-centroid's aggressive band-pass sitting in between,
distorting the signal enough to move some verifiers but not decisively
enough (or not in a verifier-consistent direction) to count as a clean
reversal.

The demographic pattern is mostly, but not entirely, as consistent as
the POP/attack-vc result. In aggregate, women's mean EER exceeds men's
under every technique and every cloak method tested in this section,
including every one of the five AntiFake-reversing techniques -- the
partial reversal lowers every group's EER by a broadly similar amount
rather than disproportionately helping any one gender recover. But at
the level of individual descent $\times$ gender cells, AntiFake produces
a genuinely systematic exception that we did not expect and report
plainly: \textbf{Asian(South) is the one descent group whose gender
direction inverts under AntiFake-cloaked source, in six of the ten
techniques scored} (quantization, ensemble averaging, mel-spectrogram
inversion, spectral subtraction, high-pass filtering, and
second-cloak-noise -- \S\ref{sec:restoration-quantization},
\S\ref{sec:restoration-ensemble}, \S\ref{sec:restoration-mel},
\S\ref{sec:restoration-spectral}, \S\ref{sec:restoration-highpass},
\S\ref{sec:restoration-secondcloak}), with men's mean EER exceeding
women's by 1--4 percentage points in every one of those six cases. This
inversion never occurs for Asian(South) under POP or attack-vc, in any
technique, and it does not occur for AntiFake's other four descent
groups (African, Asian(East), Caucasian(American), Caucasian(European)),
which keep the usual women-higher direction under all ten AntiFake
techniques scored. The remaining four AntiFake techniques
(low-pass+gain, adaptive-filter-centroid, downsampling, simulated
re-recording) do \emph{not} show this specific inversion, though two of
them (re-recording, adaptive-filter-centroid) show different,
unrelated local inversions of their own
(\S\ref{sec:restoration-rerecording},
\S\ref{sec:restoration-adaptive-filter}). We do not have a confirmed
mechanism for why Asian(South) specifically inverts under AntiFake for
almost exactly the same six-technique subset regardless of how
different those six techniques' own signal-processing rationales are
from one another, and flag it as an open question for follow-up work
rather than force a post-hoc explanation. The pre-existing disparity
documented in \S\ref{sec:protective-efficacy-demographics} is, on this
evidence, a robust but not universal property of the downstream
cloning/verification pipeline, largely unaffected by restoration in
aggregate, with one specific, reproducible, cloak-and-descent-specific
exception.

On the evidence collected across all 30 (technique, cloak-method)
combinations now scored -- all of POP's, attack-vc's, and AntiFake's
ten techniques each, including AntiFake's upsampling variant
(\S\ref{sec:restoration-resampling}), which turned out \emph{not} to
reverse the cloak (it clears the criterion on neither Resemblyzer nor,
marginally, either other verifier) -- we can state the finding of this
section directly and without overclaiming in either direction. For POP
and attack-vc, across ten qualitatively different restoration
strategies (linear filtering in three variants, lossy-basis projection,
nonlinear quantization, nonlinear spectral masking, stochastic
averaging, physical-channel simulation, and cloak-agnostic noise
injection), a signal-domain adaptive attacker gains nothing -- several
techniques make the attacker's outcome measurably worse than doing
nothing, and the mildest (spectral subtraction, high-pass filtering,
second-cloak-noise) are indistinguishable from the no-restoration
baseline rather than an improvement on it. For AntiFake, that
conclusion holds for four of the ten techniques tested (upsampling,
quantization, high-pass filtering, second-cloak-noise show no
significant effect) and is genuinely ambiguous for two more
(mel-spectrogram inversion, adaptive-filter-centroid); for the other
five (low-pass filtering -- now confirmed across its full
$M=1.0$--$2.0$ gain sweep, not just the single point originally
reported -- downsampling, ensemble averaging, spectral subtraction, and
especially simulated re-recording), restoration produces a real,
reproducible, statistically-supported partial move toward the ceiling
-- evidence that AntiFake's protective effect, while by far the
strongest measured in this project against a naive attacker
(\S\ref{sec:protective-efficacy-antifake}), is measurably less robust to
signal-domain restoration than POP's or attack-vc's is against every
technique we tested. \S\ref{sec:restoration-seedvc-f5tts} extends all
90 (technique, cloak-method, synthesizer) combinations to Seed-VC and
F5-TTS as the downstream cloning system: none of the five
AntiFake-reversing techniques clears the same all-three-verifiers bar
on either newer synthesizer, so the partial-reversal effect identified
here for SV2TTS does not extend, at least not to the point of
statistical significance, to either of the other two cloning systems.

\subsection{Extending the restoration study to Seed-VC and F5-TTS (POP, attack-vc, and AntiFake)}
\label{sec:restoration-seedvc-f5tts}

Every restoration result reported so far in this section clones the
restored candidate with SV2TTS only. All three cloak methods now have
\emph{all three} downstream systems' cloaked-source cloning complete
(Tables~\ref{tab:protective-efficacy}, \ref{tab:protective-efficacy-attackvc},
and the AntiFake equivalent), so we extended the restoration study to
Seed-VC and F5-TTS as well, for all 10 techniques (low-pass+gain
reported as a range across its 6 gain values, as elsewhere in this
section) against POP-, attack-vc-, and AntiFake-cloaked source -- 60
additional (technique, cloak-method, synthesizer) combinations in total
-- using the same statistical reversal criterion as
\S\ref{sec:restoration-summary} (EER drops below the no-restoration
baseline on all three verifiers with non-overlapping 95\% confidence
intervals). Quality metrics (STOI/PESQ/SI-SDR) are reported for every
combination this time, computed the same way as throughout this
section: from the restored audio itself, one stage before either clone
is made (\S\ref{sec:restoration} common protocol).

\paragraph{The baseline is per-synthesizer, not shared.} This is the
one place in the restoration study where na\"ively reusing a single
"no-restoration baseline" would be actively misleading, so we state it
plainly before presenting results. Tables~\ref{tab:protective-efficacy}
and \ref{tab:protective-efficacy-attackvc} already establish that both
cloaks' no-restoration protective effect is real and moderate against
SV2TTS (SB-ECAPA EER 0.2065 for POP, 0.2112 for attack-vc) but close to
the clean-audio ceiling against Seed-VC (0.0359 / 0.0329) and F5-TTS
(0.0256 / 0.0212) -- i.e.\ neither cloak protects much against
Seed-VC/F5-TTS even for a \emph{naive} attacker who applies no
restoration at all. A restoration-technique cell's EER against Seed-VC
or F5-TTS is therefore compared against \emph{that cloak method's and
that synthesizer's own} baseline, never against SV2TTS's or the other
cloak's -- comparing across synthesizers would make ordinary
restoration EER values (e.g.\ Seed-VC's 0.03--0.08 range below) look
like a dramatic reversal when they are in fact close to where the cloak
already sat before any restoration was applied. AntiFake is the
exception to this pattern in the other direction: unlike POP and
attack-vc, its no-restoration baseline stays substantially elevated on
\emph{all three} synthesizers, not just SV2TTS (SB-ECAPA EER 0.5578 for
SV2TTS, 0.4931 for Seed-VC, 0.3817 for F5-TTS -- all far above the
near-ceiling 0.02--0.04 range POP and attack-vc fall to against
Seed-VC/F5-TTS), consistent with AntiFake's SV2TTS-only protective
effect already being the strongest of the three cloaks
(\S\ref{sec:results-antifake}). We apply the identical
per-synthesizer-baseline comparison to AntiFake regardless.

\begin{table}[H]
\centering
\caption{POP restoration study extended to Seed-VC and F5-TTS: EER
(SB-ECAPA) per technique and synthesizer (own-synthesizer baseline:
SV2TTS 0.2065, Seed-VC 0.0359, F5-TTS 0.0256), plus signal-quality
metrics for the restored audio. Low-pass+gain reports the range across
all 6 gain values ($M=1.0$--$2.0$); quality shown for $M=1.0$ only.}
\label{tab:restoration-pop-seedvc-f5tts}
\begin{tabular}{lcccccc}
\toprule
Technique & SV2TTS & Seed-VC & F5-TTS & STOI & PESQ & SI-SDR (dB) \\
\midrule
Adaptive-filter-centroid  & 0.4937 & 0.4762 & 0.4998 & 0.550 & 1.69 & $-$9.3 \\
Downsampling              & 0.2722 & 0.1369 & 0.0508 & 0.966 & 2.83 & 14.7 \\
Upsampling                & 0.2054 & 0.0352 & 0.0234 & 0.971 & 2.93 & 16.3 \\
Ensemble averaging        & 0.2345 & 0.0416 & 0.0275 & 0.970 & 3.01 & 12.5 \\
Mel-spectrogram inversion & 0.2880 & 0.0491 & 0.0245 & 0.877 & 1.84 & $-$30.5 \\
Quantization              & 0.2210 & 0.0562 & 0.0613 & 0.944 & 2.15 & 15.7 \\
Simulated re-recording    & 0.2903 & 0.1308 & 0.1361 & 0.681 & 1.44 & $-$19.8 \\
Spectral subtraction      & 0.2118 & 0.0405 & 0.0126 & 0.960 & 2.57 & 15.0 \\
High-pass filtering       & 0.2084 & 0.0360 & 0.0214 & 0.971 & 2.96 & 15.3 \\
Second-cloak-noise        & 0.2092 & 0.0419 & 0.0383 & 0.963 & 2.29 & 15.9 \\
Low-pass + gain ($M$=1.0--2.0) & 0.289--0.298 & 0.159--0.172 & 0.062--0.066 & 0.936$^{\ast}$ & 2.81$^{\ast}$ & 13.4$^{\ast}$ \\
\bottomrule
\end{tabular}
\end{table}

\begin{table}[H]
\centering
\caption{attack-vc restoration study extended to Seed-VC and F5-TTS:
EER (SB-ECAPA) per technique and synthesizer (own-synthesizer baseline:
SV2TTS 0.2112, Seed-VC 0.0329, F5-TTS 0.0212), plus signal-quality
metrics for the restored audio. Low-pass+gain reports the range across
all 6 gain values ($M=1.0$--$2.0$); quality shown for $M=1.0$ only.}
\label{tab:restoration-attackvc-seedvc-f5tts}
\begin{tabular}{lcccccc}
\toprule
Technique & SV2TTS & Seed-VC & F5-TTS & STOI & PESQ & SI-SDR (dB) \\
\midrule
Adaptive-filter-centroid  & 0.5030 & 0.4790 & 0.4966 & 0.345 & 1.84 & $-$41.0 \\
Downsampling              & 0.2756 & 0.1333 & 0.0445 & 0.544 & 3.12 & $-$33.9 \\
Upsampling                & 0.2171 & 0.0353 & 0.0224 & 0.547 & 3.58 & $-$34.0 \\
Ensemble averaging        & 0.2373 & 0.0376 & 0.0240 & 0.547 & 3.51 & $-$33.8 \\
Mel-spectrogram inversion & 0.2910 & 0.0510 & 0.0229 & 0.503 & 2.09 & $-$36.7 \\
Quantization              & 0.2299 & 0.0533 & 0.0584 & 0.535 & 2.29 & $-$34.0 \\
Simulated re-recording    & 0.2902 & 0.1163 & 0.1248 & 0.433 & 1.48 & $-$35.1 \\
Spectral subtraction      & 0.2101 & 0.0370 & 0.0099 & 0.540 & 3.13 & $-$33.8 \\
High-pass filtering       & 0.2139 & 0.0346 & 0.0202 & 0.547 & 3.53 & $-$34.0 \\
Second-cloak-noise        & 0.2132 & 0.0428 & 0.0416 & 0.546 & 2.65 & $-$33.9 \\
Low-pass + gain ($M$=1.0--2.0) & 0.320--0.331 & 0.185--0.199 & 0.080--0.083 & 0.516$^{\ast}$ & 2.92$^{\ast}$ & $-$33.8$^{\ast}$ \\
\bottomrule
\end{tabular}

{\small $^{\ast}$Quality reported for $M=1.0$ only; not re-measured per
gain value, since low-pass+gain's own signal transformation (not the
downstream synthesizer) is what quality measures, and the gain
multiplier's effect on that is already characterized in
\S\ref{sec:restoration-lowpass-gain}.}
\end{table}

\begin{table}[H]
\centering
\caption{AntiFake restoration study extended to Seed-VC and F5-TTS: EER
(SB-ECAPA) per technique and synthesizer (own-synthesizer baseline:
SV2TTS 0.5578, Seed-VC 0.4931, F5-TTS 0.3817), plus signal-quality
metrics for the restored audio. Low-pass+gain reports the range across
all 6 gain values ($M=1.0$--$2.0$); quality shown for $M=1.0$ only.
Bold: combinations that satisfy the all-three-verifiers-significant
reversal criterion (partial reversal -- see discussion below).}
\label{tab:restoration-antifake-seedvc-f5tts}
\begin{tabular}{lcccccc}
\toprule
Technique & SV2TTS & Seed-VC & F5-TTS & STOI & PESQ & SI-SDR (dB) \\
\midrule
Adaptive-filter-centroid  & 0.5066 & 0.4963 & 0.5133 & 0.395 & 1.29 & $-$20.1 \\
Downsampling              & \textbf{0.5254} & 0.4191 & 0.3204 & 0.698 & 1.74 & 1.9 \\
Upsampling                & 0.5651 & 0.4883 & 0.3711 & 0.703 & 1.49 & 2.0 \\
Ensemble averaging        & \textbf{0.5316} & 0.4924 & 0.3654 & 0.702 & 1.67 & 1.2 \\
Mel-spectrogram inversion & 0.5487 & 0.4460 & 0.3793 & 0.633 & 1.33 & $-$33.9 \\
Quantization              & 0.5601 & 0.4887 & 0.3733 & 0.702 & 1.45 & 1.7 \\
Simulated re-recording    & \textbf{0.4781} & 0.4298 & 0.3541 & 0.523 & 1.31 & $-$24.8 \\
Spectral subtraction      & \textbf{0.5410} & 0.5167 & 0.4529 & 0.701 & 1.67 & 3.1 \\
High-pass filtering       & 0.5599 & 0.4915 & 0.3856 & 0.703 & 1.46 & 1.1 \\
Second-cloak-noise        & 0.5568 & 0.4908 & 0.3752 & 0.703 & 1.45 & 1.7 \\
Low-pass + gain ($M$=1.0--2.0) & \textbf{0.504--0.513} & 0.409--0.412 & 0.310--0.318 & 0.691$^{\ast}$ & 1.85$^{\ast}$ & 1.6$^{\ast}$ \\
\bottomrule
\end{tabular}
\end{table}

\paragraph{Result: no full reversal for POP or attack-vc, but AntiFake
shows genuine, if still partial, reversal -- across 60 new
combinations.} Applying the same all-three-verifiers-significant
criterion used throughout \S\ref{sec:restoration-summary}, \textbf{zero
of the 40 POP/attack-vc (technique, cloak-method, synthesizer)
combinations reverse either cloak's protection}. Only four move below
their own baseline on a majority or partial subset of verifiers, all
involving F5-TTS specifically: POP's spectral subtraction (2 of 3
verifiers significant) and, more marginally, upsampling and high-pass
filtering (1 of 3 each); attack-vc's spectral subtraction (2 of 3).
These are exactly the mildest, most near-neutral techniques under
SV2TTS too (\S\ref{sec:restoration-spectral},
\S\ref{sec:restoration-resampling}, \S\ref{sec:restoration-highpass})
-- consistent with F5-TTS's own near-ceiling baseline
(Tables~\ref{tab:protective-efficacy}, \ref{tab:protective-efficacy-attackvc})
leaving very little room for a gentle restoration technique to push EER
any lower without a real, if partial, effect showing up. Every other
POP/attack-vc combination -- including \emph{all twelve} low-pass+gain
values and adaptive-filter-centroid on every synthesizer and cloak
method -- sits clearly \emph{above} its own baseline;
adaptive-filter-centroid remains the single most damaging technique
regardless of downstream synthesizer or cloak method, exactly as
established for SV2TTS (\S\ref{sec:restoration-adaptive-filter}).

AntiFake is different, and consistent with its own SV2TTS-only pattern
(\S\ref{sec:restoration-summary}): the same four techniques that
reverse AntiFake/SV2TTS under this criterion -- downsampling, ensemble
averaging, spectral subtraction, and simulated re-recording, plus the
full low-pass+gain sweep (bolded in
Table~\ref{tab:restoration-antifake-seedvc-f5tts}) -- remain the ones
closest to reversing on Seed-VC and F5-TTS too, though \emph{none}
actually clears the all-three-verifiers bar on either new synthesizer:
each of these techniques satisfies the criterion on only 1 of 3
verifiers under Seed-VC and F5-TTS (2 of 3 for ensemble averaging under
F5-TTS, 2 of 3 for mel-spectrogram inversion under both), never all 3.
So AntiFake's Seed-VC/F5-TTS extension adds zero \emph{new}
all-three-verifiers reversals on top of the 5 already established for
SV2TTS -- reversal, where it happens at all in this restoration study,
remains confined to AntiFake/SV2TTS and remains \textbf{partial}, not
complete: none of these cells approaches the clean-audio ceiling
(\S\ref{sec:restoration-summary}). We use ``partial reversal'', not
``reversal'', throughout to keep this distinction explicit.

Combined with \S\ref{sec:restoration-summary}'s original 30-combination
SV2TTS-only result for all three cloak methods (10 techniques $\times$
POP, attack-vc, and AntiFake), this extends the paper's central
restoration finding to 90 total (technique, cloak-method, synthesizer)
combinations: a signal-domain adaptive attacker gains essentially
nothing against POP or attack-vc regardless of which of the three
downstream cloning systems they target, and against AntiFake gains, at
best, a partial, still-far-from-complete recovery of what SV2TTS could
already partially recover -- with no evidence that switching to a
newer downstream cloning system opens up any \emph{additional} avenue
of attack.

\paragraph{Does descent or gender predict the outcome here too?} We
computed the same descent/gender breakdown used throughout this paper
for all six (cloak-method, synthesizer) groupings introduced in this
subsection (POP/Seed-VC, POP/F5-TTS, attack-vc/Seed-VC,
attack-vc/F5-TTS, AntiFake/Seed-VC, AntiFake/F5-TTS), pooling across
all 16 technique-variants. In every one of the six, the aggregate
gender gap survives restoration in the same direction documented
throughout this paper: women's mean EER exceeds men's (POP/Seed-VC:
17.98\% vs.\ 15.37\%; POP/F5-TTS: 11.51\% vs.\ 10.34\%;
attack-vc/Seed-VC: 19.07\% vs.\ 15.76\%; attack-vc/F5-TTS: 12.04\% vs.\
10.56\%; AntiFake/Seed-VC: 42.92\% vs.\ 40.30\%; AntiFake/F5-TTS:
33.52\% vs.\ 32.55\%). The descent ranking, however, does \emph{not}
generalize to AntiFake the way it does across POP and attack-vc: for
POP and attack-vc, on both new synthesizers, Caucasian(American) and
Caucasian(European) remain lowest and Asian(South) remains highest, the
same ranking documented for SV2TTS throughout this paper. Under
AntiFake, that ranking breaks down, and inconsistently even between
AntiFake's own two new synthesizers -- pooled by descent,
AntiFake/Seed-VC's lowest group is Asian(East) (39.26\%) and highest is
Caucasian(American) (43.20\%), the reverse of the usual ordering, while
AntiFake/F5-TTS instead has Caucasian(American) lowest (27.12\%) and
Asian(South) highest (38.45\%), partially matching it. We read this as
further evidence, alongside the SV2TTS-only AntiFake descent findings
already discussed (\S\ref{sec:restoration-summary}), that AntiFake's
descent-conditioned behavior is a less stable, less generalizable
property of its cloak/verifier interaction than the pattern POP and
attack-vc share. At the level of individual techniques,
adaptive-filter-centroid's gender-direction inversion
(\S\ref{sec:restoration-adaptive-filter}) turns out to be far more
general than its original attack-vc/SV2TTS discovery suggested: it
inverts under attack-vc on \emph{all three} synthesizers (SV2TTS,
Seed-VC, F5-TTS) and under POP/Seed-VC as well, the only exceptions
being POP/SV2TTS and POP/F5-TTS, where the usual direction holds. We
also newly observe simulated re-recording's gender direction fully
inverting under F5-TTS, for \emph{all three} cloak methods (POP,
attack-vc, and now AntiFake as well) -- a full inversion, not the
partial per-descent-group inversion re-recording showed under SV2TTS
(\S\ref{sec:restoration-rerecording}). AntiFake contributes two further,
synthesizer-specific inversions not seen under POP or attack-vc at all:
spectral subtraction inverts under AntiFake/Seed-VC (men's mean EER
43.76\% vs.\ women's 43.33\%), and mel-spectrogram inversion inverts
under AntiFake/F5-TTS (33.90\% vs.\ 31.74\%). We do not have a
confirmed mechanism for any of these patterns and report them as open
questions rather than force them into a single account; the
gender-inverting exceptions in this study now number nine (technique,
cloak-method, synthesizer) cells in total (four for
adaptive-filter-centroid, three for simulated re-recording, one for
spectral subtraction, and one for mel-spectrogram inversion), still
clustered mostly around the two techniques whose motivating rationale
(aggressive band-pass filtering; physical-channel simulation) departs
furthest from every other, gentler technique in this section, with
AntiFake's two new exceptions the only ones that break that pattern.
Tables~\ref{tab:restoration-pop-seedvc-demo}--\ref{tab:restoration-antifake-f5tts-demo}
report the full per-technique breakdown behind the pooled figures
above, for all six (cloak-method, synthesizer) groupings.

\begin{table}[H]
\centering
\caption{POP/Seed-VC restoration: mean EER (\%) by gender and descent
$\times$ gender (FakeAVCeleb subset), per technique, averaged across
the three verifiers. Baseline row is the no-restoration reference.}
\label{tab:restoration-pop-seedvc-demo}
\begin{tabular}{lccccccc}
\toprule
Technique & Men & Women & African & Asian & Asian & Cauc. & Cauc. \\
 & & & & (East) & (South) & (Amer.) & (Euro.) \\
\midrule
Baseline (no restoration) & 7.59 & 9.25 & 9.20 & 7.95 & 10.88 & 7.06 & 7.01 \\
Adaptive-filter-centroid & 49.20 & 46.43 & 46.44 & 47.31 & 51.99 & 45.38 & 47.94 \\
Downsampling & 14.80 & 19.33 & 18.80 & 16.28 & 23.07 & 12.22 & 14.93 \\
Upsampling & 7.78 & 9.35 & 9.23 & 8.38 & 11.03 & 7.24 & 6.94 \\
Ensemble averaging & 8.14 & 10.51 & 10.38 & 9.27 & 12.01 & 7.58 & 7.38 \\
Mel-spectrogram inversion & 8.88 & 11.52 & 10.69 & 10.18 & 12.86 & 9.02 & 8.22 \\
Quantization & 8.88 & 11.32 & 10.84 & 9.85 & 11.95 & 10.03 & 7.84 \\
Simulated re-recording & 15.59 & 19.46 & 17.82 & 17.89 & 17.06 & 19.38 & 15.47 \\
Spectral subtraction & 9.31 & 10.33 & 10.88 & 9.26 & 13.39 & 7.85 & 7.70 \\
High-pass filtering & 7.37 & 9.50 & 8.98 & 8.33 & 11.12 & 6.90 & 6.83 \\
Second-cloak-noise & 7.97 & 10.08 & 10.10 & 9.11 & 10.43 & 8.38 & 7.09 \\
\bottomrule
\end{tabular}
\end{table}

\begin{table}[H]
\centering
\caption{POP/F5-TTS restoration: mean EER (\%) by gender and descent
$\times$ gender (FakeAVCeleb subset), per technique, averaged across
the three verifiers. Baseline row is the no-restoration reference.}
\label{tab:restoration-pop-f5tts-demo}
\begin{tabular}{lccccccc}
\toprule
Technique & Men & Women & African & Asian & Asian & Cauc. & Cauc. \\
 & & & & (East) & (South) & (Amer.) & (Euro.) \\
\midrule
Baseline (no restoration) & 4.99 & 5.52 & 6.50 & 6.70 & 6.08 & 2.50 & 4.51 \\
Adaptive-filter-centroid & 50.42 & 51.83 & 53.34 & 49.61 & 54.94 & 48.27 & 49.46 \\
Downsampling & 7.39 & 9.86 & 10.92 & 11.08 & 10.59 & 4.17 & 6.37 \\
Upsampling & 4.23 & 4.98 & 5.54 & 6.95 & 4.87 & 1.99 & 3.69 \\
Ensemble averaging & 4.35 & 5.70 & 6.53 & 7.18 & 5.39 & 2.27 & 3.76 \\
Mel-spectrogram inversion & 4.62 & 4.90 & 6.13 & 6.61 & 5.18 & 2.26 & 3.63 \\
Quantization & 8.84 & 9.59 & 10.10 & 9.19 & 14.44 & 5.39 & 6.96 \\
Simulated re-recording & \textbf{17.17} & \textbf{14.79} & 16.79 & 13.81 & 25.82 & 10.35 & 13.13 \\
Spectral subtraction & 3.13 & 3.68 & 4.43 & 5.10 & 2.69 & 1.82 & 2.99 \\
High-pass filtering & 4.09 & 4.73 & 5.37 & 6.75 & 4.48 & 2.02 & 3.42 \\
Second-cloak-noise & 6.38 & 7.01 & 7.90 & 7.75 & 9.87 & 2.96 & 4.99 \\
\bottomrule
\end{tabular}

{\small Bold: simulated re-recording's gender direction inverts here
(men's mean EER exceeds women's), the exception discussed above.}
\end{table}

\begin{table}[H]
\centering
\caption{attack-vc/Seed-VC restoration: mean EER (\%) by gender and
descent $\times$ gender (FakeAVCeleb subset), per technique, averaged
across the three verifiers. Baseline row is the no-restoration
reference.}
\label{tab:restoration-attackvc-seedvc-demo}
\begin{tabular}{lccccccc}
\toprule
Technique & Men & Women & African & Asian & Asian & Cauc. & Cauc. \\
 & & & & (East) & (South) & (Amer.) & (Euro.) \\
\midrule
Baseline (no restoration) & 7.57 & 8.90 & 9.30 & 8.04 & 11.14 & 6.54 & 6.14 \\
Adaptive-filter-centroid & \textbf{49.46} & \textbf{43.85} & 44.09 & 46.52 & 48.43 & 45.00 & 49.24 \\
Downsampling & 14.46 & 19.31 & 18.07 & 15.98 & 23.92 & 11.51 & 14.96 \\
Upsampling & 7.65 & 9.20 & 9.75 & 8.27 & 10.70 & 6.68 & 6.72 \\
Ensemble averaging & 7.99 & 10.31 & 10.46 & 9.39 & 11.61 & 6.88 & 7.42 \\
Mel-spectrogram inversion & 8.93 & 11.75 & 10.99 & 10.48 & 13.10 & 8.49 & 8.64 \\
Quantization & 8.39 & 11.38 & 10.63 & 9.83 & 11.62 & 9.05 & 8.29 \\
Simulated re-recording & 14.58 & 18.71 & 16.80 & 17.14 & 16.90 & 17.89 & 14.50 \\
Spectral subtraction & 8.91 & 10.40 & 10.51 & 8.96 & 13.95 & 7.34 & 7.54 \\
High-pass filtering & 7.23 & 9.57 & 9.26 & 8.38 & 10.76 & 6.69 & 6.92 \\
Second-cloak-noise & 7.91 & 10.49 & 10.11 & 9.51 & 11.07 & 7.84 & 7.47 \\
\bottomrule
\end{tabular}

{\small Bold: adaptive-filter-centroid's gender direction inverts here
(men's mean EER exceeds women's), the same inversion first documented
under attack-vc/SV2TTS (\S\ref{sec:restoration-adaptive-filter}).}
\end{table}

\begin{table}[H]
\centering
\caption{attack-vc/F5-TTS restoration: mean EER (\%) by gender and
descent $\times$ gender (FakeAVCeleb subset), per technique, averaged
across the three verifiers. Baseline row is the no-restoration
reference.}
\label{tab:restoration-attackvc-f5tts-demo}
\begin{tabular}{lccccccc}
\toprule
Technique & Men & Women & African & Asian & Asian & Cauc. & Cauc. \\
 & & & & (East) & (South) & (Amer.) & (Euro.) \\
\midrule
Baseline (no restoration) & 3.96 & 4.73 & 5.77 & 6.14 & 4.33 & 2.25 & 3.24 \\
Adaptive-filter-centroid & \textbf{51.99} & \textbf{47.64} & 51.74 & 46.37 & 53.79 & 47.84 & 49.33 \\
Downsampling & 6.92 & 9.37 & 10.18 & 10.49 & 9.65 & 4.18 & 6.20 \\
Upsampling & 3.72 & 5.06 & 5.68 & 6.45 & 4.43 & 2.01 & 3.36 \\
Ensemble averaging & 3.72 & 5.05 & 6.02 & 6.31 & 3.89 & 2.11 & 3.57 \\
Mel-spectrogram inversion & 4.36 & 5.07 & 6.54 & 6.58 & 4.57 & 2.45 & 3.43 \\
Quantization & 8.25 & 9.07 & 10.06 & 9.13 & 13.10 & 4.53 & 6.49 \\
Simulated re-recording & \textbf{15.81} & \textbf{14.40} & 16.56 & 13.25 & 24.17 & 9.01 & 12.54 \\
Spectral subtraction & 2.64 & 3.34 & 4.16 & 4.23 & 1.89 & 1.90 & 2.77 \\
High-pass filtering & 3.69 & 4.54 & 5.68 & 6.14 & 3.70 & 2.24 & 2.82 \\
Second-cloak-noise & 5.75 & 6.86 & 7.49 & 7.66 & 9.11 & 2.85 & 4.40 \\
\bottomrule
\end{tabular}

{\small Bold: both adaptive-filter-centroid and simulated re-recording
invert here -- one of three (cloak-method, synthesizer) groupings in
this study where two techniques invert simultaneously, the others being
AntiFake/F5-TTS below (mel-spectrogram inversion and simulated
re-recording) and, in the opposite direction, AntiFake/Seed-VC
(spectral subtraction alone).}
\end{table}

\begin{table}[H]
\centering
\caption{AntiFake/Seed-VC restoration: mean EER (\%) by gender and
descent $\times$ gender (FakeAVCeleb subset), per technique, averaged
across the three verifiers. Baseline row is the no-restoration
reference.}
\label{tab:restoration-antifake-seedvc-demo}
\begin{tabular}{lccccccc}
\toprule
Technique & Men & Women & African & Asian & Asian & Cauc. & Cauc. \\
 & & & & (East) & (South) & (Amer.) & (Euro.) \\
\midrule
Baseline (no restoration) & 41.28 & 43.45 & 41.46 & 39.66 & 40.48 & 46.62 & 43.61 \\
Adaptive-filter-centroid & 49.88 & 53.74 & 51.87 & 51.23 & 55.93 & 48.84 & 51.19 \\
Downsampling & 38.65 & 42.72 & 42.32 & 37.82 & 42.36 & 41.35 & 39.59 \\
Upsampling & 40.96 & 43.40 & 41.14 & 39.65 & 40.85 & 46.57 & 42.70 \\
Ensemble averaging & 40.91 & 42.89 & 41.54 & 39.73 & 39.33 & 46.46 & 42.43 \\
Mel-spectrogram inversion & 38.28 & 39.10 & 38.91 & 36.48 & 36.93 & 42.62 & 38.49 \\
Quantization & 40.84 & 43.24 & 40.93 & 39.95 & 40.13 & 46.84 & 42.37 \\
Simulated re-recording & 38.33 & 38.89 & 40.04 & 36.57 & 36.28 & 42.68 & 37.47 \\
Spectral subtraction & \textbf{43.76} & \textbf{43.33} & 42.57 & 39.51 & 42.36 & 48.21 & 45.06 \\
High-pass filtering & 40.45 & 43.45 & 40.47 & 39.58 & 40.67 & 46.41 & 42.61 \\
Second-cloak-noise & 41.39 & 42.80 & 41.00 & 39.73 & 40.34 & 46.40 & 43.00 \\
\bottomrule
\end{tabular}

{\small Bold: spectral subtraction's gender direction inverts here
(men's mean EER exceeds women's), a new inversion not observed for this
technique under any other cloak method or synthesizer in this study.}
\end{table}

\begin{table}[H]
\centering
\caption{AntiFake/F5-TTS restoration: mean EER (\%) by gender and
descent $\times$ gender (FakeAVCeleb subset), per technique, averaged
across the three verifiers. Baseline row is the no-restoration
reference.}
\label{tab:restoration-antifake-f5tts-demo}
\begin{tabular}{lccccccc}
\toprule
Technique & Men & Women & African & Asian & Asian & Cauc. & Cauc. \\
 & & & & (East) & (South) & (Amer.) & (Euro.) \\
\midrule
Baseline (no restoration) & 32.82 & 32.95 & 36.09 & 28.70 & 40.08 & 26.26 & 33.29 \\
Adaptive-filter-centroid & 49.28 & 54.86 & 51.83 & 50.73 & 56.46 & 49.57 & 51.77 \\
Downsampling & 28.45 & 32.38 & 32.78 & 30.61 & 34.26 & 25.01 & 29.40 \\
Upsampling & 31.19 & 32.19 & 34.25 & 27.97 & 38.33 & 26.10 & 31.81 \\
Ensemble averaging & 31.16 & 31.82 & 34.13 & 28.61 & 37.87 & 24.67 & 32.17 \\
Mel-spectrogram inversion & \textbf{33.90} & \textbf{31.74} & 36.07 & 28.88 & 39.73 & 26.14 & 33.28 \\
Quantization & 31.91 & 32.19 & 34.76 & 28.57 & 39.10 & 25.41 & 32.43 \\
Simulated re-recording & \textbf{34.20} & \textbf{31.28} & 36.03 & 27.14 & 42.21 & 24.85 & 33.45 \\
Spectral subtraction & 35.08 & 36.79 & 35.24 & 32.37 & 39.62 & 34.14 & 38.31 \\
High-pass filtering & 32.93 & 33.08 & 36.02 & 28.74 & 40.26 & 26.63 & 33.39 \\
Second-cloak-noise & 32.10 & 32.84 & 35.58 & 28.87 & 39.57 & 25.83 & 32.50 \\
\bottomrule
\end{tabular}

{\small Bold: both mel-spectrogram inversion and simulated re-recording
invert here (men's mean EER exceeds women's) -- mel-spectrogram
inversion's inversion is new to this (cloak-method, synthesizer)
grouping, while simulated re-recording's continues the same
full-inversion-under-F5-TTS pattern already documented for POP and
attack-vc above.}
\end{table}

\paragraph{Quality metrics: POP is measurably milder than attack-vc at
every technique, and SI-SDR's sensitivity to phase/alignment artifacts
(already documented for attack-vc's own cloak, \S\ref{sec:results-attackvc})
again makes it the least interpretable of the three in absolute terms.}
POP's STOI (0.68--0.97, excluding adaptive-filter-centroid's outlying
0.55) and SI-SDR (mostly positive, 12--16~dB) sit far above attack-vc's
corresponding figures for the \emph{same} techniques (STOI 0.43--0.55;
SI-SDR uniformly negative, $-$34 to $-$42~dB) -- consistent with POP's
own cloak-level signal quality already being substantially higher than
attack-vc's elsewhere in this paper (Tables~\ref{tab:results},
\ref{tab:results-attackvc}: POP STOI~0.971 vs.\ attack-vc STOI~0.547),
a difference restoration preserves rather than erases. Within each
cloak method, PESQ agrees in relative ranking with the EER pattern:
adaptive-filter-centroid and simulated re-recording are consistently
the lowest-PESQ, highest-EER techniques for both cloaks, while
upsampling and spectral subtraction/high-pass filtering are
consistently the highest-PESQ, lowest-EER ones. SI-SDR is far more
uniform \emph{within} attack-vc (a narrow $-$34 to $-$42~dB band
regardless of technique) than within POP (a much wider $-$31 to
$+$16~dB spread) -- consistent with the sample-alignment sensitivity
already flagged for attack-vc specifically in \S\ref{sec:results-attackvc}:
attack-vc's own cloak already introduces enough of a phase/timing shift
that SI-SDR saturates near its collapsed floor regardless of which
restoration technique is layered on top, whereas POP's gentler cloak
leaves enough headroom for SI-SDR to still discriminate between
techniques. We continue to read PESQ, not SI-SDR, as the more reliable
relative-quality signal in this section, consistent with
\S\ref{sec:results-attackvc}.

AntiFake's restored-audio quality (Table~\ref{tab:restoration-antifake-seedvc-f5tts})
sits closer to POP's profile than to attack-vc's: STOI 0.52--0.70
(excluding adaptive-filter-centroid's outlying 0.40) and PESQ 1.29--1.85
are both in the same range as POP's own figures above, and SI-SDR is
positive for seven of the ten techniques (1.1--3.1~dB). The three
exceptions -- adaptive-filter-centroid ($-$20.1~dB), simulated
re-recording ($-$24.8~dB), and mel-spectrogram inversion
($-$33.9~dB) -- are exactly the same three techniques that produce
negative SI-SDR under POP, and no others do under either cloak; this
cross-cloak consistency suggests SI-SDR's sign here is a property of
\emph{which restoration technique} is applied (specifically, how much
it disturbs sample-level phase/timing) rather than of the cloak being
restored. attack-vc breaks this pattern only because its own cloak
already introduces enough phase/timing distortion
(\S\ref{sec:results-attackvc}) that SI-SDR saturates negative for
\emph{every} technique, masking the technique-level distinction that
remains visible under POP and AntiFake's gentler cloaks. Despite this
quality similarity to POP, AntiFake is the one cloak where several
techniques still produce a statistically-supported partial reversal on
SV2TTS (\S\ref{sec:restoration-summary}) -- reinforcing that restored
signal quality and downstream verification-leakage risk are related but
distinct measurements, not interchangeable proxies for one another.

\section{Discussion and Limitations}
\label{sec:limitations}

\paragraph{This evaluation measures imperceptibility, not protective
efficacy.} We wish to state this limitation prominently rather than
in passing. POP's threat model (\S\ref{sec:related}) is that a
downstream attacker \emph{trains or fine-tunes a TTS model on the
cloaked recording}; the defense's success criterion is that
\emph{newly synthesized speech from that fine-tuned model} fails to
sound like the true speaker. The evaluation in this paper -- comparing
the cloaked recording directly against the genuine original -- instead
measures whether the perturbation is detectable on its own, which is a
\emph{precondition} for a usable cloak (an easily detected or
disruptive-sounding perturbation would be rejected by users or
platforms) but is not itself evidence that the cloak protects against
its stated threat. A complete protective-efficacy evaluation requires
an additional synthesis stage -- feeding the cloaked recordings into
one of the voice-cloning/conversion systems of
Section~\ref{sec:synth} as fine-tuning or reference data, and
comparing \emph{that} system's output against the genuine original
using the same trial-construction protocol. We report this evaluation
in Section~\ref{sec:protective-efficacy} for all three systems
(SV2TTS, Seed-VC, and F5-TTS).

\paragraph{Whisper-derived transcripts.} Two of our three source
corpora carry no official transcript, so 4,367 of 4,867 transcripts
used to condition POP's surrogate model are ASR-generated rather than
ground truth. Transcription errors could weaken the surrogate's
text/audio alignment and, in turn, the crafted perturbation's
effectiveness for the affected recordings; we quantify and exclude
only the most severe failures (hallucinations) rather than every
possible minor transcription error, so some residual noise from this
source likely remains in the reported numbers. A related phenomenon
affects every F5-TTS clone in this paper, where the reference
transcript is Whisper-generated on the fly from whatever audio is
actually being cloned, and can be degraded by the very perturbation
being evaluated -- \S\ref{sec:ablation-f5tts-transcript} argues this
is correctly read as part of a real attacker's experience rather than
a confound to remove, but flags the one place (23 crashed AntiFake
clips) where we deviated from that by necessity.

\paragraph{Single surrogate model.} POP's own paper argues for
transferability to unseen downstream models via the choice of a
generic, widely used surrogate (VITS); we evaluate only against that
same surrogate rather than an independent, held-out target model, so
we cannot yet speak to POP's transferability claim specifically,
only to its behavior on the surrogate it was designed against.

\section{Conclusion}
\label{sec:conclusion}
We reproduced POP's official training-time voice-cloaking attack end
to end and applied it, for the first time at this scale and
demographic breadth, to a 4,867-recording, tri-corpus, descent- and
gender-balanced evaluation set. We documented, precisely and
reproducibly, every metric (verification-leakage and signal-quality),
every model choice (three verifiers, three cloning/conversion systems,
four cloaking techniques), and the dataset construction underlying
this and the three companion evaluations (ProtectYourAudio, attack-vc,
AntiFake) in our broader project, along with the practical engineering
obstacles -- a silent data-loader bug, GPU memory ceilings, missing
system dependencies -- encountered building all four from scratch on
shared, resource-constrained hardware. Aggregated across all 4,864
successfully-cloaked recordings, POP preserves speaker-verification
identity almost perfectly (EER~$\leq$~0.16\%, TAR@1\%FAR~=~100.00\%
for all three verifiers) with high acoustic fidelity (STOI~0.971,
PESQ~2.92, SI-SDR~16.3~dB), and this holds uniformly across all ten
FakeAVCeleb descent$\times$gender groups (Table~\ref{tab:ablation}).
We also completed attack-vc's equivalent imperceptibility evaluation
(Section~\ref{sec:results-attackvc}), and, distinct from either
imperceptibility result, the resynthesis-based protective-efficacy
evaluation for POP against all three systems
(Section~\ref{sec:protective-efficacy}): POP measurably degrades what
a downstream SV2TTS attacker can extract (EER up 15--86\% relative
over the clean-audio ceiling), has a real but more modest effect
against a newer zero-shot TTS attacker whose security-relevant
TAR@1\%FAR barely moves (F5-TTS, EER up 13--54\% relative, TAR@1\%FAR
down under 10\% relative), and has only a marginal effect against a
Seed-VC attacker (EER up 5--15\% relative) -- a system-dependent
protective effect that an imperceptibility result alone cannot
predict.

For attack-vc, all three protective-efficacy runs also completed
(Section~\ref{sec:protective-efficacy-attackvc}), reproducing the same
qualitative pattern despite attack-vc and POP being mechanistically
unrelated attacks: strong against SV2TTS (EER up 15--95\% relative),
negligible against Seed-VC (EER up 2--6\% relative), and a nuanced
result for F5-TTS whose EER rises moderately (10--28\% relative) while
its TAR@1\%FAR -- the more security-relevant, fixed-operating-point
metric -- barely moves, underscoring why we report both rather than
EER alone. The demographic breakdown
(\S\ref{sec:protective-efficacy-demographics-attackvc}) further shows
that women and, depending on the downstream system, certain Asian
descent groups are consistently harder to re-identify after cloaking
than men or Caucasian descent groups -- but this disparity already
exists in each system's own clean-audio ceiling, so it is inherited
from the synthesizer/verifier pipeline rather than introduced by
either cloaking technique.

We also tested, rather than merely speculated about, why Seed-VC is so
much less affected than SV2TTS by either cloak
(\S\ref{sec:ablation-seedvc-mechanism}). An architectural check found
no shared encoder family between either cloak's surrogate/target and
Seed-VC's or SV2TTS's own conditioning mechanism -- and POP's surrogate
has no generalizable speaker encoder at all, only a closed lookup table
over its training speakers -- ruling out simple white-box transfer as
the explanation. An embedding-level probe (n=500 per cloak, no
synthesis involved) then showed the opposite of our original
hypothesis: both cloaks' perturbations reach Seed-VC's own CAMPPlus
identity embedding at least as much as SV2TTS's GE2E embedding (mean
cosine similarity for attack-vc: 0.90 for Seed-VC vs.\ 0.98 for
SV2TTS), yet this barely affects Seed-VC's end-to-end EER. Seed-VC's
resistance therefore looks like a property of its decoder tolerating a
perturbed conditioning signal, not of the perturbation failing to
reach that signal -- a more specific and better-supported claim than
our original ``TTS vs.\ voice conversion'' framing, though confirming
the decoder-robustness mechanism directly remains future work.

AntiFake's cloak generation is now complete for 4,844 of the
4,867-recording corpus (\S\ref{sec:cloaking}), closing out this
project's four-technique comparison. It behaves unlike either other
technique. Its imperceptibility fails outright
(Section~\ref{sec:results-antifake}: EER 22--66\%, TAR@1\%FAR only
5--25\%, roughly two orders of magnitude worse than POP or attack-vc),
and every signal-quality metric agrees the perturbation is genuinely
audible (STOI 0.703, the lowest PESQ of the three techniques at 1.46,
SI-SDR 1.9~dB). Its protective efficacy is correspondingly the
strongest measured in this project and, uniquely, effective against
\emph{all three} downstream systems rather than only SV2TTS
(Section~\ref{sec:protective-efficacy-antifake}): EER rises 179--814\%
relative against SV2TTS, 116--1{,}485\% relative against Seed-VC, and
226--2{,}199\% relative against F5-TTS -- the largest effect of any
(system, verifier) cell in this project, against the very system where
POP and attack-vc showed only a moderate effect with a
security-relevant TAR@1\%FAR that barely moved. The demographic pattern
is also inverted relative to POP/attack-vc for SV2TTS and Seed-VC:
Caucasian descent groups show both the \emph{worst} imperceptibility
and the \emph{strongest} protection under those two systems, consistent
with a single embedding-displacement mechanism driving both
measurements together, rather than the downstream-pipeline-inherited
disparity found for POP and attack-vc
(\S\ref{sec:protective-efficacy-demographics-antifake}) -- though
F5-TTS breaks even this pattern, with its gender gap nearly vanishing
and its descent ranking partially inverting, an open question we do not
force into the same account.

ProtectYourAudio's cloak generation and evaluation are likewise now
complete (4,850 of the 4,867-recording corpus). Its imperceptibility
(Section~\ref{sec:results-protectyouraudio}: EER below 2.1\%,
TAR@1\%FAR above 96.2\% for every verifier) is essentially
indistinguishable from attack-vc's, verifier by verifier, despite the
two techniques being mechanistically unrelated (a discrete
frequency-band-masking search vs.\ a continuous embedding-attack
perturbation) -- though its signal-quality profile is neither POP-like
nor attack-vc-like, pairing POP-level STOI (0.941) with an
attack-vc-level negative SI-SDR ($-$25.4~dB), a combination no other
technique in this paper produces. Its protective efficacy
(Section~\ref{sec:protective-efficacy-protectyouraudio}) is where
ProtectYourAudio genuinely distinguishes itself from POP and attack-vc:
where those two techniques showed only a negligible-to-moderate effect
against Seed-VC and F5-TTS (EER up 2--28\% relative, barely above
noise), ProtectYourAudio shows a real effect against both (Seed-VC EER
up 15--100\% relative; F5-TTS EER up 62--151\% relative, the largest
relative shift against F5-TTS reported anywhere in this paper outside
AntiFake), while its SV2TTS effect (32--132\% relative) remains
comparable to POP's and attack-vc's own strongest results. The
demographic pattern, unlike AntiFake's, follows the same
pipeline-inherited direction as POP and attack-vc
(\S\ref{sec:protective-efficacy-demographics-protectyouraudio}): women
and Caucasian descent groups are consistently easier to re-identify
after cloaking, a disparity present in each system's own clean-audio
ceiling rather than introduced by ProtectYourAudio's own search
objective.

Read as a whole, this project's four-technique comparison suggests a
real tradeoff between imperceptibility and protective efficacy, but
not a strictly binary one: POP and attack-vc stay close to
imperceptible but protect unevenly across downstream systems, AntiFake
sacrifices imperceptibility entirely for protection that is stronger
and more uniform across attacker types, and ProtectYourAudio sits
between these extremes -- matching POP's and attack-vc's near-perfect
imperceptibility while protecting measurably more broadly across
downstream systems than either, without AntiFake's audible cost. No
technique evaluated here achieves AntiFake's degree of cross-system
protection \emph{and} POP's degree of imperceptibility simultaneously,
but ProtectYourAudio demonstrates that the two objectives are not as
sharply opposed as POP and attack-vc's results alone would suggest --
a discrete, budget-constrained search over a restricted perturbation
space can apparently generalize across downstream systems better than
a continuous, unconstrained-direction perturbation of similar
imperceptibility, a distinction future cloak designs may be able to
exploit deliberately rather than as an incidental side effect.

\bibliographystyle{plain}
\begin{thebibliography}{99}

\bibitem{zhang2024mitigating}
Z.~Zhang, Q.~Yang, D.~Wang, P.~Huang, Y.~Cao, K.~Ye, and J.~Hao,
``Mitigating Unauthorized Speech Synthesis for Voice Protection,''
in \emph{Proc. 1st ACM Workshop on Large AI Systems and Models with
Privacy and Safety Analysis (LAMPS)}, 2024.

\bibitem{yu2023antifake}
Z.~Yu, S.~Zhai, and N.~Zhang,
``AntiFake: Using Adversarial Audio to Prevent Unauthorized Speech
Synthesis,''
in \emph{Proc. ACM SIGSAC Conference on Computer and Communications
Security (CCS)}, 2023.

\bibitem{liu2023protecting}
Z.~Liu, Y.~Zhang, and C.~Miao,
``Protecting Your Voice from Speech Synthesis Attacks,''
in \emph{Proc. Annual Computer Security Applications Conference
(ACSAC)}, 2023.

\bibitem{huang2021defending}
C.~Huang, Y.~Y.~Lin, H.~Lee, and L.~Lee,
``Defending Your Voice: Adversarial Attack on Voice Conversion,''
in \emph{Proc. IEEE Spoken Language Technology Workshop (SLT)}, 2021.

\bibitem{huang2021unlearnable}
H.~Huang, X.~Ma, S.~M.~Erfani, J.~Bailey, and Y.~Wang,
``Unlearnable Examples: Making Personal Data Unexploitable,''
in \emph{Proc. International Conference on Learning Representations
(ICLR)}, 2021.

\bibitem{jia2018transfer}
Y.~Jia, Y.~Zhang, R.~Weiss, Q.~Wang, J.~Shen, F.~Ren, Z.~Chen,
P.~Nguyen, R.~Pang, I.~L.~Moreno, and Y.~Wu,
``Transfer Learning from Speaker Verification to Multispeaker
Text-To-Speech Synthesis,''
in \emph{Advances in Neural Information Processing Systems
(NeurIPS)}, 2018.

\bibitem{kim2021vits}
J.~Kim, J.~Kong, and J.~Son,
``Conditional Variational Autoencoder with Adversarial Learning for
End-to-End Text-to-Speech,''
in \emph{Proc. International Conference on Machine Learning (ICML)},
2021.

\bibitem{chen2024f5tts}
Y.~Chen, Z.~Niu, Z.~Ma, K.~Deng, C.~Wang, J.~Zhao, K.~Yu, and X.~Chen,
``F5-TTS: A Fairytaler that Fakes Fluent and Faithful Speech with
Flow Matching,''
\emph{arXiv preprint}, 2024.

\bibitem{desplanques2020ecapa}
B.~Desplanques, J.~Thienpondt, and K.~Demuynck,
``ECAPA-TDNN: Emphasized Channel Attention, Propagation and
Aggregation in TDNN Based Speaker Verification,''
in \emph{Proc. Interspeech}, 2020.

\bibitem{chen2022wavlm}
S.~Chen, C.~Wang, Z.~Chen, Y.~Wu, S.~Liu, Z.~Chen, J.~Li, N.~Kanda,
T.~Yoshioka, X.~Xiao, J.~Wu, L.~Zhou, S.~Ren, Y.~Qian, Y.~Qian,
J.~Wu, M.~Zeng, and F.~Wei,
``WavLM: Large-Scale Self-Supervised Pre-Training for Full Stack
Speech Processing,''
\emph{IEEE Journal of Selected Topics in Signal Processing}, 2022.

\bibitem{panayotov2015librispeech}
V.~Panayotov, G.~Chen, D.~Povey, and S.~Khudanpur,
``Librispeech: An ASR Corpus Based on Public Domain Audio Books,''
in \emph{Proc. ICASSP}, 2015.

\bibitem{yamagishi2021asvspoof}
J.~Yamagishi, X.~Wang, M.~Todisco, M.~Sahidullah, J.~Patino,
A.~Nautsch, X.~Liu, K.~A.~Lee, T.~Kinnunen, N.~Evans, and H.~Delgado,
``ASVspoof 2021: Accelerating Progress in Spoofed and Deepfake Speech
Detection,''
in \emph{Proc. ASVspoof 2021 Workshop}, 2021.

\bibitem{khalid2021fakeavceleb}
H.~Khalid, S.~Tariq, M.~Kim, and S.~S.~Woo,
``FakeAVCeleb: A Novel Audio-Video Multimodal Deepfake Dataset,''
in \emph{Proc. NeurIPS Datasets and Benchmarks Track}, 2021.

\bibitem{asvspoof2021doi}
ASVspoof Consortium,
``ASVspoof Challenge Series,''
\url{https://www.asvspoof.org}, accessed 2026.

\end{thebibliography}

\end{document}
